\documentclass[11pt,a4paper]{article}
\usepackage[T1]{fontenc}
\usepackage[utf8]{inputenc}
\usepackage{lmodern}
\usepackage[margin=27mm,headheight=14pt]{geometry}
\usepackage{amsmath,amssymb,amsthm,mathtools}
\usepackage{booktabs,longtable,array,tabularx}
\usepackage{microtype,xcolor,enumitem,fancyhdr,listings}
\usepackage{hyperref}
\usepackage[nameinlink,noabbrev]{cleveref}
\definecolor{inkblue}{HTML}{173B57}
\definecolor{codegray}{HTML}{F4F6F8}
\hypersetup{colorlinks=true,linkcolor=inkblue,citecolor=inkblue,urlcolor=inkblue,
 pdftitle={Certified Prime Lifting from Known Primes},pdfauthor={Keio},
 pdfsubject={Revised edition with eighteen additional algorithms and exact verifiers}}
\urlstyle{same}
\pagestyle{fancy}
\fancyhf{}\fancyhead[L]{\small Certified Prime Lifting}\fancyhead[R]{\small Keio | Revised edition}
\fancyfoot[C]{\thepage}\renewcommand{\headrulewidth}{0.3pt}
\setlength{\parskip}{3pt}\setlength{\parindent}{1.3em}
\setlength{\emergencystretch}{2em}
\setlist{itemsep=2pt,topsep=4pt}
\newtheorem{theorem}{Theorem}[section]
\newtheorem{proposition}[theorem]{Proposition}
\newtheorem{corollary}[theorem]{Corollary}
\newtheorem{lemma}[theorem]{Lemma}
\theoremstyle{definition}\newtheorem{definition}[theorem]{Definition}
\theoremstyle{remark}\newtheorem{remark}[theorem]{Remark}
\newcommand{\Z}{\mathbb Z}\newcommand{\F}{\mathbb F}
\newcommand{\ord}{\operatorname{ord}}\newcommand{\rad}{\operatorname{rad}}
\newcommand{\lcm}{\operatorname{lcm}}\newcommand{\isqrt}{\operatorname{isqrt}}
\newcommand{\bits}{\operatorname{bits}}\newcommand{\Norm}{\operatorname{N}}
\newcommand{\sha}[1]{\texttt{\small #1}}
\lstset{language=Python,basicstyle=\ttfamily\scriptsize,columns=fullflexible,
 keepspaces=true,showstringspaces=false,breaklines=true,breakatwhitespace=false,
 backgroundcolor=\color{codegray},frame=single,rulecolor=\color{black!15},
 framesep=5pt,numbers=none,tabsize=4,aboveskip=9pt,belowskip=9pt}
\title{\color{inkblue}\bfseries Certified Prime Lifting from Known Primes\\[5pt]
 \large Eighteen Additional Algorithms for Construction,\\
 Arithmetic, and Certificate Synthesis\\[3pt]
 \normalsize Revised and expanded edition}
\author{Keio}
\date{6 November 2026}

\begin{document}
\maketitle
\begin{abstract}
We study how certified primes can be reused to construct and prove larger
primes, and how several incomplete certificates can be combined into a
complete proof. The baseline investigation uses multiplicative orders,
Lucas sequences, cyclotomic transfers, and elliptic curves; its inventory
of 32 distinct certified integers and its complete verifier are retained.
Two subsequent chain extensions bring the incorporated computation to
18,885 decimal digits. This revised edition adds eighteen explicit
algorithms with input contracts, correctness proofs, and executable
demonstrations. They include coefficient-sensitive closures for
$ap^2+1$, $ap^3+1$, and $ap^2-1$; digit-constrained generation; exact
batch arithmetic; witness and elliptic-order fusion; and an intersection
of residue classes with integer Hasse windows. In particular, a full
Fermat congruence and short nonunity certify $ap^2+1$ for positive even
$a<4p-2$, while a short-index Lucas certificate applies to nonsquare
positive even $a<4p+2$ in $ap^2-1$. The construction prototypes certify
outputs through 117 digits, and the positive quadratic criterion also
recertifies the existing 18,882-digit chain member. Independent finite
checks include exact trial division, direct recurrences, invalid-input
controls, and exhaustive small optimization problems. All numerical
constants and six executable verification listings are embedded in this
single source. These are additions and refinements developed for this
manuscript, not established claims of historical priority, universal speed
superiority, a new world record, or a resolution of the Collatz or Mills
transcendence problems.
\end{abstract}

\noindent\textbf{Keywords.} Provable prime generation; Pocklington certificate;
Lucas sequence; cyclotomic polynomial; elliptic primality certificate;
complex multiplication; reproducible computation.

\noindent\textbf{Scope.} This is a computational research manuscript, not a
claim of publication or peer review. It consolidates an interactive
investigation and records which conclusions have explicit proofs and which
remain proposals.

\medskip
\noindent\textbf{Reading the revision.}
The baseline constructions and their original verifier are preserved.
Section~\ref{sec:revision-ledger} records the later chain and the distinction
between prime counts and algorithm counts.
Sections~\ref{sec:new-constructions}, \ref{sec:arith18}, and
\ref{sec:cert-synthesis} contain A01--A18.
Section~\ref{sec:revision-reproduce} explains how to extract every complete
program from this one TEX file.

\clearpage
\begingroup\small\setcounter{tocdepth}{2}\setlength{\parskip}{0pt}\tableofcontents\endgroup
\clearpage

\section{Introduction and the precise problem}
\label{sec:intro}

The input is one or more integers whose primality is already proved. The aim is
to construct a larger integer together with a short, independently checkable
reason that it is prime. A useful construction must therefore specify three
different objects: a family of candidates, a search within that family, and a
certificate for each accepted candidate. Conflating these objects leads to
incorrect conclusions, especially when a formula is short but its numerical
evaluation involves millions of digits.

\begin{definition}
A \emph{certified prime lift} is a tuple consisting of certified prime inputs,
an explicitly defined output integer $N$ larger than its designated input, and
a finite certificate whose validity implies that $N$ is prime. Different
inputs may be used together; an input prime may occur to a high power.
\end{definition}

This terminology describes a workflow, not a new primality theorem. Recursive
generation of provable primes is well established, notably in Maurer's work
\cite{maurer}. The $N-1$ and $N+1$ methods used here belong to the classical
Pocklington, Lucas, Lehmer, Morrison, and Brillhart--Lehmer--Selfridge framework
\cite{pocklington,bls,morrison}. The elliptic argument is the
Goldwasser--Kilian certificate principle \cite{gk,sutherland}; constructing
curves with prescribed order is related to reverse CM methods
\cite{broker}. Our contribution is an explicit synthesis, elementary proofs
for selected specializations, and a reproducible finite collection of
certificates.

\subsection{What the computations establish}

The retained baseline investigation separates three cohorts.
\begin{enumerate}[label=\textup{(\roman*)}]
\item The first experiment starts from two Mersenne primes and one certified
Proth prime and produces ten larger primes. The largest has 4,153 digits.
\item A previously reported cyclotomic seed and four successive squaring
lifts are reconstructed and certified, ending at 9,439 digits. The coefficients
of this chain came from the preceding conversation; their verification is not
reported as a fresh search.
\item The second experiment produces nine principal outputs and two auxiliary
cyclotomic outputs. These include three elliptic certificates, two pairs of
cyclotomic transfers, a negative cubic Lucas certificate, and a linear lift
from the 4,718-digit chain member.
\end{enumerate}
An additional experiment in \cref{sec:collatz} supplies two primes linked by
an exact Collatz orbit segment. After duplicate seeds are removed, the combined
inventory contains 32 distinct primes.
All integer identities and digit counts are exact. The certificate checks are
deterministic integer computations, although search routines may use probable
prime tests to reject candidates cheaply. No probable-prime pass is accepted
as the final evidence of primality.

\subsection{What ``immediate'' can and cannot mean}

A formula such as $N=cP^2+1$ can be written immediately, and its factorization
$N-1=cP^2$ immediately identifies potentially useful information. Neither fact
proves $N$ prime. For example, $P=3$ and $c=6$ give $N=55$. Conversely, an
already supplied certificate can often be checked much faster than a suitable
candidate can be found. The numerical results below distinguish these two
costs explicitly.

Large known primes can also be inserted into expressions that are immediately
proved composite by small congruences. This is useful pruning, but is not a
new giant-prime discovery. Likewise, an identity that rewrites an already known
larger prime in terms of a smaller one reuses existing knowledge; it does not
transfer primality by algebra alone.

\section{Three reusable certificate mechanisms}
\label{sec:mechanisms}

We write $v_q(m)$ for the exponent of a prime $q$ in $m$. A \emph{fully factored}
integer means that all of its prime factors, their multiplicities, and proofs
of the primality of those factors are available. A proof may depend on other
proofs, but the dependency graph must terminate in independently verified
seeds or small primes proved by trial division.

\subsection{Multiplicative orders from a factored part of \texorpdfstring{$N-1$}{N-1}}

\begin{theorem}[Pocklington order principle]\label{thm:minus}
Let $N>1$ be odd and let $F>1$ divide $N-1$. Suppose that $F$ is fully factored.
For every distinct prime $q\mid F$, suppose that an integer $a_q$ satisfies
\begin{equation}\label{eq:pock}
 a_q^{N-1}\equiv1\pmod N,\qquad
 \gcd\bigl(a_q^{(N-1)/q}-1,N\bigr)=1.
\end{equation}
Then every prime $\ell\mid N$ satisfies $F\mid\ell-1$. Consequently each of
the following conditions is sufficient for the primality of $N$:
\begin{equation}\label{eq:minus-bounds}
 (F+1)^2>N;\qquad
 (2F+1)^2>N\quad\text{if $F$ is odd}.
\end{equation}
In particular, $F^2>N$ is a convenient sufficient bound.
\end{theorem}

\begin{proof}
Fix a prime $\ell\mid N$ and a prime $q\mid F$. The first congruence makes
$a_q$ a unit modulo $\ell$. Put $h=\ord_{\ell}(a_q)$. Then $h\mid N-1$,
whereas the greatest-common-divisor condition implies $h\nmid(N-1)/q$.
Thus $v_q(h)=v_q(N-1)$, so $q^{v_q(F)}\mid h\mid\ell-1$. Applying this
argument to every $q\mid F$ gives $F\mid\ell-1$. Hence $\ell\ge F+1$.
If $F$ is odd, then $\ell$ is odd and $\ell=uF+1$ with a positive even $u$,
so $\ell\ge2F+1$. A composite positive integer has a prime divisor at most
its square root, contradicting either strict bound in \eqref{eq:minus-bounds}.
\end{proof}

The gcd exponent is $(N-1)/q$, not $(N-1)/q^{v_q(F)}$. Repeated powers are
already accounted for by the order argument. The formulation above does not
require $\gcd(F,(N-1)/F)=1$; some textbook formulations impose this condition
to express the factored part using full prime powers. Different primes $q$
may have different witnesses. When a common witness works, its full Fermat
power need only be computed once.

\begin{corollary}[Proth certificate]\label{cor:proth}
Let $N=k2^s+1$, where $k$ is positive and odd and $k<2^s$. If
$a^{(N-1)/2}\equiv-1\pmod N$, then $N$ is prime.
\end{corollary}
\begin{proof}
Squaring gives the Fermat congruence. Since $N$ is odd, the reduced power minus
one has gcd $1$ with $N$. Apply \cref{thm:minus} with $F=2^s$; the inequality
$(F+1)^2>N$ follows from $k<F$.
\end{proof}

For the Mersenne seeds we also use the classical Lucas--Lehmer theorem:
if $r$ is an odd prime, $2^r-1$ is prime precisely when the recurrence
$s_0=4$, $s_{j+1}=s_j^2-2$ has $s_{r-2}\equiv0\pmod{2^r-1}$
\cite{lehmertest}. The small exponents used here are checked by trial division
before this criterion is applied.

\subsection{Lucas ranks from a factored part of \texorpdfstring{$N+1$}{N+1}}

Let $U_j(a,b)$ be defined by
\[
 U_0=0,\quad U_1=1,\quad U_{j+2}=aU_{j+1}-bU_j,
 \qquad D=a^2-4b.
\]
Here $a,b$ are Lucas parameters, not the capital-letter seed primes used later.

\begin{theorem}[Lucas order principle]\label{thm:plus}
Let $N>1$ be odd and let the fully factored integer $F>1$ divide $N+1$.
Fix a single pair of integers $(a,b)$, put $D=a^2-4b$, and suppose that,
for every distinct prime $q\mid F$,
\begin{equation}\label{eq:lucas}
 \gcd(N,2bD)=1,\qquad U_{N+1}(a,b)\equiv0\pmod N,
 \qquad \gcd\bigl(U_{(N+1)/q}(a,b),N\bigr)=1
 \quad(q\mid F).
\end{equation}
Every prime $\ell\mid N$ then satisfies
\[
 F\mid\ell-\left(\frac D\ell\right),\qquad
 \ell\equiv1\ \text{or}\ -1\pmod F.
\]
Thus $(F-1)^2>N$ suffices for primality. If $F$ is odd, the stronger sufficient
bound is $(2F-1)^2>N$.
\end{theorem}

\begin{proof}
For a prime $\ell\mid N$, let $\alpha,\beta$ be the distinct nonzero roots of
$X^2-aX+b$ in $\F_{\ell^2}$. The gcd hypothesis justifies
$U_j=(\alpha^j-\beta^j)/(\alpha-\beta)$. If $D$ is a square modulo $\ell$,
the ratio $\alpha/\beta$ lies in $\F_\ell^*$ and has order dividing $\ell-1$.
Otherwise Frobenius interchanges the roots, so this ratio has norm one and
order dividing $\ell+1$. In both cases, $U_j=0$ modulo $\ell$ exactly when
that order divides $j$. The same valuation argument as in
\cref{thm:minus} shows that $F$ divides this order. This proves the residue
restriction. If $F$ is odd, the coefficients in the possible forms $uF\pm1$
must be positive and even, giving $\ell\ge2F-1$. The square-root divisor
bound finishes the proof.
\end{proof}

A common discriminant matters when $F$ has several distinct prime factors:
the sign for a fixed divisor $\ell$ must be the same for all those factors.
Using one pair $(a,b)$ ensures this automatically. A Jacobi condition
$(D/N)=-1$ is helpful in searching for parameters and is checked in the
reported examples, but is not needed as an additional hypothesis in the
particular order statement above. A generic Lucas probable-prime pass without
the factor-by-factor gcd conditions is insufficient.

\subsection{Elliptic points of a known order}

\begin{definition}
A projective point modulo $N$ is \emph{strongly nonzero} if it reduces to a
nonidentity point modulo every prime divisor of $N$. In a projective
representation $(X:Y:Z)$, a sufficient condition is $\gcd(Z,N)=1$. An affine
point automatically has this property.
\end{definition}

\begin{theorem}[Known-order elliptic certificate]\label{thm:ec}
Let $N>1$ satisfy $\gcd(N,6)=1$, and let
$E:y^2=x^3+\alpha x+\beta$ have discriminant coprime to $N$.
Let $F>1$ be fully factored. Suppose that arithmetic valid modulo every prime
divisor of $N$ establishes
\[
 [F]Q=O,\qquad [F/q]Q\text{ is strongly nonzero for every prime }q\mid F.
\]
If $F>(N^{1/4}+1)^2$, then $N$ is prime. For a prime $F=P$, it is enough that
$Q$ be affine and $[P]Q=O$.
\end{theorem}

\begin{proof}
For every prime $\ell\mid N$, the order of the reduced point divides $F$ but
does not divide $F/q$ for any $q\mid F$. Its order is therefore exactly $F$.
Hasse's bound gives
\[
 F\le\#E(\F_\ell)\le(\sqrt\ell+1)^2.
\]
If $N$ were composite, choose $\ell\le\sqrt N$. The displayed inequality
contradicts the stated lower bound on $F$.
\end{proof}

The numerical certificates use the stronger, entirely integral comparison
\begin{equation}\label{eq:hasse-int}
 \bigl(\isqrt(P)-1\bigr)^4>N,
 \qquad \isqrt(P)=\lfloor\sqrt P\rfloor.
\end{equation}
Indeed, \eqref{eq:hasse-int} implies $N^{1/4}+1<\isqrt(P)\le\sqrt P$.
Every nontrivial denominator in the affine point additions is required to be
coprime to $N$. A failed inverse or an unresolved exceptional case rejects the
certificate. These checks ensure that the operations reduce to genuine group
operations over every $\F_\ell$; a nonzero-looking pair of residues alone does
not provide an elliptic primality proof. The general fully factored-order
version is classical; our three numerical elliptic outputs use prime orders.

\section{Arithmetic constructions and cubic extensions}
\label{sec:arithmetic}

\subsection{Linear, power, and mixed-product lifts}

For a known prime $P$, the candidates $cP+1$ and $cP^e+1$ expose the factored
parts $P$ and $P^e$ of $N-1$. When $P$ is odd, $c$ is taken even to make $N$
odd. Taking several certified inputs gives
\[
 N=c\prod_{i=1}^rP_i^{e_i}+1.
\]
The primality of the product is unnecessary: what matters is its complete
factorization. If that product is $F$, the elementary sufficient condition is
$F^2>N$, with the sharper bounds of \cref{thm:minus} available when useful.
The analogous negative candidates expose a known factor of $N+1$ and use
\cref{thm:plus}.

For fixed small $c$, the decimal size of $cP^e+1$ is approximately $e$ times
that of $P$. This is a statement about size, not about the probability of
primality. Increasing $e$ also increases the cost of the modular computations.
An arbitrarily high power can be specified without supplying a successful
certificate or an efficient search.

\begin{corollary}[Near-square lift]\label{cor:near}
Let $P$ be a known odd prime and $0<h<P$. Set $N=2P(P-h)+1$. If
\eqref{eq:pock} holds for $q=P$, then $N$ is prime.
\end{corollary}
\begin{proof}
Every prime divisor is at least $2P+1$, whereas
$N<2P^2+1<(2P+1)^2$.
\end{proof}

The known factor $P$ can be smaller than $\sqrt N$ in this construction.
Discarding parity information would miss a valid elementary certificate.

\subsection{A positive cubic specialization}

For an odd integer $F\ge3$, define
\begin{equation}\label{eq:delta}
 \delta(F)=\begin{cases}0,&F\equiv3\pmod4,\\2,&F\equiv1\pmod4.\end{cases}
\end{equation}

\begin{proposition}[Positive cubic lift]\label{prop:positive}
Let $F\ge3$ be odd and fully factored. Let $d$ be odd with $1\le d<F$, and set
\[
 N_+(F,d)=F^3-\delta(F)F^2-dF+1.
\]
If \eqref{eq:pock} holds for every prime $q\mid F$, then $N_+(F,d)$ is prime.
\end{proposition}

\begin{proof}
The integer $N=N_+(F,d)$ is odd and greater than one. By
\cref{thm:minus}, each of its prime divisors has the form $uF+1$ with $u\ge2$
even. Since $N<F^3<(2F+1)^3$, a composite $N$ would have exactly two prime
factors, counted with multiplicity. Write
\[
 N=(uF+1)(vF+1),\qquad u,v\ge2\text{ even}.
\]
The size bound yields $uv<F$. Also $u+v\le uv<F$, because
$(u-1)(v-1)\ge1$. Hence
\[
 \frac{N-1}{F}=uvF+(u+v)
              =(F-\delta(F)-1)F+(F-d).
\]
Both displayed remainders lie strictly between $0$ and $F$. Uniqueness of
Euclidean division gives $uv=F-\delta(F)-1\equiv2\pmod4$, contrary to
$4\mid uv$.
\end{proof}

This proof isolates a simple coefficient obstruction in the cubic range of
classical $N-1$ criteria. It does not replace those criteria with a new
fundamental test. The largest output of the first experiment uses a
\emph{composite} fully factored $F$, demonstrating that two input certificates
can genuinely be merged in this specialization.

\subsection{A negative cubic specialization}

\begin{proposition}[Negative cubic lift]\label{prop:negative}
Let $F\ge3$ be odd and fully factored, let $d$ be odd with $1\le d<F$, and put
\[
 N_-(F,d)=F^3-\delta(F)F^2-dF-1>1.
\]
If the common-parameter Lucas conditions \eqref{eq:lucas} hold for $F$, then
$N_-(F,d)$ is prime.
\end{proposition}

\begin{proof}
Every prime divisor of $N=N_-(F,d)$ is $1$ or $-1$ modulo $F$ and is at least
$2F-1$. Since $N<F^3<(2F-1)^3$, a composite $N$ has exactly two prime factors,
with multiplicity. Because $N\equiv-1\pmod F$ and $F\ge3$, the signs must be
opposite. Thus
\[
 N=(uF-1)(vF+1),\qquad u,v\ge2\text{ even}.
\]
Write $T=uv$ and $w=u-v$. Then $4\mid T$ and
$|w|\le T/2-2$. The latter follows, for instance when $u\ge v$, from
$uv/2-2-u+v=(v-2)(u/2+1)\ge0$. Therefore
\[
 E:=\frac{N+1}{F}=TF+w=F^2-\delta(F)F-d<F^2.
\]
If $T\ge F$, then $T\ge F+1$, since $T$ is even and $F$ odd. This gives
$E\ge T(F-1/2)+2>F^2$, a contradiction. Hence $T<F$ and $|w|<F$.
The quotient and remainder of $E$ upon division by $F$ are
$F-\delta(F)-1$ and $F-d$. If $w\ge0$, uniqueness gives
$T=F-\delta(F)-1\equiv2\pmod4$. If $w<0$, it gives
$T=F-\delta(F)$, an odd number. Both conclusions contradict $4\mid T$.
\end{proof}

The two cubic constructions require different witnesses: ordinary
multiplicative orders for $N_+$ and Lucas ranks for $N_-$. Their analogous
coefficient patterns do not imply that the two numbers are simultaneously
prime. We do not assert historical priority for either specialization.

\subsection{An exact square-test closure throughout the cubic range}

The coefficient patterns above are useful because a possible two-factor
decomposition is ruled out by parity. More generally, that last obstruction
can be decided by one integer square test. This gives a complete elementary
closure of the argument after the order witnesses have been established.

\begin{proposition}[Exact two-factor exclusion]\label{prop:square-closure}
Let $N>1$ and $F\ge3$ be odd, with $N<F^3$.
\begin{enumerate}[label=\textup{(\alph*)}]
\item Suppose $F\mid N-1$ and every prime divisor of $N$ is $1$ modulo $F$.
Write $(N-1)/F=hF+r$, $0\le r<F$. Then $N$ is composite if and only if
$r^2-4h=z^2\ge0$ and $(r-z)/2,(r+z)/2$ are positive even integers.
\item Suppose $F\mid N+1$ and every prime divisor of $N$ is $1$ or $-1$
modulo $F$. Use centered division
\[
 (N+1)/F=TF+w,\qquad -F/2<w<F/2.
\]
Then $N$ is composite if and only if $w^2+4T=z^2\ge0$ and
$(z+w)/2,(z-w)/2$ are positive even integers.
\end{enumerate}
The hypotheses on prime divisors can be supplied by
\cref{thm:minus,thm:plus}, respectively.
\end{proposition}

\begin{proof}
In case (a), the divisor lower bound and $N<F^3$ show that a composite has
exactly two prime factors. As in \cref{prop:positive}, write
$N=(uF+1)(vF+1)$ with $u,v\ge2$ even. Then $uv<F$ and $u+v<F$, so
$uv=h$ and $u+v=r$. These are precisely the roots of $X^2-rX+h$.
Conversely, the stated positive even roots give this factorization into
integers greater than one.

In case (b), the same size argument and the sign of $N$ modulo $F$ give
$N=(uF-1)(vF+1)$ with $u,v\ge2$ even. Put $T_0=uv$ and $w_0=u-v$.
Since $|w_0|\le T_0/2-2$ and $(N+1)/F\le F^2$, the inequality
$T_0\ge F$ would imply
\[
 (N+1)/F\ge(F+1)(F-1/2)+2>F^2,
\]
a contradiction. Hence $T_0<F$ and $|w_0|<F/2$. Centered-division uniqueness
gives $T=T_0$ and $w=w_0$. Thus $w^2+4T=(u+v)^2$, with the stated root
formulas. The converse again reconstructs the two factors exactly.
\end{proof}

For odd $F$ there is no tie in centered division: from ordinary division
$E=hF+r$, take $(T,w)=(h,r)$ if $2r<F$ and $(h+1,r-F)$ otherwise.
A square discriminant alone is not enough; positivity and evenness of the
reconstructed roots are essential. The two cubic specializations force the
potential product to be $2$ modulo $4$ or odd, immediately excluding these
roots. This is a convenient BLS-style arithmetic consequence, without a
claim of historical novelty.

\subsection{Combining both signs by the Chinese remainder theorem}

\begin{proposition}[Two-sided residue certificate]\label{prop:crt}
Let $F,G>1$ be coprime odd fully factored integers. Suppose that
$F\mid N-1$, $G\mid N+1$, and that the witnesses of
\cref{thm:minus,thm:plus} hold for these two parts. Define $L=2FG$ and let
$r$ be the unique integer with
\[
 1<r<L,\qquad r\equiv1\pmod{2F},\qquad r\equiv-1\pmod G.
\]
If $N<L^2$ and either $r^2>N$ or $r\nmid N$, then $N$ is prime.
\end{proposition}

\begin{proof}
Every prime divisor $\ell$ is $1$ modulo $2F$ and $\pm1$ modulo $G$.
Consequently $\ell\equiv1$ or $r\pmod L$. A prime divisor at most $\sqrt N$
is smaller than $L$. The representative $1$ is not prime, leaving only the
integer $r$. The last hypothesis excludes it.
\end{proof}

This criterion combines information without claiming that both halves are
always necessary. In the numerical example of \cref{sec:first}, the stronger
one-sided Lucas bound already proves primality. Both sets of witnesses were
nevertheless checked, so the example is a valid realization of the combined
construction, not an example of its necessity.

\section{Cyclotomic and repunit certificate transfer}
\label{sec:cyclotomic}

Cyclotomic identities can expose a previously proved prime as a factor of the
next candidate minus one. The logical benefit comes from that factorization,
not from the mere occurrence of a cyclotomic polynomial.

\begin{proposition}[A $\Phi_3$ to $\Phi_7$ transfer]\label{prop:phi}
Let $a\ge2$ be fully factored. Set
\[
 P=\Phi_3(a)=a^2+a+1,\qquad
 Q=\Phi_7(a)=a^6+a^5+a^4+a^3+a^2+a+1.
\]
The witnesses \eqref{eq:pock} for $F=a$ prove $P$ prime. Once $P$ is proved,
the witnesses for $F=aP$ prove $Q$ prime.
\end{proposition}

\begin{proof}
For the first stage, $P-1=a(a+1)$ and $(a+1)^2>P$, so
\cref{thm:minus} applies. Note carefully that $a^2<P$; the sharper integral
lower bound on a prime divisor is essential. For the second stage,
\[
 Q-1=a(a+1)\Phi_3(a)\Phi_6(a),\qquad
 (aP)^2-Q=a^5+2a^4+a^3-a-1>0.
\]
Thus the known factor $aP$ exceeds $\sqrt Q$.
\end{proof}

\begin{proposition}[Repunit transfer]\label{prop:repunit}
For integers $a,m\ge2$, define $R_m(a)=(a^m-1)/(a-1)$. If both $a$ and
$R_m(a)$ are fully factored, then
\[
 N=R_{2m+1}(a),\qquad F=aR_m(a)
\]
satisfy $F\mid N-1$ and $F^2>N$. Consequently valid $N-1$ witnesses for the
prime factors of $F$ certify $N$.
\end{proposition}

\begin{proof}
Direct factorization gives $N-1=aR_m(a)(a^m+1)$. Moreover,
\begin{equation}\label{eq:repdiff}
 (a-1)^2(F^2-N)
 =a^{m+2}(a^{m-1}-2)+a^2+a-1>0.
\end{equation}
The first term is nonnegative for $a,m\ge2$, and the remaining term is
positive. Apply \cref{thm:minus}.
\end{proof}

Neither proposition guarantees prime values. If $2m+1=uv$ with $u,v>1$, then
$R_u(a)$ is a proper divisor of $R_{uv}(a)$. Thus the tempting index iteration
$3\to7\to15$ stops at $15$. The longer pattern $5\to11\to23\to47\to95$
likewise stops at $95$. Prime indices are necessary, not sufficient, for prime
repunit values. A related valid but unexecuted transfer is
\[
 \Phi_5(a)-1=a(a+1)\Phi_4(a),\qquad \Phi_4(a)=a^2+1.
\]
If $a$ is even and fully factored and $\Phi_4(a)$ is certified prime, then
$F=a\Phi_4(a)$ is fully factored and
$F^2-\Phi_5(a)=a^6+a^4-a^3-a-1>0$ for $a\ge2$.

\section{Reverse norm constructions and elliptic certificates}
\label{sec:cm}

Suppose that a known prime $P$ is represented as $P=u^2+Dv^2$. The identity
\begin{equation}\label{eq:norm}
 N=Pt^2+2ut+1=(ut+1)^2+D(vt)^2
\end{equation}
suggests a curve-order target. In the quadratic ring, put
$\pi=1+t(u+v\sqrt{-D})$. Then
\[
 \Norm(\pi)=N,\qquad \Norm(\pi-1)=Pt^2,
 \qquad N+1-\operatorname{Tr}(\pi)=Pt^2.
\]
These relations motivate a CM search for a curve and a point $R$ whose
multiple $Q=[t^2]R$ has order $P$. They do not by themselves show that $N$ is
prime, that a particular twist has the predicted group order, or that $Q$ is
nonzero. The reported certificates finish with direct point arithmetic and
\cref{thm:ec}, independently of these predictions.

\subsection{Gaussian norms: discriminant \texorpdfstring{$-4$}{-4}}

Let
\begin{equation}\label{eq:C}
 C=401713\cdot2^{3300}+1.
\end{equation}
This 1,000-digit seed passes Proth's test with base $3$. Cornacchia reduction
provides positive integers $u,v$, with $u$ even and $v$ odd, such that
$C=u^2+v^2$. The complete integers are embedded in the verifier appendix.
For $t\equiv2\pmod4$, \eqref{eq:norm} gives $N\equiv5\pmod8$, convenient for
modular square-root extraction during the search. Four coefficient choices
$y^2=x^3-dx$, $d\in\{1,2,4,8\}$, were tried after a sieve and a Fermat
prefilter.

Two successful parameters were
\[
 t=3018,\qquad t=2^{1650}+94.
\]
The outputs have 1,006 and 1,993 digits, respectively. In both cases the
successful curve is
\[
 E:y^2=x^3-4x,
\]
and the supplied affine point $Q$ satisfies $[C]Q=O$, the smoothness test,
and \eqref{eq:hasse-int}. The original independent checks also verify the
recorded construction $Q=[t^2]R$; this extra construction check is unnecessary
for the final primality implication.

\subsection{\texorpdfstring{The norm $x^2+3y^2$: CM field discriminant $-3$}{The norm x squared plus 3y squared: CM field discriminant -3}}

Let $A=2^{1279}-1$. A representation $A=s^2+3r^2$ is obtained with $s>0$
even and $r>0$ odd. With
\[
 t=2^{638}+794,\qquad
 N=At^2+2st+1=(st+1)^2+3(rt)^2,
\]
the resulting 770-digit integer is certified on
$E:y^2=x^3+1$. Its explicit affine point has order $A$; the curve
discriminant is coprime to $N$ and the strict Hasse bound holds.
The search used affine arithmetic, whereas an independent verifier used
Jacobian coordinates and a different scalar-multiplication order, with the
necessary unit conditions checked.
Here $-3$ denotes the discriminant of the CM field $\mathbb Q(\sqrt{-3})$;
the displayed integral binary quadratic form $x^2+3y^2$ has discriminant $-12$.

\subsection{Extensions that have not been executed}

Other norm forms may be useful when a particular seed does not have the
desired representation. For example, discriminant $-8$ uses $u^2+2v^2$;
the curve $y^2=x^3-30x-56$ has $j=8000$. The 9,439-digit seed below is
$3$ modulo $8$, making this a relevant proposed direction. No prime based
on this proposal was searched for or certified in the present work.

Likewise, the fully factored order $F=\prod P_i^{e_i}$ in \cref{thm:ec}
allows several input certificates to be merged elliptically. This remains a
classical theoretical option here, not a fourth computed elliptic example.
Knowing a proposed curve order is insufficient: the required strongly
nonzero multiples must actually be exhibited.

\section{The numerical inventory and its provenance}
\label{sec:experiments}

\subsection{Foundational seeds}

Besides $C$ from \eqref{eq:C}, write
\[
 H=2^{521}-1,\qquad A=2^{1279}-1,\qquad B=2^{4423}-1,
 \qquad p_0=14^{512}-14^{256}+1.
\]
The first three are established Mersenne primes \cite{gimps}, but the
verification does not rely on a database lookup: it repeats their
Lucas--Lehmer tests. The complete list of foundational inputs is as follows.

\begin{center}
\begin{tabular}{@{}lrl@{}}
\toprule
Seed & Decimal digits & Independent certificate used here\\
\midrule
$H$ & 157 & Lucas--Lehmer, exponent $521$\\
$A$ & 386 & Lucas--Lehmer, exponent $1279$\\
$B$ & 1,332 & Lucas--Lehmer, exponent $4423$\\
$C$ & 1,000 & Proth, witness $3$\\
$p_0$ & 587 & $N-1$ witnesses for $2^{256}7^{256}$\\
\bottomrule
\end{tabular}
\end{center}

For $p_0$, put $w=14^{256}$. Then $p_0=w^2-w+1=\Phi_{1536}(14)$,
$p_0-1=w(w-1)$, and $w^2>p_0$. The witnesses are $a_2=11$ and $a_7=2$.
Both full power residues are $1$, and both required gcds are $1$. Base $2$
alone is not the advertised certificate for the $2$-part. This seed appeared
in a PrimeGrid discussion on 27 December 2014, including subsequent reports
of proof by existing software \cite{primegrid}. It is not presented as a newly
discovered seed.

\subsection{First experiment: ten outputs}
\label{sec:first}

\Cref{tab:first} preserves the ten outputs behind the earlier statement that
the largest constructed result had 4,153 digits. Every row is a proved prime.
The notation $T$ in row I5 is defined exactly by
\begin{equation}\label{eq:T}
 \begin{split}
 L&=2AC,\qquad r=1+2A\bigl((-A^{-1})\bmod C\bigr),\\
 T&=r+\bigl(\lfloor C/A\rfloor+964\bigr)L.
 \end{split}
\end{equation}
The inverse and the least nonnegative residue are taken modulo $C$.

\begin{table}[htbp]
\centering\small
\caption{The first ten certified outputs. Times are observed wall seconds;
search includes its successful certificate, while verification is a separate
standard-library Python execution.}\label{tab:first}
\begin{tabularx}{\linewidth}{@{}lXrrr@{}}
\toprule
ID & Exact integer & Digits & Search & Verify\\
\midrule
I1 & $776A^2-1$ & 773 & 0.384 & 0.540\\
I2 & $A2^{2158}+1$ & 1,035 & 1.158 & 0.084\\
I3 & $A^3-1971A+1$ & 1,156 & 1.501 & 0.202\\
I4 & $2482A^4+1$ & 1,544 & 4.515 & 0.543\\
I5 & $T$ of \eqref{eq:T} & 1,999 & 5.112 & 6.837\\
I6 & $722BC+1$ & 2,334 & 3.434 & 1.757\\
I7 & $2B(B-868)+1$ & 2,664 & 9.850 & 1.821\\
I8 & $C^3-2C^2-1315C+1$ & 2,998 & 9.157 & 2.975\\
I9 & $B^3-15751B+1$ & 3,995 & 211.692 & 7.122\\
I10 & $(AC)^3-6529AC+1$ & 4,153 & 88.080 & 13.183\\
\bottomrule
\end{tabularx}
\end{table}

The ordinary $N-1$ and positive cubic rows use base $2$ for every distinct
prime factor in their selected factored parts. Row I1 uses $F=A^2$ and Lucas
parameters $(a,b)=(3,1)$, so $D=5$. Row I2 uses the Proth witness $3$.
Row I5 uses base $2$ on the $A$-part of $T-1$ and Lucas parameters $(5,1)$,
with $D=21$, on the $C$-part of $T+1$.

For I10, set $F=AC\equiv3\pmod4$. The complete nontrivial witness conditions
are
\[
 2^{N-1}\equiv1\pmod N,\qquad
 \gcd(2^{(N-1)/A}-1,N)=\gcd(2^{(N-1)/C}-1,N)=1.
\]
Together with \cref{prop:positive}, they certify the 4,153-digit output.
The final successful certificate in the search took approximately $0.94$
seconds; locating that candidate took $88.08$ seconds. These costs answer
different questions.

For I5, the calculations establish $L^2>T$ and $r^2>T$, as required by
\cref{prop:crt}. They also establish $(2C-1)^2>T$. Thus the stronger one-sided
Lucas criterion already suffices, even though $A^2<T$ and $C^2<T$. The latter
two inequalities alone must not be used to claim an essential advantage of
the two-sided method.

\subsection{The squaring chain through 9,439 digits}
\label{sec:chain}

The earlier conversation supplied the following chain:
\begin{align*}
 p_1&=438p_0^2+1,& p_2&=13696p_1^2+1,\\
 p_3&=5494p_2^2+1,& p_4&=13258p_3^2+1.
\end{align*}
The respective digit counts are
\[
 587\longrightarrow1177\longrightarrow2357\longrightarrow4718
 \longrightarrow9439.
\]
For each step $N=cP^2+1$, the selected fully factored part is $F=P^2$;
$P^4>N$, and base $2$ satisfies
\begin{equation}\label{eq:chain-cert}
 2^{N-1}\equiv1\pmod N,\qquad \gcd(2^{cP}-1,N)=1.
\end{equation}
The gcd exponent is $cP$, not $c$. Each $P$ is proved before its child is
accepted. The last integer has 31,354 bits. Its SHA-256 identifier, computed
from its ASCII decimal expansion without a terminal newline, is
\begin{center}\footnotesize\ttfamily
82ddc0fab24f40e0c8c777c2f2e2138fec8\\
d9bc505e21eb414858111ca910089
\end{center}
The printed line break is not part of the digest.

The original search times of these coefficients are unavailable, and no
search-time inference is made from a later verification. The coefficients
are therefore attributed to the prior conversation, followed by explicit
independent reconstruction. In particular, the 9,439-digit result is not
accepted merely because it was reported in a chat message. It is included
in the full standalone verification in \cref{sec:reproduce}.

Absence from a selected prime database would not establish first discovery.
The present manuscript makes no global novelty claim for $p_3$, $p_4$, or
the other specific integers. The relevant achievement is a verifiable
certificate, not the historical priority of a decimal string.

\subsection{Second experiment: nine principal and two auxiliary outputs}
\label{sec:second}

Let $a_H=78299H$ and $a_A=48909A$. The multipliers factor as
$78299=13\cdot19\cdot317$ and $48909=3\cdot7\cdot17\cdot137$.
The second experiment gives the nine principal primes in
\cref{tab:second}. The symbols $u,s$ refer to the explicit norm
representations in \cref{sec:cm}.

\begin{table}[htbp]
\centering\small
\caption{Principal outputs of the second experiment. A time spanning two
rows is one shared search, not two independent timings.}\label{tab:second}
\begin{tabularx}{\linewidth}{@{}lXrl@{}}
\toprule
ID & Exact integer or parameter & Digits & Search seconds\\
\midrule
II1 & $\Phi_3(a_H)$ & 324 & $1.352$ (pair)\\
II2 & $\Phi_7(a_H)$ & 971 & same run\\
II3 & $\Phi_3(a_A)$ & 780 & $6.151$ (pair)\\
II4 & $\Phi_7(a_A)$ & 2,339 & same run\\
II5 & $Ct^2+2ut+1$, $t=3018$ & 1,006 & $2.734$\\
II6 & $Ct^2+2ut+1$, $t=2^{1650}+94$ & 1,993 & $4.351$\\
II7 & $At^2+2st+1$, $t=2^{638}+794$ & 770 & $0.394$\\
II8 & $C^3-2C^2-9053C-1$ & 2,998 & $71.453$\\
II9 & $2214p_3+1$ & 4,721 & $57.640$\\
\bottomrule
\end{tabularx}
\end{table}

Two additional first-stage primes arose during the cyclotomic searches:
\[
 \Phi_3(260H)\quad(319\text{ digits}),\qquad
 \Phi_3(4683A)\quad(778\text{ digits}).
\]
They are retained as auxiliary outputs, making eleven rather than nine
certified integers in this experiment. Their searches are part of the shared
runs in \cref{tab:second}; those times must not be multiplied by the number
of outputs.

For II1--II2 all factor witnesses are base $2$. For II3 every witness is
base $2$; in II4 the factor $3$ uses base $5$, while the remaining factors
use base $2$. In the first auxiliary, the factor $5$ uses base $3$; in the
second auxiliary, the factor $7$ uses base $5$. All other auxiliary
witnesses use base $2$.

The negative cubic II8 uses $(a,b)=(9,1)$ and $D=77$. Its verified conditions
are
\[
 \gcd(154,N)=1,\quad (77/N)=-1,\quad
 U_{N+1}(9,1)\equiv0\pmod N,\quad
 \gcd(U_{(N+1)/C}(9,1),N)=1.
\]
The additional Lucas proof on the successful candidate took about $0.715$
seconds. Independent generic-matrix Lucas verification took $27.469$
seconds. This output is different from the 2,998-digit positive cubic I8.

For II9, $p_3^2>N$ and the witness is base $2$:
$2^{N-1}\equiv1\pmod N$ and $\gcd(2^{2214}-1,N)=1$.
This is a modest increase over $p_3$, not an improvement on the size of $p_4$.
The two CM($-4$) searches tested 91 and 3 candidates surviving their sieves,
respectively; the CM($-3$) search tested 10. These favorable finite runs do
not imply constant expected search time.

\subsection{Search and verification protocol}

The principal search arithmetic used \texttt{gmpy2 2.3.2} with
\texttt{GMP 6.3.0}. Independent verifiers used Python 3.12.14 standard-library
integers, modular \texttt{pow}, and \texttt{math.gcd}. Lucas search code used
specialized doubling identities; independent verification used generic
$2\times2$ matrix powers. The CM($-3$) check also changed the point-coordinate
system and scalar-multiplication order. These are independent implementations,
not a claim of formal verification of Python, GMP, hardware, or the operating
system.

The search workflow is: specify a finite parameter window; sieve against
small primes; apply a cheap rejection test to survivors; attempt the full
certificate; then verify that certificate with a separate implementation.
The search does not promote a sieve survivor or a Fermat probable prime to
the status of a prime. Different routines include different setup costs in
their timers, and concurrent jobs shared resources. The reported wall times
are observations from individual runs, not a controlled benchmark of the
mathematical methods.

For perspective, independent checks of the three elliptic outputs took
$2.936$, $22.350$, and $0.707$ seconds, including their specified ancestor
and construction checks. The first packet's complete independent verification
took $35.327$ seconds. Summing every per-output time from different scripts
would double count some seed and chain checks.

\section{An implementation improvement: shared exponents}
\label{sec:optimization}

The proof obligations in \eqref{eq:chain-cert} need not be evaluated as two
unrelated exponentiations. More generally, for $N=cP^e+1$, define
\[
 z=a^{cP^{e-1}}\bmod N.
\]
Then $a^{N-1}\equiv z^P\pmod N$. The two exact checks are therefore
\begin{equation}\label{eq:shared}
 z^P\equiv1\pmod N,\qquad \gcd(z-1,N)=1.
\end{equation}
This is algebraically identical to the original certificate. It is
intermediate-result reuse, not a weaker test or a new primality theorem.
A ladder $a^c,(a^c)^P,\ldots$ supplies another implementation of the same
relations.

On the fixed 9,439-digit candidate $p_4$ and base $2$, one GMP measurement
gave
\begin{center}
\begin{tabular}{@{}lrr@{}}
\toprule
Implementation & Seconds & Sum of exponent bit lengths\\
\midrule
Two separate powers & 4.363427 & 47,038\\
Shared power of \eqref{eq:shared} & 3.330762 & 31,354\\
\bottomrule
\end{tabular}
\end{center}
The measured reduction is approximately $23.7\%$, or a factor of $1.310$ in
elapsed time. There was one repetition per method. Exponent bit-length sums
are a rough workload indicator, not an exact operation count; windows,
multiplication algorithms, and modular reduction affect runtime. In
particular, this measurement does not establish a $23.7\%$ reduction in
total search cost or a universal speedup on other hardware.

A complementary practical improvement is to retain a proof dependency graph
and cache each validated ancestor. The same prime should not be reproved for
every child in a batch. Multiple witnesses using one base can also share its
full Fermat residue. These optimizations reduce duplicated work while leaving
every mathematical proof obligation explicit.

\section{Record-scale inputs and stopped searches}
\label{sec:records}

\subsection{Mersenne coefficients inside Proth candidates}

Let $M_r=2^r-1$ and consider $N=M_r2^s+1$. If $s\ge r$, this is a Proth
number since its odd coefficient is smaller than $2^s$. If $r\ge3$ is odd,
odd $s$ makes $3\mid N$. For $r\ge3$ odd and $s\ge r$ even, one has
$N\equiv5\pmod{12}$.
Consequently, when $N$ is prime, quadratic reciprocity gives $(3/N)=-1$;
For $r\ge3$ odd and $s\ge r$ even, Euler's criterion and
\cref{cor:proth} show that
\[
 N\text{ is prime}\quad\Longleftrightarrow\quad
 3^{(N-1)/2}\equiv-1\pmod N.
\]
Thus the witness base is fixed for this family. The modular exponentiation
still has to be carried out.

If $M_r$ is already certified prime and $1\le s<r$, then $M_r^2>N$.
The same candidate can instead be certified using its known $M_r$ factor in
$N-1$. This gives a clean choice of proof mechanism on both sides of the
parameter boundary, not automatic primality at either side.

\begin{proposition}[A Fermat-number obstruction]\label{prop:fermat-filter}
Let $r>1$ be odd and set $h=v_2(r-1)$. For any positive $s$ with
$j=v_2(s)<h$,
\[
 2^{2^j}+1\mid (2^r-1)2^s+1.
\]
This divisor is proper, so the candidate is composite.
\end{proposition}
\begin{proof}
Modulo $F_j=2^{2^j}+1$, one has $2^{2^j}\equiv-1$ and
$2^{2^{j+1}}\equiv1$. Since $r\equiv1\pmod{2^{j+1}}$, we obtain
$2^r-1\equiv1$. Since $s/2^j$ is odd, $2^s\equiv-1$. The product plus
one vanishes. Finally $2^j<r$ and the candidate exceeds $2^r$, so the
displayed divisor is smaller than the candidate. The argument does not
assume that the Fermat number itself is prime.
\end{proof}

For $r=136279841$, one has $h=5$. Therefore every exponent $s$ not divisible
by $32$ is excluded by one of $3,5,17,257,65537$. The first exponent satisfying
both $s\ge r$ and $32\mid s$ is $136279872$. Surviving this exact filter
does not imply primality.

\subsection{What can be decided immediately at record scale}

The source snapshot consulted on 5 November 2026 records
\[
 M=2^{136279841}-1\quad(41{,}024{,}320\text{ digits}),\qquad
 R=10223\cdot2^{31172165}+1\quad(9{,}383{,}761\text{ digits})
\]
as a known Mersenne record and a known Proth record, respectively
\cite{gimps,prothrecord}. They are contextual external inputs, not additional
nodes in the locally reverified inventory. Small modular arithmetic proves
the following transformed candidates composite.

\begin{center}
\begin{tabular}{@{}lrlr@{}}
\toprule
Candidate & Proper divisor & Candidate & Proper divisor\\
\midrule
$2M+1$ & 3 & $2R+1$ & 5\\
$4M+1$ & 5 & $2R-1$ & 3\\
$M^2-M+1$ & 7 & $(R-1)^2+1$ & 397\\
$2M(M-1)+1$ & 2953 & $R^3-R+1$ & 4483\\
\bottomrule
\end{tabular}
\end{center}
The modular computations do not require writing the enormous decimal
integers. For example, $M\bmod\ell$ can be computed as
$2^{136279841}\bmod\ell-1$.

The previously known prime $7\cdot2^{20267500}+1$ \cite{primelist} fits the combined form
$(2^3-1)2^{20267500}+1$ and has 6,101,127 digits. Such an existing example
demonstrates that the family reaches large sizes, not that the construction
has produced a new record. Similarly,
\[
 2^{53689908}\bigl((2^{82589933}-1)+1\bigr)-1
 =2^{136279841}-1
\]
is an exact re-expression of an already known output prime.

\subsection{Bounded experiments that did not produce a prime}

A separate Proth search considered
$N_k=k2^{32000000}+1$ for odd $1\le k<2^{24}$. Among the $8{,}388{,}608$
initial parameters, sieving by primes at most $10^9$ eliminated $7{,}933{,}515$
and retained $455{,}093$. No record-sized primality proof finished. A short
PRST \cite{prst} run estimated approximately $12.17$ hours for one full test under its
measured settings; that was an extrapolation, not an observed completed test.
The attempted computation was stopped after $61{,}440$ of roughly $32$ million
iterations. It supplies no prime certificate. The sieve bound also makes no
statement about larger possible factors.

At the much smaller 9,439-digit seed $p_4$, bounded searches of
$2kp_4+1$ and $2kp_4-1$ tested 15 and 13 sieve survivors in approximately
$57.30$ and $41.86$ seconds, respectively, without obtaining a certified
output. These are finite unsuccessful attempts. They neither prove
compositeness of every member of either family nor prove that the chain
cannot be continued. All reported searches are completed or stopped; no
ongoing background search is implied.

\section{A constructive connection with Collatz prime trajectories}
\label{sec:collatz}

\subsection{Repository provenance and the usable construction}

The \emph{Verifiable Mathematical Reasoning Corpus} was consulted at the
following commit, dated 5 November 2026, 13:06:15 UTC:
\begin{center}\small\ttfamily
5936298ed4bacb7fb1423b889f02278fa521ea16
\end{center}
Its manuscript numbered [072],
\texttt{unified\_collatz\_recurrences.tex}, contains a useful affine
array and a discussion of primality-certificate obstacles in inverse
Mersenne orbits \cite{corpus}. The item numbered [048] has the same Git blob
and is a duplicate source, not independent confirmation. No theorem elsewhere
in the repository is needed to justify any numerical primality certificate
in this paper.

Define the standard Collatz map by $\mathcal C(n)=3n+1$ for odd $n$ and
$\mathcal C(n)=n/2$ for even $n$. The corpus array can be specialized so
that its orbit relation and a useful known factor of $N-1$ are designed
together.

\begin{proposition}[Certificate-compatible Collatz array]\label{prop:collatz}
Let $P>3$ be a certified prime, let $s,K,r$ be positive integers, and put
$F=P^s$. Suppose that
$2r4^K<F$. For $0\le i\le K$, define
\[
 N_i=2rF\,4^i3^{K-i}+1.
\]
Then $F\mid N_i-1$, $N_i<F^2$, and
\begin{equation}\label{eq:collatz-array}
 3N_i+1=4N_{i-1}\quad(1\le i\le K).
\end{equation}
For any index $i$, the $N-1$ witnesses for the single prime $P$ certify $N_i$
prime. If several indices pass, the corresponding primes lie on one genuine
Collatz trajectory in decreasing index order, with exactly three standard
steps per index difference.
\end{proposition}

\begin{proof}
Write $k_i=2r4^i3^{K-i}$. Then $k_i\le2r4^K<F$ and
$N_i=k_iF+1\le(F-1)F+1<F^2$. Apply \cref{thm:minus} to $F=P^s$.
Identity \eqref{eq:collatz-array} follows by direct expansion. Each $N_i$ is
odd, so the right side has exactly two factors of $2$. Thus the steps are
$N_i\mapsto4N_{i-1}\mapsto2N_{i-1}\mapsto N_{i-1}$, with no unverified
valuation assumption.
\end{proof}

This construction reuses $P$ as a certificate factor; it does not assert that
the trajectory ends at $P$. It also does not assert that all the $N_i$ are
prime. Choosing $s$ larger makes the finite certifiable range larger, while
increasing the size and cost of the candidates.

\subsection{An additional executed example: two 391-digit primes}

The specialization $P=A=2^{1279}-1$, $s=K=1$ was searched while preparing
this manuscript. At $r=98723$ it produced
\begin{equation}\label{eq:collatz-pair}
 \boxed{Q_0=592338A+1,\qquad Q_1=789784A+1.}
\end{equation}
Both integers have 391 decimal digits and both are prime. For each
$j\in\{0,1\}$, the verified certificate is
\[
 A^2>Q_j,\qquad 2^{Q_j-1}\equiv1\pmod{Q_j},\qquad
 \gcd\bigl(2^{(Q_j-1)/A}-1,Q_j\bigr)=1.
\]
The reduced exponents are the small integers $592338$ and $789784$.
Their precise orbit relation is
\[
 Q_1\longmapsto4Q_0\longmapsto2Q_0\longmapsto Q_0.
\]
The bounded search took approximately $1.1151$ seconds and examined 1,484
parameters surviving the joint sieve before success. The first candidate's
Fermat condition passed at 42 of those parameters; one pair passed both.
A separate standard-library Python implementation rechecked the seed by
Lucas--Lehmer, both exact certificates, hashes, sizes, and the orbit relation
in approximately $0.0148$ seconds. These are measured finite computations,
not expected-time estimates or claims of the smallest possible multiplier.

The baseline inventory is therefore
\[
 \underbrace{5}_{\text{foundational seeds}}
 +\underbrace{4}_{\text{prior chain lifts}}
 +\underbrace{10}_{\text{first experiment}}
 +\underbrace{11}_{\text{second experiment}}
 +\underbrace{2}_{\text{Collatz pair}}
 =32.
\]
The pair supplies a concrete use of the suggested corpus. Its importance is
the simultaneous orbit relation and checkable primality, not a new largest
prime or a result on the eventual behavior of all Collatz orbits.

\subsection{Why inverse ancestors of a fixed Mersenne seed are different}

\begin{proposition}[A one-step obstruction]\label{prop:inverse}
Let $r\ge3$ be an odd integer, put $M_r=2^r-1$, and let $a\ge1$ be an
integer. The odd integer
$N=(4^aM_r-1)/3$ has accelerated Collatz successor $M_r$, but
\[
 M_r\nmid N-1,\qquad M_r\nmid N+1,
 \qquad
 v_2(N-1)=\begin{cases}3,&a=1,\\2,&a\ge2.\end{cases}
\]
\end{proposition}
\begin{proof}
Since $M_r\equiv1\pmod3$, $N$ is an odd integer and
$3N+1=4^aM_r$ has exactly $2a$ factors of $2$. Moreover,
\[
 3(N-1)=4^aM_r-4,\qquad3(N+1)=4^aM_r+2.
\]
Divisibility by $M_r>4$ would force $M_r\mid4$ or $M_r\mid2$, respectively,
which is impossible. Finally
$v_2(N-1)=2+v_2(4^{a-1}M_r-1)$. This is $2$ for $a\ge2$ and $3$ for
$a=1$, since $v_2(M_r-1)=1$.
\end{proof}

The proposition explains a structural limitation of this particular inverse
orbit construction: neither the large seed factor nor a large power of two
is automatically available in $N-1$, and the seed factor is absent from
$N+1$ as well. It does not preclude other factorizations or an elliptic
certificate for the same candidate.

The corpus also discusses existence of long finite sequences of prime visits
using prime-distribution results for admissible linear forms, related to
Maynard's work \cite{maynard}. Those statements allow the starting and ending
primes to vary. They do not supply an arbitrary prescribed giant prime as
the endpoint, a search bound within $N<F^2$, or an infinite decreasing
sequence of positive integers. The finite construction and proof above do
not depend on such an existence theorem, nor do they address the Collatz
conjecture.

\section{The original Mills-constant motivation and its limits}
\label{sec:mills}

The investigation began with prime-representing constants. For an exponent
sequence $(C_k)$, write $\xi(C_k)$ for the least real $\xi>1$ such that
$\lfloor\xi^{C_k}\rfloor$ is prime for every relevant $k$, whenever that
minimum exists. Saito proved that Mills' constant $\xi(3^k)$ is irrational
\cite{saito-irr}. The version-v3 preprint on variants proves its
transcendence under the Density Hypothesis and gives unconditional arithmetic
results for several other exponent families \cite{saito-var}. These results
concern constants, not a numerical primality oracle.

An expression $3^k+n$ or $3^k-n$ is an integer exponent and is therefore
algebraic. A transcendence question in this context concerns the associated
\emph{real constant} $\xi(3^k\pm n)$, with positive indices and a well-defined
minimum, rather than the integer $3^k\pm n$. Different shifts normally
define different minimization problems. Neither the certificates in this
paper nor the act of combining four shifted constants establishes a new
transcendence result for any previously unresolved shift.

Even if a prime-representing constant is known abstractly, computing an
individual integer part requires adequate certified precision. For a rational
enclosure $L<\xi<U$, the equality
$\lfloor L^E\rfloor=\lfloor U^E\rfloor$ identifies that integer part.
An irrationality or transcendence theorem alone does not supply such an
enclosure or the computational resources to raise it to a huge exponent.

Nor is primality preserved by arbitrary combinations of prime outputs.
Products of two primes are composite; sums of two odd primes are even and
exceed two; a nontrivial power of a prime is composite. Towers built from odd
prime bases remain odd, so adding or subtracting one produces an even integer
larger than two in the nontrivial cases. Useful combinations must instead
carry a structural certificate, such as the mixed factored part $AC$,
the cyclotomic transfer, or a point of known elliptic order. That is the
precise connection between the original motivation and the methods developed
here.

\section{A catalogue of candidate designs}
\label{sec:catalogue}

\Cref{tab:catalogue} collects the diverse designs considered during the
investigation. It is a catalogue of constructions, not a count of independent
new primality theorems. Here $P$ and the $P_i$ are certified primes,
$F=\prod P_i^{e_i}$ is fully factored, and all candidates must be odd and
greater than one. Every successful row requires actual witnesses and its
applicable bound or exclusion argument.

{\small
\begin{longtable}{@{}>{\raggedright\arraybackslash}p{0.40\linewidth}>{\raggedright\arraybackslash}p{0.31\linewidth}>{\raggedright\arraybackslash}p{0.22\linewidth}@{}}
\caption{Candidate designs and the status of their use in this work.}
\label{tab:catalogue}\\
\toprule
Design & Certificate information & Status here\\
\midrule\endfirsthead
\toprule Design & Certificate information & Status here\\\midrule\endhead
\bottomrule\endfoot
$cP+1$ & Known factor $P\mid N-1$ & Executed: II9 and Collatz pair\\
$cP^e+1$ & Known prime power in $N-1$ & Executed: chain and I4\\
$c\prod P_i^{e_i}+1$ & Merge several certified factors & Executed: I6\\
$P^e2^s+1$, $P^e<2^s$ & Proth certificate; coefficient need not be prime & Executed special case I2\\
$(aP)^{2^r}+1$, $2\le a<P$ even, $r\ge0$ integral & Use $P^{2^r}\mid N-1$ & Valid proposal; not searched\\
$2P(P-h)+1$, $0<h<P$ & Parity-strengthened divisor bound & Executed: I7\\
$F^3-\delta F^2-dF+1$ & $N-1$ orders and cubic exclusion & Executed: I3, I8--I10\\
$cP^e-1$ & Known factor of $N+1$ & Executed: I1\\
$c\prod P_i^{e_i}-1$ & Common-discriminant Lucas ranks & Valid proposal; not searched\\
$P^e2^s-1$ & Appropriate Lucas/LLR conditions and bounds & Family considered; no new output\\
$F^3-\delta F^2-dF-1$ & Lucas ranks and cubic exclusion & Executed: II8\\
$\Phi_3(a)$, $\Phi_6(a)$ with factored $a$ & Factor $a$ in $N-1$; exact divisor bound & $\Phi_3$ executed; $\Phi_6$ seed\\
$R_m(a)\to R_{2m+1}(a)$ & Transfer factor $aR_m(a)$ & Executed: $\Phi_3\to\Phi_7$\\
$\Phi_4(a)\to\Phi_5(a)$ & Transfer factor $a\Phi_4(a)$ & Proved setup; not searched\\
$N\equiv1\pmod F$, $N\equiv-1\pmod G$ & Combine two residue restrictions & Executed: I5; one side also suffices\\
$Pt^2+2ut+1$, $P=u^2+v^2$ & Prime-order point on a $j=1728$ curve & Executed: II5--II6\\
$Pt^2+2ut+1$, $P=u^2+3v^2$ & Prime-order point on a $j=0$ curve & Executed: II7\\
Other CM norm forms, e.g. $u^2+2v^2$ & Appropriate curve, twists, and actual point & Proposed CM field discriminant $-8$\\
Elliptic fully factored order $F$ & Strongly nonzero $[F/q]Q$ for every $q$ & Classical theorem; not executed\\
$2rP^s4^i3^{K-i}+1$ & $N-1$ certificate and Collatz relation & Executed: two additional primes\\
Factor extraction from $2^P-1$ & A prime divisor has order $P$ & Larger factors guaranteed; extraction not executed\\
\end{longtable}}

For the generalized Fermat-shaped row, take an integer $r\ge0$, an even
integer $2\le a<P$, and put $e=2^r$ and $F=P^e$.
Then $F^2>(aP)^e+1$, so the factored-part setup is valid.
For the factor-extraction row, a prime divisor $\ell$ of $2^P-1$ has
$\ord_\ell(2)=P$, because $P$ is prime and order $1$ is impossible.
Thus $\ell\equiv1\pmod{2P}$ for odd $P$, and in particular $\ell>P$.
This guarantees larger prime divisors, but finding and certifying such a
divisor may be much harder than specifying the enormous input to be factored.

\section{Limitations and concrete further questions}
\label{sec:limits}

The finite results leave several substantive questions open. First, which
construction minimizes the cost per accepted prime at a prescribed output
size? Answering this requires repeated experiments with matched hardware,
parameter windows, sieve policies, and verification rules. The present
single-run timings cannot rank all methods universally.

Second, the norm approach could use other discriminants or products of
known prime orders. The difficulty is constructive: one must find an
appropriate candidate and an actual point with the required order. The
Hasse inequality alone does not create that point. Similarly, the
cyclotomic transfer graph is constrained by algebraic compositeness of
repunit indices; a formal iteration must be checked before launching an
expensive search.

Third, one can search for several certified prime entries in the same
Collatz-compatible array, perhaps with prescribed spacing. The certificate
theorem remains exact, but simultaneous primality in a particular bounded
array is not guaranteed. The existence of arbitrarily long finite prime
patterns in a varying family must not be read as a construction from a
fixed input in a fixed search interval.

Finally, no finite collection of successful $cP^2+1$ lifts proves that the
procedure continues indefinitely with a bounded multiplier. Dirichlet's
theorem on an unrestricted arithmetic progression does not by itself put a
prime in the particular finite range where the selected known factor is
large enough for the desired certificate. The examples show that rigorous
reuse works, and the formulas suggest further experiments; neither point
settles an infinitude conjecture for these specific families.

\section{Reproducibility, provenance, and the complete verification record}
\label{sec:reproduce}

The single TEX source contains every numerical datum needed to reproduce
the certificate verification: symbolic constructions of all 32 integers,
small witness bases and Lucas parameters, the two norm representations,
the three affine elliptic points, exact digit counts, and decimal SHA-256
identifiers. The long decimal expansions of the candidate integers are
unnecessary because their formulas are exact. The source compiles without
external figures, a bibliography file, or an input data file.

The verifier in \cref{app:code} uses only the Python standard library and
performs no network request or probable-prime acceptance. It first checks
formula reconstruction and metadata, then verifies the seeds and walks
the certificate dependency graph. All small factors are checked by trial
division. A named large factor is available only after its own certificate
has passed. Errors raise exceptions explicitly; the decisive conditions
are not implemented using assertions that disappear under Python's
optimization flag.

A final execution of the exact consolidated listing completed successfully
with status \texttt{ALL\_32\_PRIMES\_CERTIFIED}. It checked all 32 nodes in
approximately $138.746$ seconds with Python standard-library arithmetic,
including $p_4$ and the new Collatz pair. This is one observed verification
run; it is distinct from the historical search timings and from earlier
per-packet verification runs.

The original two computational packets were also audited for integrity:
their 31 and 57 manifest entries, respectively, matched the stored file
hashes. Their 32 decimal files represented 30 distinct integers because
$A$ and $C$ were repeated inputs. The extra Collatz pair is supplied directly
in the consolidated source. Hashes identify the exact intended integers;
they are not themselves primality evidence.

For historical attribution, the cyclotomic seed $p_0$ is public prior work,
the four chain coefficients were supplied in the earlier conversation, and
the first and second construction experiments were explicitly executed and
independently checked. The final Collatz pair was searched and certified
during manuscript preparation. No claim that any of these specific
integers was previously unknown is needed for the mathematical conclusions.

\subsection*{How to compile and verify}

Compile the one source with \texttt{pdflatex} twice, or use:
\begin{lstlisting}[language={}]
latexmk -pdf Certified_Prime_Lifting.tex
\end{lstlisting}
To recover the verifier,
copy the Python listing in \cref{app:code} into a file, or extract its marked
listing from the TEX source as follows:
\begin{lstlisting}
from pathlib import Path
import re
s = Path("Certified_Prime_Lifting.tex").read_text()
pat = (r"\\begin\{lstlisting\}\[label=lst:verifier\]\n"
       r"(.*?)\\end\{lstlisting\}")
m = re.search(pat, s, re.S)
if m is None:
    raise ValueError("Verifier listing was not found")
Path("verify_prime_lifting.py").write_text(m.group(1))
\end{lstlisting}
Run \texttt{python3 verify\_prime\_lifting.py} with Python 3.11 or later.
Success prints \texttt{ALL\_32\_PRIMES\_CERTIFIED}.

\clearpage
\section{Revision ledger and the continuing computation}
\label{sec:revision-ledger}

This revised edition retains the 32 certified integers of the original
inventory and adds eighteen explicitly specified algorithmic modules,
labelled A01--A18. The counts describe different objects: the baseline has
32 distinct prime integers and 21 rows of candidate designs; it did not
contain 32 independent primality algorithms. An additional algorithm need
not output a new prime: some modules construct candidates, some exclude
composites, some certify a supplied candidate, and some optimize or audit
the computation. Each module states its output contract.

\subsection{Two later certified extensions}

The later computation continued the chain of \cref{sec:chain} by
\[
 p_5=22330p_4^2+1,\qquad p_6=5936p_5+1.
\]
These integers have 18,882 and 18,885 decimal digits, respectively. Both
have base-2 Pocklington certificates. For $p_5$ the known factored part is
$p_4^2$; for $p_6$ it is $p_5$. In both cases its square exceeds the
candidate, the full Fermat residue is $1$, and the reduced-power gcd is
$1$. The seven-node chain starts from $p_0$ and the small primes $2,7$,
so no prime-table lookup is a hidden premise. Appendix
\ref{app:later-chain} includes an independently executable reconstruction
and exact verifier.

The decimal SHA-256 identifiers, without a trailing newline, are
\begin{quote}\small\ttfamily
p5: 7407dab961f68eb15854c74700a8007bf30d89410a6beeb1d2fc7d365f4b7ce8\\
p6: 1fd5cb8c8470b74c24c163501a2775693b19707db62c65e6b074a1e56ff7cbd3
\end{quote}
The identifiers select the intended integers; the order and gcd checks
prove primality. These are incorporated prior computations, not new
discoveries made by the eighteen-module experiment. The baseline inventory
and these two endpoints together contain 34 distinct certified integers.

\subsection{The interrupted 37,775-digit search}

The next bounded search used $N=2mp_6^2+1$ for
$1\le m\le10^6$. Sieving by primes through $10^8$ left 61,123 parameters;
extending the sieve through $10^9$ left 54,270. At interruption,
2,813 distinct completed Fermat-composite results were preserved in the
journal, and no new prime had been certified. Of these rejections, 2,683
belong to the final survivor set; the other 130 had subsequently been
removed by the deeper sieve. Thus 51,587 members of that set remained
without a preserved completed rejection.

The last live counter was 2,820, but its seven-entry excess was not present
in the durable journal. The journal count is used here. This discrepancy
does not promote any unrecorded computation to evidence. The stopped search
is not resumed by this revision. The implemented prototypes below are new,
bounded experiments with their own inputs and records.

\section{A common contract for the eighteen algorithms}
\label{sec:algorithm-contract}

\begin{definition}[Certificate-aware algorithm]
A certificate-aware algorithm has four explicit components:
an input domain, a finite procedure or a stated finite search window,
an output type, and a verification rule for that output. Arithmetic modules
return exact residues, and planning modules return plans with stated model
costs. Any claim about the primality of an integer is tagged as
\emph{candidate}, \emph{proved composite}, \emph{proved prime}, or
\emph{inconclusive}. An exhausted window is a statement about that window,
not a proof that the corresponding infinite family has no prime.
\end{definition}

A composite output includes an exact proper factor or a theorem-backed
contradiction such as a failed Fermat congruence. A prime output includes
all necessary factorizations, witnesses, size inequalities, and predecessor
certificates. A failed gcd check, a missing witness, a parameter-search
failure, or an unavailable elliptic inverse is not silently relabelled as
a proof of compositeness. If an inverse failure supplies a gcd strictly
between $1$ and $N$, that gcd itself is valid compositeness evidence.

\begin{definition}[Prime-divisor constraint]
A prime-divisor constraint for $N$ is a proved statement that every prime
divisor $\ell$ of $N$ lies in a specified set $\mathcal D$. Order
certificates, common-discriminant Lucas certificates, and valid elliptic
point calculations provide examples. Candidate-generation restrictions,
such as a decimal digit pattern, are a different type of condition and do
not by themselves restrict the prime divisors.
\end{definition}

\begin{proposition}[Sound composition of divisor constraints]
Let $N>1$ be an integer. Suppose independently checked certificates establish
$\ell\in\mathcal D_i$ for every prime $\ell\mid N$, for each
$i=1,\ldots,t$. If no integer in
\[
 [2,\lfloor\sqrt N\rfloor]\cap\bigcap_{i=1}^t\mathcal D_i
\]
divides $N$, then $N$ is prime.
\end{proposition}
\begin{proof}
A composite $N>1$ has a prime divisor at most $\sqrt N$. That divisor
would lie in every $\mathcal D_i$ and contradict the final exclusion.
\end{proof}

This elementary rule is the semantic foundation for combining certificates:
the constraints may be intersected only after their premises have been
verified. The intersection can be powerful when it has few representatives,
but enumerating a large intersection is not automatically efficient.

\subsection{Scope of the additional algorithms}

A01--A18 are additions to this manuscript with explicit algorithms,
correctness arguments, and reproducible demonstrations. Their ingredients
include established order tests, residue automata, Chinese remaindering,
product trees, exact integer reduction, finite optimization, and elliptic
groups. New notation or a new composition does not establish historical
priority. We distinguish a classical criterion, a constructed specialization,
an implementation rule, and an experimental scheduling policy. None of
these labels promises a faster asymptotic general-purpose primality test.

The purpose of the revision is to make stronger compositions concrete and
falsifiable. An exact operation-count improvement is reported as such.
Measured timing claims require matched inputs and repeated measurements.
An exact finite-population optimum is not extrapolated to an unseen prime
distribution, and a toy demonstration is not represented as a record-scale
benchmark.


\subsection{Catalogue of the eighteen additions}
\label{sec:new-catalogue}

The identifiers in Table~\ref{tab:new-catalogue} are stable throughout the
text and the executable appendices. They count procedures with different
input--output contracts, not eighteen interchangeable filters. The complete
hypotheses are stated in the corresponding subsections; in particular,
the quadratic and cubic constructions require a certified odd prime $p$
and a positive even coefficient $a$.

{\small
\begin{longtable}{@{}p{0.06\textwidth}p{0.41\textwidth}p{0.45\textwidth}@{}}
\caption{Eighteen additional algorithms and their principal contracts.}
\label{tab:new-catalogue}\\
\toprule
ID & Procedure & Output or decisive obligation\\
\midrule\endfirsthead
\toprule ID & Procedure & Output or decisive obligation\\
\midrule\endhead
\bottomrule\endfoot
A01 & Extended quadratic Fermat closure &
$N=ap^2+1$, $a<4p-2$; one full Fermat congruence and short nonunity.\\
A02 & Cubic trace exclusion &
$N=ap^3+1$; a coefficient-sized list of exact discriminant checks.\\
A03 & Joint positive and negative lifting &
Certify both $kp-1$ and $kp+1$ using a shared seed and separate proofs.\\
A04 & Inverse construction of digit blocks &
Prescribed prefix and suffix, with an adjustable middle block and known
factor of $N-1$.\\
A05 & Prime-power cyclotomic divisor mining &
Exact modular divisibility plus a universal lower bound on prime divisors.\\
A06 & Short-index negative quadratic closure &
$N=ap^2-1$, nonsquare $a<4p+2$; one long Lucas value and $\gcd(U_a,N)$.\\
\midrule
A07 & Sparse digit-language join &
Meet-in-the-middle generation satisfying both a language and a certificate
congruence.\\
A08 & Polynomial parameter root compiler &
Exact forbidden residue classes, audited by finite differences.\\
A09 & Two-tree batch sieve &
Route a small factor to each rejected candidate; retain unresolved cases.\\
A10 & Signed cofactor reduction &
Canonical modular products for $k2^n\pm1$ with at most one correction.\\
A11 & Common-order CRT packet &
One combined order certificate for several component integers.\\
A12 & Exact finite-population screening plan &
Minimum cost over ordered subsets in an explicitly specified finite model.\\
\midrule
A13 & Prime-power witness fusion &
Project separate witnesses to coprime primary orders and multiply them.\\
A14 & Batched unit verification &
One gcd on the successful path, with factor localization on failure.\\
A15 & Elliptic order fusion &
Project points to primary components and realize the lcm of known orders.\\
A16 & Residue--Hasse window intersection &
Combine universal arithmetic constraints with exact elliptic intervals.\\
A17 & Shared certificate planning &
Minimum-cost verified-set path; shared predecessors are counted once.\\
A18 & Hereditary divisor certification &
An exact divisor inherits a verified lower bound on all its prime factors.\\
\end{longtable}}

\subsection{A logical restriction on composition}

A universal divisor constraint and an existential one are different types
of evidence. For example, the Fermat pair used in A01 and A02 initially
shows that \emph{at least one} prime divisor is $1$ modulo $p$. The
coefficient argument then finishes the primality proof. The pair by itself
does not prove that \emph{every} prime divisor lies in that class, and must
not be supplied as a universal residue input to A16 or A18.

In contrast, the gcd conditions in the Pocklington order principle exclude
insufficient order at every prime divisor. Valid common-parameter Lucas
certificates and strongly nonzero elliptic multiples also provide universal
statements. A composition must retain this quantifier information, as well
as the numerical parameters. A short formula with the wrong quantifier is
not a certificate.


\section{Six construction algorithms and their exact acceptance rules}
\label{sec:new-constructions}

The six algorithms in this section add explicit acceptance rules and search
procedures to the earlier inventory. They are not six claims that a completely
new theory of primality has been discovered. A01, A02, and A06 give elementary
coefficient-sensitive refinements of order arguments; A03 combines two
different certificate mechanisms in a simultaneous search; A04 solves a
prescribed-block construction problem before testing primality; and A05 turns
a small divisor of a potentially enormous cyclotomic value into a prime
certificate. Their mathematical claims are proved below. Historical priority
for the refinements and combinations has not been established.

A prime supplied as an input always comes with an already checked proof.
The five larger numerical experiments use
\begin{equation}\label{eq:new-common-seed}
 P_\star=2^{127}-1
 =170141183460469231731687303715884105727.
\end{equation}
The verification listing first proves the exponent $127$ prime by trial
division and repeats all $125$ Lucas--Lehmer steps, obtaining residue zero.
This is a familiar 39-digit seed, not a newly discovered input. The small
seeds $3,5,11$ used in examples and controls are checked directly.

The output convention is deliberately three-valued. A valid certificate
returns \texttt{CERTIFIED}; a displayed factor or a failed necessary Fermat
condition returns \texttt{COMPOSITE}; a failed sufficient certificate returns
\texttt{INCONCLUSIVE}. Exhausting a finite search window is a fourth search
outcome, \texttt{NO\_CERTIFICATE\_IN\_WINDOW}, and does not establish that the
window contains no primes unless every entry was separately proved composite.

\subsection{A01: parity-strengthened quadratic order closure}
\label{sec:a01}

\paragraph{Input and purpose.}
The input is a certified odd prime $p$, a positive even coefficient $a$, and
a witness base $1<b<N$, where $N=ap^2+1$. The goal is to replace the long
reduced-exponent gcd of a conventional certificate by the short condition
$b^a\not\equiv1\pmod N$. For the range $a<p$, this principle appears in
\cite[Theorem 4.4]{philemon2021}. The following parity argument enlarges the
coefficient range to $a<4p-2$.

\begin{lemma}[Quadratic positive coefficient barrier]\label{lem:a01-barrier}
Let $p$ be an odd prime, let $a>0$ be even, and set $N=ap^2+1$.
If $N$ is composite and has a prime divisor congruent to $1$ modulo $p$,
then $a\ge4p-2$.
\end{lemma}
\begin{proof}
Choose such a divisor $\ell=up+1$, and write its complementary divisor as
$N/\ell=vp+1$, using $N\equiv1\pmod p$. Both divisors are odd and greater
than one, so $u,v\ge2$ are even. Comparison with $ap^2+1$ gives
\[
 a=uv+\frac{u+v}{p},\qquad p\mid u+v.
\]
The sum is even, hence $u+v=2pt$ for an integer $t\ge1$. Since
$(u-2)(v-2)\ge0$, we have $uv\ge2(u+v)-4$. Therefore
\[
 a\ge(4p+2)t-4\ge4p-2.
\]
The complementary divisor need not be prime; the argument includes repeated
prime factors and products with more than two prime factors.
\end{proof}

\begin{theorem}[A01 acceptance rule]\label{thm:a01}
Let $p$ be a certified odd prime and let $a$ be positive and even with
$a<4p-2$. Put $N=ap^2+1$. If
\begin{equation}\label{eq:a01-tests}
 b^{N-1}\equiv1\pmod N,
 \qquad b^a\not\equiv1\pmod N,
\end{equation}
then $N$ is prime. In particular, if $2^a<N$, then
\[
 N\text{ is prime}\quad\Longleftrightarrow\quad
 2^{N-1}\equiv1\pmod N.
\]
\end{theorem}
\begin{proof}
The first congruence makes $b$ a unit modulo $N$. Its multiplicative order
$h$ divides $ap^2$ but does not divide $a$, so $p\mid h\mid\phi(N)$.
Also $p\nmid N$. In the prime-power formula for $\phi(N)$, this forces at
least one prime divisor $\ell$ of $N$ to satisfy $p\mid\ell-1$.
If $N$ were composite, \cref{lem:a01-barrier} would contradict $a<4p-2$.
For the final statement, $1<2^a<N$ makes the second congruence automatic;
the converse is Fermat's little theorem.
\end{proof}

\paragraph{Executable procedure.}
Check the input proof and the strict coefficient bound. Compute
$f=b^{N-1}\bmod N$. A value $f\ne1$ proves compositeness because
$1<b<N$. If $f=1$, either use the exact integer inequality $b^a<N$ or
compute $s=b^a\bmod N$. Return \texttt{CERTIFIED} when $s\ne1$, and
\texttt{INCONCLUSIVE} when $s=1$. For base two and odd $N$, the inexpensive
test $a<\bits(N)$ is equivalent to $2^a<N$.

The acceptance rule uses one full modular power and at most one short power.
There is no unconditional promise that a chosen base will certify every prime
in the larger interval. This is not a general claim that a Fermat pass proves
primality outside the displayed family.

\paragraph{Sharpness and examples.}
The boundary admits the factorization
\[
 (4p-2)p^2+1=(2p+1)(2p^2-2p+1).
\]
At $p=3$, $a=10$, and $b=9$, the composite $91=7\cdot13$ satisfies
$9^{90}\equiv1\pmod{91}$ and $9^{10}\equiv9\not\equiv1\pmod{91}$.
Thus equality cannot be included in \cref{thm:a01}. The short test is also
essential: $N=55=6\cdot3^2+1$ and $b=54$ have both powers equal to one,
so the verifier correctly returns \texttt{INCONCLUSIVE}.

The small prime $151=6\cdot5^2+1$ illustrates the enlarged range $a\ge p$.
The executed larger search returned
\begin{equation}\label{eq:a01-example}
 N_{01}=180P_\star^2+1,
 \qquad \text{79 decimal digits},\qquad b=2.
\end{equation}
Its full residue is $1$, and its short residue is exactly $2^{180}<N_{01}$.
Eighteen candidates survived the small-prime sieve before this certificate
was found. This is an observed finite search count, not a distributional
claim.

\subsection{A02: cubic closure by a bounded trace scan}
\label{sec:a02}

\paragraph{Input and purpose.}
Let $p$ be a certified odd prime, $a>0$ even, and $N=ap^3+1$.
The order implication used in A01 still supplies a prime factor that is
$1$ modulo $p$, but the quadratic coefficient barrier no longer excludes
all composite values. A02 explicitly enumerates the small set of possible
factor traces. Unlike the earlier cubic closure in
\cref{prop:square-closure}, it does not begin by proving that every prime
factor is $1$ modulo a large fully factored $F$.

\begin{proposition}[Coefficient-sized cubic obstruction list]
\label{prop:a02}
Assume
\[
 b^{N-1}\equiv1\pmod N,
 \qquad b^a\not\equiv1\pmod N.
\]
Set
\[
 T=\left\lfloor\frac{ap+4}{2p+1}\right\rfloor.
\]
Then $N$ is composite if and only if, for some even integer $t$ with
$2\le t\le T$, the integer
\begin{equation}\label{eq:a02-discriminant}
 \Delta_t=(tp)^2-4(ap-t)
\end{equation}
is a nonnegative square $z^2$ and both
\[
 u=\frac{tp-z}{2},\qquad v=\frac{tp+z}{2}
\]
are positive even integers. In that case
$N=(up+1)(vp+1)$ is an explicit proper factorization.
\end{proposition}
\begin{proof}
As in the proof of A01, the two power conditions imply $p\mid\phi(N)$
and hence the existence of a prime factor $\ell\equiv1\pmod p$.
If $N$ is composite, write $N=(up+1)(vp+1)$ with even $u,v\ge2$.
Expanding and dividing by $p$ gives
\[
 ap^2=uvp+(u+v).
\]
Thus $u+v=tp$ with $t>0$ even, and $uv=ap-t$.
The inequality $uv\ge2(u+v)-4$ gives
$(2p+1)t\le ap+4$, which is the asserted finite bound.
The two coefficients are exactly the roots of
$X^2-tpX+(ap-t)$, giving \eqref{eq:a02-discriminant}.
Conversely, roots with the stated signs and parity reconstruct two
integers exceeding one whose product is $N$.
\end{proof}

\paragraph{Executable procedure and cost.}
First run the same full and short power checks as in A01. When they pass,
enumerate $t=2,4,\ldots,2\lfloor T/2\rfloor$. Compute the integer square
root of each nonnegative $\Delta_t$, and check the reconstructed factors
exactly. A factor gives \texttt{COMPOSITE}; absence of all such factors
gives \texttt{CERTIFIED}. The closure takes at most $\lfloor T/2\rfloor$
integer square tests, approximately $a/4$, rather than trial division up
to $\sqrt N$. These square tests can become expensive when $a$ grows;
no general speed advantage over a well-implemented Pocklington certificate
is inferred. The coefficient controls the extra work.

\paragraph{Executed example and adversarial control.}
The search returned
\begin{equation}\label{eq:a02-example}
 N_{02}=82P_\star^3+1,
 \qquad \text{117 decimal digits},\qquad b=2.
\end{equation}
Here $T=40$. The full power is $1$, the short power is $2^{82}<N_{02}$,
and none of the twenty even traces supplies a factor. Ten candidates
survived the preliminary sieve before this certificate was obtained.

The trace scan is a logical obligation, not a cosmetic extra check. At
$p=3$, $a=8$, $N=217$, and $b=191$, both power conditions pass, but
$t=4$ gives $\Delta_t=64$, $u=2$, $v=10$, and
$217=7\cdot31$. The executable verifier includes this example to detect
an accidental promotion of a mere order condition to a prime certificate.

\subsection{A03: simultaneous opposite-sign certificates}
\label{sec:a03}

\paragraph{Construction problem.}
Given a certified odd prime $p$, seek an even $k>0$ for which both
\[
 N_-=kp-1,\qquad N_+=kp+1
\]
are prime. This produces a certified pair at distance two, with the same
parent prime supplying a factor of $N_-+1$ and $N_+-1$.
It is a simultaneous finite search, not a proof of infinitely many twin
primes and not a promise of a pair for every fixed parent or search window.

\begin{proposition}[Joint certificate contract]\label{prop:a03}
Suppose $N_+<(2p-1)^2$. If a Pocklington witness for the factor $p$ proves
$N_+$ and common-parameter Lucas conditions for the factor $p$ prove $N_-$,
then both numbers are prime. Explicitly, sufficient conditions are
\begin{align*}
 b^{N_+-1}&\equiv1\pmod{N_+},&
 \gcd(b^k-1,N_+)&=1,\\
 U_{N_-+1}(\alpha,\beta)&\equiv0\pmod{N_-},&
 \gcd(U_k(\alpha,\beta),N_-)&=1,
\end{align*}
together with $\gcd(N_-,2\beta(\alpha^2-4\beta))=1$.
\end{proposition}
\begin{proof}
The two known factors are immediate from the definitions. The shared size
bound implies the parity-strengthened bounds in
\cref{thm:minus,thm:plus}. Apply those theorems separately.
\end{proof}

\paragraph{Algorithm.}
Choose a finite interval for $k$. Reject a parameter whenever either
$kp-1$ or $kp+1$ has a displayed proper small factor. For $p\ne3$ and
both candidates exceeding $3$, divisibility by three forces any surviving
even $k$ to be a multiple of six. On the survivors, attempt both exact
certificates and emit the pair only after both have passed. The input
certificate and small-parameter sieve can be shared; modular powers with
the two different moduli are still separate computations.

\paragraph{Executed 42-digit pair.}
With $p=P_\star$ and $k=1104$, the output is
\begin{align*}
 N_-&=187835866540358031831782783302336052722607,\\
 N_+&=187835866540358031831782783302336052722609.
\end{align*}
For $N_+$, base $b=2$ gives full residue $1$ and reduced gcd $1$.
For $N_-$, $(\alpha,\beta)=(3,1)$ gives discriminant $5$, full Lucas
residue $0$, and reduced gcd $1$. The parameter was reached after examining
184 multiples of six, of which 24 pairs survived trial division by the
odd primes below $101$. Neither integer is asserted to be historically new.
This joint output is different from the earlier Collatz-linked pair: its
specified relation is the exact difference $N_+-N_-=2$.

\subsection{A04: inverse construction with prescribed digit blocks}
\label{sec:a04}

\paragraph{Construction problem.}
Fix a radix $B\ge2$, a nonzero prefix $A$, a suffix $0\le C<B^m$ of
width $m$, and a free middle block $0\le X<B^t$. The candidate is
\begin{equation}\label{eq:a04-block}
 N=A B^{m+t}+X B^m+C.
\end{equation}
The leading zeros of the middle and suffix blocks, when present, count
toward their specified widths. The task is to enforce these digit constraints
while also making a known factor large enough to certify $N$.

\begin{proposition}[Certificate-compatible block splice]\label{prop:a04}
Let $F>1$ be fully factored and satisfy $\gcd(B,F)=1$. Set
\[
 K=A B^{m+t}+C,
 \qquad r\equiv(1-K)(B^m)^{-1}\pmod F,
 \qquad 0\le r<F.
\]
The middle blocks giving $F\mid N-1$ are exactly
\[
 X=r+jF\quad\text{with}\quad0\le r+jF<B^t.
\]
Restricting further to
\begin{equation}\label{eq:a04-limit}
 X\le H:=\min\left(B^t-1,
 \left\lfloor\frac{F^2-1-K}{B^m}\right\rfloor\right)
\end{equation}
guarantees $N<F^2$. For every odd $N>1$ in this restricted progression,
the ordinary Pocklington witness conditions for $F$ are therefore sufficient
to certify primality.
\end{proposition}
\begin{proof}
The congruence $K+XB^m\equiv1\pmod F$ has the unique residue class
displayed because $B^m$ is invertible modulo $F$. The digit restriction
is exactly the stated interval for $X$. The second upper bound is
equivalent to $K+XB^m\le F^2-1$. The final implication is
\cref{thm:minus}.
\end{proof}

\paragraph{Algorithm and failure states.}
Compute the one modular inverse, the residue $r$, and the clipped endpoint
$H$. If $H<r$, the admissible interval is empty. Otherwise enumerate
$j=0,\ldots,\lfloor(H-r)/F\rfloor$, or a user-supplied finite prefix
of this range. Reject even candidates, candidates no larger than the
designated input prime, and candidates with a displayed small factor.
Attempt a complete certificate on the survivors. A prescribed suffix that
forces divisibility by two or five in radix ten can make every candidate
composite; there is no existence guarantee for arbitrary blocks.

The construction is not independent primality testing of the digit blocks.
It solves the global congruence before generating them. It also differs
from merely appending zeros and one to a prime: both visible end blocks can
be fixed while the middle block is selected by an inverse congruence.
The number of admissible middle blocks is determined exactly by the
progression and interval; certificate attempts remain bounded by the chosen
search budget.

\paragraph{Executed 101-digit example.}
Take $B=10$, $A=314159$, $C=27183$, $m=5$, $t=90$, and
$F=P_\star^2$. In the notation above, the successful middle block is
$X=r+195F$. Equations \eqref{eq:a04-block}--\eqref{eq:a04-limit}, together
with the exact formula for $r$, specify the entire integer without an
external data file. It is a 101-digit prime beginning with
\texttt{314159} and ending with \texttt{27183}. Base two gives full
Pocklington residue $1$ and the required gcd $1$ for the prime $P_\star$.
Fifty-nine candidates survived the preliminary sieve through this successful
parameter. Its decimal SHA-256 identifier is
\begin{center}\ttfamily\scriptsize
136c7dd47d894a5a130e9d29463fba3cfc5407e1f7e27dcb9de347738cdf2ca0.
\end{center}
The hash identifies the reconstructed integer; the modular conditions prove
its primality.

\subsection{A05: cyclotomic divisor mining with a certified size threshold}
\label{sec:a05}

\paragraph{Input and purpose.}
Let $p$ be a certified odd prime, $t\ge1$, and $b\ge2$. The symbolic source
\[
 C=\Phi_{p^t}(b)
   =1+b^{p^{t-1}}+\cdots+b^{(p-1)p^{t-1}}
\]
may have far more digits than a prospective divisor $R$. A05 uses the
restricted prime factors of this source to turn a verified small divisor
into a proof of primality. It extends the earlier factor-extraction
observation by specifying a prime-power index, an essential exceptional-prime
guard, an explicit acceptance threshold, and a procedure that need not expand $C$.

\begin{theorem}[Cyclotomic divisor threshold]\label{thm:a05}
Suppose $1<R\mid\Phi_{p^t}(b)$, $\gcd(R,p)=1$, and
\begin{equation}\label{eq:a05-bound}
 R<(2p^t+1)^2.
\end{equation}
Then $R$ is prime. In particular, every accepted $R$ exceeds the input $p$.
\end{theorem}
\begin{proof}
Put $m=p^t$. The value $C$ is odd, whether $b$ is odd or even. For any
prime $\ell\mid R$, the constant term of $C$ shows that $\ell\nmid b$.
The identity $(b^{m/p}-1)C=b^m-1$ implies that the order of $b$ modulo
$\ell$ divides $m$. If it were a proper divisor of $m$, it would divide
$m/p$, and reduction of the geometric sum would give $C\equiv p\pmod\ell$.
That would force $\ell=p$, contrary to $\gcd(R,p)=1$. Thus
$\ord_\ell(b)=p^t$ and $p^t\mid\ell-1$. Since $\ell$ and $p^t$ are odd,
\[
 \ell\equiv1\pmod{2p^t},\qquad \ell\ge2p^t+1.
\]
A composite $R$ has at least two prime factors counted with multiplicity,
which would imply $R\ge(2p^t+1)^2$. This contradicts
\eqref{eq:a05-bound}.
\end{proof}

\paragraph{Division-free evaluation of the source.}
To verify $R\mid C$, first compute $x=b^{p^{t-1}}\bmod R$. For
$S_j(x)=1+x+\cdots+x^{j-1}$, maintain the pair
$(x^j,S_j(x))$ modulo $R$ using
\[
 x^{j+k}=x^j x^k,\qquad
 S_{j+k}(x)=S_j(x)+x^j S_k(x).
\]
Binary doubling evaluates $S_p(x)$ using $O(\log p)$ modular
multiplications, without division by $x-1$. A zero result proves the exact
divisibility $R\mid C$. This avoids the invalid assumption that $x-1$ has
an inverse modulo a still-unproved candidate. It also avoids constructing
the potentially enormous integer $C$.

\paragraph{Algorithm.}
Under a finite candidate budget, enumerate prospective
$R=2p^t j+1$ satisfying \eqref{eq:a05-bound}, or accept a supplied
prospective divisor from another calculation. Check the exceptional-prime
guard and the modular geometric sum. A passing $R$ is immediately a
certified prime by the theorem. A failed divisibility check rejects it as a
divisor of this source, not necessarily as a prime. If no candidate succeeds,
the algorithm makes no claim that the source has a divisor below the
threshold. Its verification cost is polynomial in the bit lengths of the
modular input and exponent, but the number of candidates required to find a
divisor is not bounded by this theorem.

\paragraph{Executed divisor certificates.}
For $(b,p,t)=(2,5,2)$, direct integer reconstruction gives
\[
 \Phi_{25}(2)=1082401=601\cdot1801.
\]
Both divisors are coprime to $5$ and smaller than $51^2=2601$.
Thus both $601$ and $1801$ are prime by the source-order argument.
For $(b,p,t)=(2,11,2)$ the same procedure certifies $727$ as a divisor of
$\Phi_{121}(2)$, because $727<243^2$. The implementation first checks
the modular geometric sums and then cross-checks these small source values
by direct integer evaluation; separate trial division confirms the three
small outputs.

The exceptional-prime guard cannot be deleted. Indeed,
\[
 \Phi_9(4)=4161=3\cdot19\cdot73,
 \qquad 57=3\cdot19<19^2,
\]
so $57$ is a composite divisor below the proposed threshold, excluded exactly
by $\gcd(57,3)\ne1$. Nor does the primality of all small divisors imply
primality of the entire cyclotomic value: $1082401$ itself is composite.

\subsection{A06: short-index Lucas closure for negative quadratic lifts}
\label{sec:a06}

\paragraph{Input and purpose.}
For $N=ap^2-1$, the conventional $N+1$ certificate for the known prime
power $p^2$ examines a reduced Lucas index $ap$. A06 instead uses the much
shorter index $a$, at the cost of an explicit nonsquare condition on $a$
and a coefficient bound. This is the negative-sign companion to the
positive coefficient barrier in A01.

\begin{lemma}[Quadratic negative coefficient barrier]\label{lem:a06-barrier}
Let $p$ be an odd prime, let $a>0$ be even, and let $N=ap^2-1$.
If $N$ is composite and has a prime divisor congruent to $1$ or $-1$
modulo $p$, then either $a$ is a perfect square or $a\ge4p+2$.
\end{lemma}
\begin{proof}
Choose a prime divisor $\ell=up+\varepsilon$, where
$\varepsilon\in\{1,-1\}$. Since $N\equiv-1\pmod p$, the complementary
divisor is $vp-\varepsilon$. Both are odd and greater than one, so
$u,v\ge2$ are even. Expanding gives
\[
 a=uv+\varepsilon\frac{v-u}{p},\qquad p\mid v-u.
\]
If $u=v$, then $a=u^2$. Otherwise, write
$|u-v|=2ph$ for an integer $h\ge1$. With $w=\min(u,v)\ge2$,
\[
 uv=w(w+2ph)\ge4+4ph.
\]
It follows that $a\ge uv-2h\ge4+(4p-2)h\ge4p+2$.
Again, the complementary divisor need not be prime.
\end{proof}

\begin{theorem}[A06 acceptance rule]\label{thm:a06}
Let $p$ be a certified odd prime. Let $a>0$ be even and not a square,
with $a<4p+2$, and set $N=ap^2-1$. Choose Lucas parameters
$(\alpha,\beta)$, put $D=\alpha^2-4\beta$, and suppose
\begin{equation}\label{eq:a06-tests}
 \gcd(N,2\beta D)=1,\qquad
 U_{N+1}(\alpha,\beta)\equiv0\pmod N,\qquad
 \gcd(U_a(\alpha,\beta),N)=1.
\end{equation}
Then $N$ is prime.
\end{theorem}
\begin{proof}
For a prime divisor $\ell$ of $N$, the rank-of-apparition argument in
\cref{thm:plus} applies. The order of the ratio of the two roots of
$X^2-\alpha X+\beta$ divides $ap^2$ and does not divide $a$.
It is therefore divisible by $p$, and consequently
\[
 p\mid\ell-\left(\frac D\ell\right),
 \qquad \ell\equiv1\ \text{or}\ -1\pmod p.
\]
If $N$ were composite, \cref{lem:a06-barrier} would force either a square
coefficient or $a\ge4p+2$, contrary to the hypotheses.
\end{proof}

\paragraph{Algorithm and cost.}
Check the parent proof, parity, nonsquare condition, and strict coefficient
bound. Check nondegeneracy, compute $U_a$, and take its gcd with $N$.
A proper gcd proves compositeness; a gcd equal to $N$ leaves the attempt
inconclusive. When the gcd is one, compute the full Lucas residue and accept
only if it is zero. A bounded search may try several parameter pairs;
the Jacobi condition $(D/N)=-1$ is useful for selecting them but is not an
additional hypothesis of the theorem.

The two indices are $a$ and $ap^2$, rather than $ap$ and $ap^2$.
This is a shorter explicit proof obligation. A well-designed implementation
can share index calculations, however, so the changed indices do not by
themselves establish a universal speedup or a factor-of-two improvement.
The number of Lucas doubling operations remains logarithmic in $N$ for
each fixed certificate attempt.

\paragraph{Executed example.}
The search returned
\begin{equation}\label{eq:a06-example}
 N_{06}=678P_\star^2-1,
 \qquad\text{80 decimal digits},
 \qquad(\alpha,\beta)=(4,1).
\end{equation}
Here $D=12$, the full Lucas residue is zero, and
$\gcd(U_{678}(4,1),N_{06})=1$. The coefficient is even and nonsquare,
and the size condition is immediate from $678<4P_\star+2$.
Eighty-four candidates survived the small-prime sieve through this output.

Both additional hypotheses have genuine counterexamples. If the nonsquare
condition is omitted, $p=3$, $a=4$, and $N=35$ pass all the Lucas conditions
with $(\alpha,\beta)=(1,1)$: $U_{36}=0$ and $U_4=-1$.
If equality is allowed in the coefficient bound, $p=3$, $a=14$, and
$N=125$ also pass with the same parameters: $U_{126}=0$ and $U_{14}=1$.
Thus neither deletion is justified by the modular conditions alone.

\subsection{Finite verification of the six constructions}
\label{sec:construction-verification}

The executable standard-library verifier reconstructs every displayed large
integer from \eqref{eq:new-common-seed}, proves that seed, and checks the
complete acceptance conditions. The six algorithms supply nine distinct
output integers in these experiments: the four single outputs of A01, A02,
A04, and A06 together with the two outputs of A03 and the three small outputs
of A05, counted correctly as $1+1+2+1+3+1=9$.
The common seed is a separate input. These outputs are a demonstration
inventory, not nine assertions of first discovery.

The finite controls compare A01 against exact trial division on 4,144
witness attempts, with 918 accepting attempts; A02 on 672 attempts, with
160 acceptances; and A06 on 1,981 attempts, with 244 acceptances.
All accepted integers in these small ranges are prime. Acceptances are
counted per witness and are not counts of distinct primes.
The generic matrix implementation of Lucas sequences is checked against
1,953 direct recurrence values, and the division-free geometric-sum
implementation is checked against 744 directly summed examples.
Seven hypothesis-sensitive negative controls include the composite examples
$91$, $55$, $217$, $35$, $125$, $57$, and $1082401$ described above.
These finite checks support implementation reliability; the proofs, rather
than the checked ranges, establish the general theorems.

As an additional check on the larger examples, the listing certifies A01
and A02 again with the classical Pocklington theorem, using the known
factors $P_\star^2$ and $P_\star^3$ respectively. It certifies A06 again
with the classical Lucas prime-power criterion, this time checking the
longer reduced index $678P_\star$. These are different certificate
implications applied to the same integers, not claims of independent
hardware or formally verified arithmetic libraries.

The complete executed listing returned
\texttt{A01\_A06\_ALL\_EXACT\_CHECKS\_PASSED}.
Its observed running time was below $0.1$ seconds in this environment,
including the small controls and the bounded example searches. These are
small examples with fast standard-library arithmetic, not a benchmark of
record-size searches and not evidence that a 117-digit result generalizes
to an instantaneous prime of millions of digits.

\section{Six exact arithmetic and search modules}
\label{sec:arith18}

Algorithms A07--A12 add a computational layer to the certificate constructions.
They have different contracts: A07 and A08 enumerate candidates, A09 produces
small-factor rejection certificates, A10 performs exact modular arithmetic,
A11 certifies a whole packet of integers, and A12 optimizes a stated finite
screening problem. Calling all six primality tests would obscure their
different outputs. The ingredients include classical meet-in-the-middle
search, polynomial congruences, product trees, special-modulus reduction,
Chinese remaindering, and dynamic programming. We claim neither historical
priority for those ingredients nor a measured improvement over optimized
GMP or PRST. The contribution here is their explicit integration with the
known-prime lifting contracts, with proofs and executable checks.

\begin{center}\small
\begin{tabularx}{\textwidth}{@{}lXl@{}}
\toprule
ID & Additional algorithm & Exact output\\
\midrule
A07 & Certificate-constrained digit-language join & Complete bounded candidate list\\
A08 & Quadratic residue compiler with finite-difference audit & Parameter exclusions and factors\\
A09 & Two-tree batch sieve with factor routing & Composite / small prime / unresolved\\
A10 & Signed cofactor reduction with one correction & Canonical modular product\\
A11 & CRT packet of common prime-power orders & Simultaneous primality certificate\\
A12 & Exact-cost finite-population screening planner & Optimal plan in a specified model\\
\bottomrule
\end{tabularx}
\end{center}

\subsection{A07: a certificate-constrained digit-language join}
\label{alg:A07}

\paragraph{Problem and input.}
Digit constraints can describe an alphabet, a fixed prefix, forbidden
adjacent pairs, or a regular language recognized by a finite automaton.
It is wasteful to enumerate such strings, turn them into arbitrary integers,
and only afterwards discover that their neighbors have no useful known
factor. Instead, impose the certificate congruence during generation.

Fix a base $B\ge2$, a length $d$, permitted digits at each position, and a
finite automaton $\mathcal A$ accepting the desired words. Leading zeroes are
excluded when a fixed digit length is intended. Let $F>1$ be fully factored.
The target set consists of the accepted integers $N$ satisfying
\begin{equation}\label{eq:A07-target}
 N\equiv1\pmod F,\qquad
 \max\mathcal Q<N<F^2,\qquad
 \gcd\left(N,\prod_{q\in\mathcal Q}q\right)=1,
\end{equation}
where $\mathcal Q$ is a finite set of proved small primes. Oddness is enforced
either by the digit language or by including $2$ in $\mathcal Q$.
The bound makes the elementary Pocklington certificate available, but
does not itself imply primality.

\paragraph{Sparse split.}
Split a word into a prefix $A$ of length $d-t$ and a suffix $C$ of length $t$.
Its integer value is $N=AB^t+C$. For each permitted suffix, store
\[
 \bigl(C\bmod F,\ C,\ T_C\bigr),\qquad
 T_C=\{s:\mathcal A\text{ accepts }C\text{ when started in state }s\}.
\]
The records are hashed by the first component. A prefix has a value $A$
and a resulting automaton state $s_A$. It only needs the bucket
\begin{equation}\label{eq:A07-join}
 C\bmod F=(1-AB^t)\bmod F,
\end{equation}
and within that bucket it retains the suffixes for which $s_A\in T_C$.
Finally apply the interval and small-prime conditions in
\eqref{eq:A07-target}. No array with one entry for every residue modulo
the potentially enormous integer $F$ is allocated.

\begin{proposition}[Completeness of the digit join]\label{prop:A07}
The sparse join produces exactly the integers specified by
\eqref{eq:A07-target}, each once for fixed-length digit words.
Every accepted output which subsequently passes the Pocklington conditions
for $F$ is prime.
\end{proposition}
\begin{proof}
Every permitted word has a unique prefix--suffix split. Acceptance of the
concatenation is equivalent to $s_A\in T_C$ by the definition of the
automaton transition. The congruence $AB^t+C\equiv1\pmod F$ is equivalent to
\eqref{eq:A07-join}. Thus every required word reaches exactly its own
matching bucket, and an incorrect word fails one of the explicit tests.
The last assertion follows from $F^2>N$ and \cref{thm:minus}.
\end{proof}

\paragraph{Procedure and cost.}
Enumerate the two half-word sets, build the suffix buckets, join the
matching residues, filter the emitted integers, and attempt a certificate.
With $U_L$ and $U_R$ half-words and $H$ records examined in matching buckets,
the combinatorial work is $O(U_L+U_R+H)$ hash operations, in addition to
digit-transition and integer-arithmetic costs. A suffix acceptance set can
be stored as a bitset over the automaton states. The bit lengths of the
residues still matter. In a balanced unrestricted alphabet of size $a$,
the half-word counts are about $a^{d/2}$, but the collision count $H$ can
be as large as $U_LU_R$. This is an output-sensitive exponential search,
not a polynomial-time method for every digit constraint. The split-and-join
idea is classical \cite{arith-hs}; adding the certificate congruence avoids
a full $F$-state residue automaton.

\paragraph{Executed example.}
Use decimal digits $\{1,3,7,9\}$, length four, and forbid equal adjacent
digits. Set $F=107$ and $\mathcal Q=\{3,5,7,11\}$; all these candidates
are already odd. There are $256$ unrestricted words and $108$ words in the
language. Splitting after two digits produces $12$ admissible prefixes and
$12$ admissible suffixes. Exactly two residue-bucket hits remain:
\[
 1713,\qquad7919.
\]
The first is divisible by $3$. For the second,
\[
 7919-1=74\cdot107,\qquad107^2>7919,\qquad
 2^{7918}\equiv1\pmod{7919},
\]
\[
 2^{74}\equiv6451\pmod{7919},\qquad\gcd(6450,7919)=1.
\]
Thus $7919$ is certified. Direct enumeration of all $256$ words gives
exactly the same candidate set; trial division independently confirms the
small numerical example. The example demonstrates constrained generation,
not a previously unknown prime.

\subsection{A08: a quadratic residue compiler with a finite-difference audit}
\label{alg:A08}

\paragraph{Input and purpose.}
The linear coefficient sieve for $kP^2+1$ extends to polynomial parameter
families, but a root of a polynomial modulo $q$ must be identified exactly.
Take a fully factored odd $F>1$ and
\begin{equation}\label{eq:A08-family}
 N(t)=F(2t^2+2t+2)+1,\qquad 0\le t\le T,
 \qquad 2T^2+2T+2<F.
\end{equation}
All candidates are odd, $F\mid N(t)-1$, and $N(t)<F^2$.
For every small prime $q$, compile the complete set of forbidden parameter
classes $R_q=\{t\bmod q:N(t)\equiv0\pmod q\}$. Each rejected parameter is
accompanied by the actual factor $q$.

\begin{proposition}[Two-class compiler]\label{prop:A08}
For the family \eqref{eq:A08-family}, $R_2$ is empty. If $q\mid F$ then
$R_q$ is empty. For any other odd prime $q$, put
\[
 \Delta_q=-3-2F^{-1}\pmod q.
\]
If $\Delta_q$ is a nonsquare, $R_q$ is empty. If $s^2=\Delta_q$ modulo $q$,
then
\begin{equation}\label{eq:A08-roots}
 R_q=\left\{\frac{-1+s}{2},\frac{-1-s}{2}\right\}\pmod q,
\end{equation}
with duplicates removed. In particular there are at most two forbidden
classes.
\end{proposition}
\begin{proof}
The cases $q=2$ and $q\mid F$ follow directly from \eqref{eq:A08-family}.
In the remaining cases,
\[
 2N(t)=F\bigl((2t+1)^2+3\bigr)+2.
\]
Both $F$ and $2$ are invertible modulo $q$, so vanishing of $N(t)$ is
equivalent to $(2t+1)^2=\Delta_q$. Solving this square equation gives
\eqref{eq:A08-roots}, including the double-root case.
\end{proof}

\paragraph{Exact reference algorithm.}
Compute $F\bmod q$, handle the two exceptional cases, find a modular square
root by a prime-field routine such as Tonelli--Shanks, and mark
$t\equiv r\pmod q$ for each $r\in R_q$. Verify a returned square root by
squaring it. For prime $q$, the Legendre-symbol computation also decides
whether a root exists. Modular square-root routines are standard
\cite{arith-flint-sqrt}; this compiler specializes their output to the
certificate-preserving parameter family.

An independent way to audit the root table is useful. For any integer
polynomial $f$ of degree $e$, let
$\delta_j(t)=\Delta^j f(t)\pmod q$, where $\Delta f(t)=f(t+1)-f(t)$.
The identities
\begin{equation}\label{eq:A08-difference}
 \delta_j(t+1)=\delta_j(t)+\delta_{j+1}(t)\pmod q
 \quad(0\le j<e),\qquad \delta_e(t+1)=\delta_e(t)
\end{equation}
generate successive values using additions. Initialize the differences from
$f(0),\ldots,f(e)$ and update from low to high index, so that each right-hand
side still uses the old value of $\delta_{j+1}$. Enumerating $t=0,\ldots,q-1$
then gives the complete root list. This identity is over the integers; it
does not divide by $e!$ and remains valid when $q\le e$.

\paragraph{Cost and boundary of usefulness.}
For the quadratic specialization, root discovery requires small-field
arithmetic and at most two progressions per prime. Marking a parameter
window of length $W$ costs
$O(\sum_q |R_q|(W/q+1))$ elementary marks, in addition to root discovery.
The reference finite-difference audit uses $eq$ additions per prime after
initialization. It is suitable for bounded tests and small moduli; it must
not replace the square-root compiler by enumeration of every residue when
$q$ is enormous. For a general polynomial or a short interval one may
instead generate only the required consecutive values. No assertion about
the number of prime outputs is made.

\paragraph{Executed example.}
Set $F=1009$ and $T=20$. The $21$ candidates have polynomial coefficients
\[
 N(t)=2019+2018t+2018t^2.
\]
The root compiler and the independent difference computation agree for
every prime $q\le31$. The union of the forbidden classes leaves exactly
$t=4$ and $t=10$. Both pass the exact certificate:
\begin{center}\small
\begin{tabular}{@{}rrrrr@{}}
\toprule
$t$ & $N(t)$ & $(N(t)-1)/1009$ & reduced base-$2$ residue & gcd\\
\midrule
4 & 42379 & 42 & 41835 & 1\\
10 & 223999 & 222 & 182964 & 1\\
\bottomrule
\end{tabular}
\end{center}
In both rows the full Fermat residue is $1$, and the displayed gcd is
$\gcd(\text{reduced residue}-1,N(t))$. Since $1009^2>N(t)$, these are
primality proofs. The difference audit uses $320$ modular additions,
excluding its initialization; additional tests include degree-six
polynomials modulo $2$ and $3$. Those checks specifically exercise
characteristics in which division by a polynomial degree factorial would
be invalid.

\subsection{A09: two-tree batch sieving with explicit factor routing}
\label{alg:A09}

\paragraph{Input and output.}
Let $N_1,\ldots,N_s>1$ be candidates and let
$\mathcal Q=\{q_1,\ldots,q_h\}$ be distinct proved small primes. Put
$W=\prod_j q_j$. The output for each $N_i$ is one of:
a proper factor, the fact that $N_i$ itself equals a known wheel prime,
or the status \textsc{unresolved}. The candidates need not be pairwise
coprime.

Build one balanced product tree on the candidates and a second on the wheel
primes. Starting from $W\bmod\prod_iN_i$, pass remainders down the candidate
tree. At leaf $i$ this gives $w_i=W\bmod N_i$. Compute
\[
 g_i=\gcd(w_i,N_i).
\]
If $g_i=1$, no wheel prime divides that candidate. Otherwise descend the
wheel tree, retaining nodes whose product has nontrivial gcd with $N_i$,
until a prime leaf $q\mid N_i$ is reached. Output $q$ as a compositeness
certificate only when $q<N_i$.

\begin{proposition}[Sound factor routing]\label{prop:A09}
The two-tree sieve rejects exactly the candidates having a proper divisor
in $\mathcal Q$. It never rejects a prime which itself belongs to
$\mathcal Q$.
\end{proposition}
\begin{proof}
At every edge of the candidate tree, the child product divides the parent
product. Reducing a parent remainder modulo that child therefore preserves
the remainder of $W$. Hence $g_i=\gcd(W,N_i)$. A prime divisor common to
$W$ and $N_i$ belongs to at least one child of any wheel-tree node containing
it, so descent reaches a leaf dividing $N_i$. If that leaf is smaller than
$N_i$, it is a proper divisor. If it equals $N_i$, the supplied primality
proof for the wheel leaf identifies a small prime, not a composite.
Conversely every proper wheel divisor gives $g_i>1$ and is found by descent.
\end{proof}

\paragraph{Why routing is part of the algorithm.}
The equality $g_i=N_i$ does not by itself decide compositeness. It can occur
both for a wheel prime and for a squarefree product of wheel primes.
Returning composite whenever the leaf gcd exceeds $1$ would incorrectly
reject the former. Returning unresolved whenever it equals $N_i$ would
discard useful factors in the latter. The explicit descent resolves this
ambiguity and supplies an independently checkable integer divisor.

\paragraph{Relation to prior methods and cost.}
Product and remainder trees, and batch gcd computations, are established
techniques \cite{cert-batchgcd,arith-flint-crt}. Here they are used for a
known wheel product rather than interpreted as a new primality criterion.
If the total candidate bit length is $S$, each balanced product-tree level
has total bit length $O(S)$ and there are $O(\log s)$ levels. Actual time
includes the reduction of the wheel product, remainders down the tree,
leaf gcds, and wheel-tree descent; the bit length of $W$ also counts.
Batching can share work but creates large intermediate products. A production
implementation should cap the total product size and use several batches.
The reference experiment establishes exact outputs and makes no speed claim
against an optimized segmented sieve.

\paragraph{Executed example.}
For primes at most $43$ the wheel product is
\[
 W=13082761331670030.
\]
The batch deliberately includes both equality cases and a composite with
no wheel factor:
\begin{center}\small
\begin{tabular}{@{}rrll@{}}
\toprule
$N_i$ & $\gcd(W,N_i)$ & routed leaf & exact output\\
\midrule
1713 & 3 & 3 & composite\\
7919 & 1 & -- & unresolved\\
341 & 341 & 11 & composite\\
49 & 7 & 7 & composite\\
17 & 17 & 17 & known small prime\\
10403 & 1 & -- & unresolved\\
42379 & 1 & -- & unresolved\\
223999 & 1 & -- & unresolved\\
\bottomrule
\end{tabular}
\end{center}
The surviving integer $10403=101\cdot103$ verifies concretely that a wheel
survivor cannot be promoted to a prime. The three other nontrivial survivors
have independent certificates elsewhere in this section; the sieve itself
does not use that knowledge.

\subsection{A10: signed cofactor reduction with one correction}
\label{alg:A10}

This arithmetic module isolates a useful contract for moduli
\[
 N=k2^n+\varepsilon,\qquad k\ge1,\quad n\ge1,\quad
 \varepsilon\in\{1,-1\},\quad N\ge3.
\]
It does not assume $N$ prime, and it does not require the Proth inequality.
When $k$ fits a machine word, it replaces reduction of a full product by
division by the large $N$ with a bit split and division by the fixed small
cofactor $k$. Special-modulus arithmetic and invariant-divisor techniques
are classical \cite{arith-gm}; the exact signed formula and its input range
are proved explicitly below.

\begin{proposition}[One-correction cofactor reduction]\label{prop:A10}
Let $R=2^n$ and suppose
\[
 0\le x\le(N-1)^2.
\]
Write
\[
 x=aR+b,\quad 0\le b<R,\qquad
 a=kt+s,\quad 0\le s<k.
\]
Define
\begin{equation}\label{eq:A10-reduce}
 y=b+sR-\varepsilon t.
\end{equation}
For $\varepsilon=1$, the canonical remainder $x\bmod N$ is $y$ if
$y\ge0$ and $y+N$ otherwise. For $\varepsilon=-1$, it is $y$ if $y<N$
and $y-N$ otherwise. At most one correction is required.
\end{proposition}
\begin{proof}
The identity
\[
 x-y=t(kR+\varepsilon)=tN
\]
proves congruence. Moreover $t=\lfloor x/(kR)\rfloor$.
If $\varepsilon=1$, then $kR=N-1$, so $0\le t\le N-1$ and
$0\le b+sR\le kR-1=N-2$. Consequently $-N<y<N$.
If $\varepsilon=-1$, then $kR=N+1$, so
$0\le t<N-1$ and $0\le b+sR\le kR-1=N$.
Thus $0\le y<2N$. The stated single correction gives the unique integer
in $[0,N)$ congruent to $x$.
\end{proof}

\paragraph{Procedure.}
For canonical input residues $0\le u,v<N$, set $x=uv$ and perform
\begin{lstlisting}
high = x >> n
low = x & ((1 << n) - 1)
t, s = divmod(high, k)
y = low + (s << n) - epsilon*t
if epsilon == 1 and y < 0:
    y += N
elif epsilon == -1 and y >= N:
    y -= N
return y
\end{lstlisting}
The product bound is automatic for these inputs, so this routine can replace
the reduction inside repeated squaring. There is no Montgomery-domain
conversion in this particular kernel. If $k$ is a machine word, division
of the multiword \texttt{high} by $k$ is a word-divisor operation; fixed
reciprocal techniques may be used \cite{arith-gm}. Its linear pass over the
limbs and the bit operations do not remove the cost of forming $uv$.
If $k$ is itself a large integer, the small-divisor advantage disappears.

\paragraph{The range restriction is essential.}
It is tempting to replace the hypothesis by $x<N^2$, but the stated
one-correction formula then fails for the plus sign. For $N=5$, $k=1$,
$R=4$, and $x=24<N^2$, it gives $y=-6$, and one addition of $N$ leaves
$-1$ instead of the correct remainder $4$. The executable reference
rejects such an out-of-contract input. This does not affect products of
canonical residues, for which $x\le16$ in this example.

\paragraph{Executed example and tests.}
For $R=256$, $k=5$, and $x=1000\cdot1200$, the common decomposition is
\[
 a=4687,\quad b=128,\quad t=937,\quad s=2.
\]
At $N=1281$, \eqref{eq:A10-reduce} gives $-297$, corrected to $984$.
At $N=1279$, it gives $1577$, corrected to $298$.
Both equal direct integer remainders. The reference script checks all
$193871$ canonical operand pairs for $1\le k\le9$, $1\le n\le4$, and
both signs whenever $N\ge3$. It also verifies $768$ fixed-seed random
products at radix sizes $64,256,1024,4096$ bits, and $128$ complete modular
powers against Python's independent built-in \texttt{pow}.
The random product choices exercise implementation arithmetic; they are not
probabilistic primality evidence. No elapsed-time superiority is claimed.

\subsection{A11: a CRT packet of common prime-power orders}
\label{alg:A11}

\paragraph{Objective.}
Several candidates built from the same certified prime have a common order
requirement. This permits an exact packet certificate: a verifier can prove
all component integers prime through powers modulo their product. The
larger modulus is a real cost, so the purpose is compositional verification
and a common arithmetic contract, not an automatic speedup.

\begin{theorem}[Common-order packet]\label{thm:A11}
Let $p$ be a certified odd prime and let $N_1,\ldots,N_s$ be integers with
\[
 1<N_i<p^4.
\]
Put $M=\prod_iN_i$. If an integer $Z$ satisfies
\begin{equation}\label{eq:A11-packet}
 Y=Z^p\bmod M,\qquad Y^p\equiv1\pmod M,\qquad
 \gcd(Y-1,M)=1,
\end{equation}
then every $N_i$ is prime.
\end{theorem}
\begin{proof}
Let $\ell$ be any prime divisor of an $N_i$. The full power condition
shows that $Z$ is a unit modulo $\ell$ whose order divides $p^2$.
The gcd condition says $Z^p\not\equiv1\pmod\ell$. Since $p$ is prime,
the order is exactly $p^2$, so $p^2\mid\ell-1$ and $\ell\ge p^2+1$.
A composite $N_i<p^4$ would have a prime divisor at most
$\sqrt{N_i}<p^2$, a contradiction.
\end{proof}

The theorem does not require the verifier to factor the $N_i-1$ or
individually run their complete certificates. It also does not assert that
the product $M$ is prime. In fact, for a packet with at least two components
it is deliberately composite.

\paragraph{Building a packet from the lifting family.}
A convenient construction has
\[
 N_i=k_i p^2+1,\qquad N_i<p^4.
\]
Choose provisional bases $2\le a_i<N_i$ and compute
$z_i=a_i^{k_i}\bmod N_i$. If the candidates are pairwise coprime, construct
the CRT representative
\[
 Z\equiv z_i\pmod{N_i}\quad\text{for every }i,\qquad0\le Z<M.
\]
The two powers in \eqref{eq:A11-packet} then simultaneously encode the local
order checks. Pairwise coprimality is required by this CRT construction,
although it is not a hypothesis of the verification theorem. A collision
must be handled by gcd and factor extraction; it cannot be ignored.
Standard implementations of CRT and product-tree reduction are described in
\cite{arith-flint-crt}.

If the packet fails, it is not valid to call every component composite.
Reduce the computed residues modulo the individual $N_i$. A failed full
Fermat condition for a base strictly between $1$ and $N_i$ proves that
component composite. A proper gcd supplies a factor. A reduced gcd equal
to $N_i$ can merely mean that the chosen base has insufficient order,
requiring a new witness. Valid components may be retained and repacked.

\paragraph{Executed certificate.}
Take $p=11$, bases $a_i=2$, and
\[
 (k_1,k_2,k_3,k_4)=(6,12,18,22).
\]
The component integers and their proposed $p^2$-order roots are
\begin{center}\small
\begin{tabular}{@{}rrr@{}}
\toprule
$k_i$ & $N_i=k_i11^2+1$ & $z_i=2^{k_i}\bmod N_i$\\
\midrule
6 & 727 & 64\\
12 & 1453 & 1190\\
18 & 2179 & 664\\
22 & 2663 & 79\\
\bottomrule
\end{tabular}
\end{center}
CRT gives the fully explicit certificate
\[
 M=6129547598087,\qquad Z=4290773793366,
\]
\[
 Y=Z^{11}\bmod M=4047432967480,\qquad
 Y^{11}\bmod M=1,\qquad\gcd(Y-1,M)=1.
\]
Together with $11$ prime and $N_i<11^4$, this proves all four component
integers prime. Direct trial division independently checks these small
examples. Negative controls replace one component by $243$, replace the
seed by $9$, use the witnesses $0$ and $1$, and violate the strict size
bound; all are rejected.

\paragraph{Operation counts and the cost of packing.}
For a positive exponent $e$, left-to-right binary powering initialized at
the leading bit uses $\bits(e)-1$ squarings and
$\operatorname{popcount}(e)-1$ multiplications. Two powers with exponent
$11$ therefore use six squarings and four multiplications. Four separate
order checks use $24$ squarings and $16$ multiplications, but their moduli
have only $10,11,12,12$ bits; the packet modulus has $43$ bits.
Under the illustrative squared-bit weight for schoolbook squaring, the
two costs are
\[
 6(10^2+11^2+12^2+12^2)=3054,\qquad6\cdot43^2=11094.
\]
Thus fewer operations are not evidence of faster arithmetic. More generally,
if $L_i=\bits(N_i)$ and $\mathcal M(L)$ denotes multiplication cost, the
powering terms scale as
\[
 O\!\left(\log p\sum_i\mathcal M(L_i)\right)
 \quad\text{and}\quad
 O\!\left(\log p\,\mathcal M\!\left(\sum_iL_i\right)\right),
\]
respectively, with CRT setup and gcds additional. Packing gives one
algebraic verification object, not a free reduction of total bit complexity.

\subsection{A12: an exact-cost finite-population screening planner}
\label{alg:A12}

\paragraph{Why a deterministic formulation is useful.}
An assertion that one candidate is most likely prime requires a probability
model and evidence supporting it. A different optimization is available
without such an assertion. For a specified finite set of candidates,
residue classes determine exactly which candidates a given sound
small-factor screen rejects. One can optimize the screening plan for that
finite population under an explicit cost model.

Let the candidate index set be $U$, let $R_j\subseteq U$ be the exactly known
rejection set of screen $j$, and let $c_j\ge0$ be its prescribed cost per
surviving candidate. Let $C\ge0$ be the terminal test cost per survivor.
The terminal procedure may still return composite or unresolved; it is
not assumed to accept every survivor. The objective is to process the
whole finite population, not to minimize the random time to the first
prime.

For a subset $A$ of screens already applied, define
\[
 S_A=U\setminus\bigcup_{j\in A}R_j.
\]
The recurrence
\begin{equation}\label{eq:A12-dp}
 V(A)=\min\left\{
 C|S_A|,\
 \min_{j\notin A}\bigl(c_j|S_A|+V(A\cup\{j\})\bigr)
 \right\}
\end{equation}
includes the option of stopping the filters immediately. Retain the
minimizing decision at each subset to reconstruct an order. A screen with
no useful additional rejection can be omitted.

\begin{proposition}[Optimality in the stated finite model]\label{prop:A12}
The value $V(\varnothing)$ is the minimum cost among every ordered subset
of the available screens followed by the terminal procedure.
\end{proposition}
\begin{proof}
After a set $A$ of screens, the survivor set is $S_A$, independently of
their previous order. A plan either terminates now at cost $C|S_A|$ or
chooses some next screen $j$, pays $c_j|S_A|$, and faces exactly the
subproblem indexed by $A\cup\{j\}$. Induction on the number of unused
screens proves \eqref{eq:A12-dp} and its optimality.
\end{proof}

\paragraph{Procedure, complexity, and accounting.}
Compute the exact rejection sets from residue classes, store them as
bitsets, evaluate the $2^h$ subsets for $h$ screens, and recover the
minimizing path. The planner is exponential in the number of available
screens and is intended for a small menu of screen groups, not for one
decision variable per prime below a billion. It uses no independence
assumption between screens. Candidate-dependent costs can be substituted
by replacing $c_j|S_A|$ and $C|S_A|$ with sums of the corresponding prescribed
costs over $S_A$.

The rejection sets are already known at planning time. If obtaining them
costs as much as running the screens, that expense is a real setup cost;
it must be added when evaluating a production implementation. In the
present experiment the cost objective deliberately describes only the
subsequent execution. No held-out or future-population speed claim follows
from optimizing these exact finite data. The optimal-substructure argument
is standard dynamic programming, applied here to certificate-search filters.

\paragraph{Executed example.}
Use the $21$ parameters of \eqref{eq:A08-family} with $F=1009$ and
$t=0,\ldots,20$. Consider the following prescribed abstract costs:
\begin{center}\small
\begin{tabular}{@{}lrrrrr@{}}
\toprule
screen prime $q$ & 17 & 3 & 13 & 5 & 7\\
cost per survivor & 5 & 1 & 4 & 1 & 2\\
initial rejection count & 3 & 14 & 2 & 8 & 6\\
\bottomrule
\end{tabular}
\end{center}
Let the terminal cost be $40$. Applying all screens in the displayed
order costs $279$ units. The optimal plan applies the screens
$3,5,7$, leaves parameters $4,7,10$, and costs
\[
 21+7+8+3\cdot40=156
\]
units. There are $32$ subset states. Enumerating all $326$ ordered subsets
independently confirms that no plan has lower cost in this model.
These are exact modeled costs, not measured seconds, and the survivor
$t=7$ is not prime: the omitted screen $11$ would eliminate it in A08.
A smaller menu can be optimal under its own costs without becoming a
complete primality test.

\subsection{Reproducibility and what these six experiments establish}
\label{sec:arith18-verification}

The standard-library reference script in Appendix~\ref{app:arithmetic}
implements every module and raises an
exception if any required identity fails. It contains the complete small
inputs, digit sets, wheel primes, CRT packet, cost vectors, and fixed
pseudorandom seed for arithmetic stress tests. No large integer is accepted
on a probable-prime pass. The new prime claims in this section are backed
by explicit order certificates; trial division is an additional check on
the small examples.

The decisive reference output is
\[
 \texttt{ALL\_SIX\_ARITHMETIC\_ALGORITHMS\_VERIFIED}.
\]
Its meaning is finite and specific: the digit join agrees with exhaustive
enumeration, the quadratic root compiler agrees with finite differences,
every routed divisor is exact, the signed reduction satisfies its range
contract on the stated tests, the packet satisfies its order equations,
and the screening planner agrees with exhaustive optimization.
It does not establish universal performance rankings, a new giant-prime
record, or an unconditional supply of prime values in any bounded family.


\section{Certificate synthesis, hybrid bounds, and proof reuse}
\label{sec:cert-synthesis}

Algorithms A13--A18 operate on evidence as well as on integers.  An order
certificate can establish a statement about every prime divisor of an
integer without establishing that the integer itself is prime.  Keeping
these two conclusions separate permits useful compositions: several weak
certificates can give a stronger one, and a composite auxiliary integer can
sometimes certify one of its own divisors.  Every prime factor used in an
order must already be certified, and all modular computations appearing in
a certificate must actually be checked.

The algebraic ingredients are classical: multiplicative orders,
primary decomposition in finite abelian groups, Hasse's bound, products
and gcds, and shortest paths.  The additions below specify concrete
transformations, exact closure rules, and reproducible demonstrations.
They are additions to this manuscript's algorithm catalogue; historical
priority for these combinations is not asserted.  In particular, they are
not six independent replacements for the classical primality criteria.

\subsection{A13: prime-power witness projection and fusion}
\label{alg:A13}

\paragraph{Purpose and input.}
A conventional $N-1$ certificate may use a different witness for each
prime factor of its known part.  This algorithm combines those witnesses
into one element whose order is a fully factored integer.  Input an odd
$N>1$, a nonempty set $S$ of certified primes dividing $E=N-1$, and witnesses
$a_q$ satisfying
\[
 a_q^E\equiv1\pmod N,\qquad
 \gcd(a_q^{E/q}-1,N)=1\quad(q\in S).
\]
The cofactor of $E$ outside $S$ need not be factored.  Repeated exact division
by each supplied $q$ determines $e_q=v_q(E)$.

\begin{proposition}[A fused factored-order witness]\label{prop:cert-fusion}
Set
\[
 Q_q=q^{e_q},\qquad H=\prod_{q\in S}Q_q,\qquad
 z_q=a_q^{E/Q_q}\bmod N,\qquad z=\prod_{q\in S}z_q\bmod N.
\]
For every prime $\ell\mid N$, the reduction of $z$ has order exactly $H$.
Consequently
\[
 z^H\equiv1\pmod N,\qquad
 \gcd(z^{H/q}-1,N)=1\quad(q\in S).
\]
Every prime divisor of $N$ is at least
\[
 \lambda(H)=\begin{cases}H+1,&H\text{ even},\\2H+1,&H\text{ odd},\end{cases}
\]
and $\lambda(H)^2>N$ proves that $N$ is prime.
\end{proposition}

\begin{proof}
Fix $\ell\mid N$.  The full-power and gcd conditions show that
$h_q=\ord_\ell(a_q)$ divides $E$ and has $q$-valuation exactly $e_q$.
Raising $a_q$ to $E/Q_q$ removes every other primary part of $h_q$, so
$\ord_\ell(z_q)=Q_q$.  The orders $Q_q$ are pairwise coprime.  If
$z^u=1$, raising this identity to $H/Q_q$ kills every component except
$z_q$; since $H/Q_q$ is coprime to $Q_q$, it follows that $Q_q\mid u$.
Also $z^H=1$ modulo $\ell$.  Thus $\ord_\ell(z)=H$.  The stronger
annihilation congruence modulo $N$ follows directly from
$z_q^{Q_q}=a_q^E\equiv1\pmod N$, since $Q_q\mid H$; it is not inferred
merely by checking prime divisors.  The reduced powers are nonidentity
modulo every $\ell\mid N$, which gives the stated gcds.  Finally,
$H\mid\ell-1$ and oddness of $N$ give the bound $\lambda(H)$.
\end{proof}

\paragraph{Algorithm.}
Check all input witnesses and supplied prime factors; compute the
$Q_q$ and projections; multiply the projections modulo $N$; output
$(N,H,z)$ with the factorization of $H$.  A separate verifier can discard
the original witnesses and check the displayed $H$-power and gcd
conditions directly.  It accepts primality only when the size bound also
holds.  Otherwise the output remains a valid partial order certificate.

\paragraph{Exact examples.}
For $N=31$, use $a_2=3$ and $a_3=21$.  Multiplying these witnesses directly
gives $3\cdot21\equiv1\pmod{31}$ and loses all useful order.  Projection
instead gives
\[
 H=6,\quad z_2=3^{15}\equiv30,\quad z_3=21^{10}\equiv5,
 \quad z\equiv26\pmod{31}.
\]
The element $26$ has order $6$ modulo every prime divisor, and
$7^2>31$ closes the certificate.  A second example is
\[
\begin{gathered}
 N=1009,\quad (Q_2,Q_3)=(16,9),\quad(a_2,a_3)=(11,2),\\
 (z_2,z_3)=(179,759),\qquad(H,z)=(144,655).
\end{gathered}
\]
Both reduced gcds are $1$ and $145^2>1009$.

\paragraph{Cost and limits.}
After the input checks, construction uses $|S|$ modular powers with
exponents $E/Q_q$ and $|S|-1$ modular multiplications.  The compact
certificate stores one witness instead of one per prime.  Verification
uses exponent $H$ and its factor quotients; intermediate powers may be
shared.  A smaller numerical exponent is not a proportional speedup,
because exponentiation cost depends on its bit length.  There is no
universal running-time advantage, and this procedure does not manufacture
valid input witnesses.  The bounded verifier independently factors the
small moduli in 86 nonempty test cases below $400$ and checks the resulting
orders at each prime divisor, including cases where the final size bound
does not hold.

\subsection{A14: batched unit checks with exact failure localization}
\label{alg:A14}

\paragraph{Purpose and input.}
Pocklington reduced powers and elliptic denominators produce a collection
of conditions $\gcd(u_i,N)=1$ for one fixed modulus $N$.  The values $u_i$
must already be the correctly computed quantities required by the proof.
This algorithm batches their unit checks; it does not batch or omit the
underlying modular powers.  Product-tree methods are classical tools in
batch arithmetic; see, for example, the related constructions in
\cite{cert-batchgcd}.

\begin{proposition}[Unit aggregation and diagnosis]\label{prop:cert-unit}
For $N>1$ and a finite nonempty collection $u_1,\ldots,u_s$,
\[
 \gcd\left(\prod_{i=1}^s u_i,N\right)=1
 \quad\Longleftrightarrow\quad
 \gcd(u_i,N)=1\quad\hbox{for every }i.
\]
All products can be reduced modulo $N$.  A product tree permits exact
localization of every failed leaf by descending only into nodes whose
product has gcd greater than $1$ with $N$.
\end{proposition}

\begin{proof}
A prime divisor of $N$ divides the product precisely when it divides at
least one factor.  Hence a gcd-$1$ subtree contains no failed leaf.
Recursively splitting every other subtree reaches precisely the indices
with nonunit $u_i$.  Reduction modulo $N$ does not change any of these gcds.
\end{proof}

\paragraph{Algorithm and output labels.}
Build a balanced tree of products modulo $N$.  Compute the root gcd.  If
it is $1$, return \texttt{ALL\_UNITS}.  Otherwise descend into the affected
subtrees and compute leaf gcds.  A leaf gcd $d$ with $1<d<N$ is a proper
factor and proves compositeness.  A leaf gcd equal to $N$ only invalidates
that unit obligation: it can arise from an unsuitable witness even when
$N$ is prime.  If no proper factor is exposed, return
\texttt{NONUNIT\_WITHOUT\_PROPER\_FACTOR}, not a compositeness assertion.

\paragraph{Exact examples.}
For the compact witness $(N,H,z)=(1009,144,655)$ in A13, the reduced
quantities are $1007$ and $373$.  One gcd of their modular product proves
both are units.  In contrast, for
\[
 N=91,\qquad (u_1,u_2,u_3,u_4)=(2,7,13,5),
\]
the root gcd is the uninformative value $91$; tree descent exposes the
proper factors $7$ and $13$.  For the prime modulus $101$, the collection
$(0,2,3,4)$ has a nonunit leaf but yields no proper factor.  This last case
tests the essential distinction between a failed certificate and a
composite integer.

\paragraph{Cost and limits.}
The successful path uses $s-1$ modular multiplications and exactly one
gcd, replacing $s$ separate gcds.  Full failure localization uses at most
$2s-1$ gcds, so the unsuccessful path need not be faster.  No wall-clock
speedup is claimed from these operation counts.  The verifier compares all
diagnosed leaves with separate gcds in 592 bounded batches and checks the
three output states above.  This is distinct from trial-dividing many
candidate integers: all guards here belong to a single proof modulus.

\subsection{A15: elliptic order fusion by primary projection}
\label{alg:A15}

\paragraph{Purpose and input.}
Several points on the same elliptic curve may each have a certified order
too small for the uniform Hasse threshold.  Their combined order
information can be much stronger.  Input $N>1$ with $\gcd(N,6)=1$, one curve
$E$ with discriminant coprime to $N$, and points $P_i$ whose fully factored
orders $f_i$ have been certified modulo every prime divisor of $N$.
Explicitly, the input checks are
\[
 [f_i]P_i=O,\qquad [f_i/q]P_i\text{ strongly nonzero for every }q\mid f_i.
\]
All point arithmetic must be valid at every such prime divisor.  Affine
implementations therefore check every denominator for invertibility and
reject unresolved exceptional cases.  These are the same essential
safeguards as in the Goldwasser--Kilian framework
\cite{gk,sutherland}.

\begin{proposition}[Fusion of certified elliptic orders]\label{prop:cert-ec-fusion}
Put $H=\lcm(f_1,\ldots,f_t)$.  For each $q\mid H$, select an index $i(q)$
with $v_q(f_{i(q)})=e_q=v_q(H)$ and form
\[
 Z_q=[f_{i(q)}/q^{e_q}]P_{i(q)},\qquad R=\sum_{q\mid H}Z_q.
\]
If the displayed arithmetic has been validated modulo every prime divisor
of $N$, then $R$ has order exactly $H$ over each such residue field.
Thus $H>(N^{1/4}+1)^2$ proves $N$ prime.  If the $f_i$ are pairwise
coprime, the simpler sum $R=P_1+\cdots+P_t$ has order $\prod_i f_i$.
\end{proposition}

\begin{proof}
The input conditions make the reduction of $P_i$ have exact order $f_i$.
Consequently $Z_q$ has exact order $q^{e_q}$.  In the abelian group
$E(\F_\ell)$, elements of pairwise coprime orders have a sum whose order
is their product: applying $H/q^{e_q}$ to an annihilating multiple isolates
the $q$-primary component.  Hence $R$ has order $H$.  Now
$H\mid\#E(\F_\ell)$ for every $\ell\mid N$, and Hasse's bound gives
$H\le(\sqrt\ell+1)^2$.  A prime divisor at most $\sqrt N$ would contradict
the stated threshold.  The coprime-order special case uses the same group
argument without primary projections.
\end{proof}

If all input certificates are retained, one may use
$H\mid\#E(\F_\ell)$ directly without forming $R$.  Constructing $R$ is
useful for emitting a compact, independently checkable single-point
certificate.  The final verifier can check $[H]R=O$ and the required
strongly nonzero reduced multiples without trusting the creator's point
construction.  Orders from different curves must not be merged by this
LCM rule.

\paragraph{Exact example and cancellation controls.}
For
\[
 N=101,\qquad E:y^2=x^3-4x-2,
\]
the points $P=(22,16)$ and $Q=(40,39)$ have certified orders $15$ and $7$,
respectively.  Their direct sum is $(69,53)$ and has order $105$.
Neither input order meets the uniform Hasse threshold, since
$(101^{1/4}+1)^2>16>15$.  The combined order does; the stronger integral
check $(\lfloor\sqrt{105}\rfloor-1)^4=6561>101$ also holds.
The generic primary-projection implementation produces the different
point $(18,1)$, again of exact order $105$.

Overlapping orders require the projections.  Using $T=(69,53)$, the points
$[7]T=(63,11)$ and $[3]T=(65,42)$ have orders $15$ and $35$, but their
unmodified sum is $[10]T$ and has order $21$, losing the factor $5$.
The projection algorithm instead returns $(75,91)$ of order $105$.
The verifier also checks the extreme cancellation $P+(-P)=O$ and confirms
that selecting one maximal primary component avoids this cancellation.
On the composite modulus $35$, adding the valid points $(0,1)$ and
$(5,14)$ on $y^2=x^3+1$ encounters denominator gcd $5$; the guard rejects
the operation and exposes the factor.

\paragraph{Cost and limits.}
Verification of the inputs uses their ordinary factored-order scalar
multiplications.  Construction adds one scalar multiplication per distinct
prime of $H$ and the corresponding point additions.  Compact output
verification depends on $H$ and its distinct prime factors rather than
on the entire construction history.  This can strengthen available
evidence from a largest individual order to their LCM, but it does not
make the input points or their orders free to find.  The examples establish
correct fusion, not an ECPP running-time improvement.

\subsection{A16: compiling residue constraints with exact Hasse windows}
\label{alg:A16}

\paragraph{Purpose and input.}
The uniform elliptic bound discards all information about possible small
prime divisors except their size.  A known point order also localizes a
small prime divisor to short intervals.  These intervals can be
intersected with residue information from checked $N-1$ or $N+1$
certificates.  This extends the earlier two-sided CRT closure by adding
an elliptic constraint.

Input $N>1$ with $\gcd(N,6)=1$, validated order witnesses implying
\[
 \ell\bmod M\in\mathcal R\quad\text{for every prime }\ell\mid N,
\]
and a validated elliptic point of fully factored order $H$ over every
prime divisor.  For example, an odd multiplicative part $F$ supplies
$M=2F$, $\mathcal R=\{1\}$.  A Lucas certificate supplies both possible
signs as in \cref{thm:plus}; distinct sign choices must be retained when
merging constraints by the generalized Chinese remainder theorem.
The order and residue premises are proof obligations, not user-supplied
assertions accepted without verification.

\begin{proposition}[Integer Hasse-window closure]\label{prop:cert-hasse-window}
Let
\[
 S=\lfloor\sqrt N\rfloor,\quad
 U=S+1+\isqrt(4S),\quad T=\lfloor U/H\rfloor,
\]
and for $1\le t\le T$ define
\[
 I_t=\bigl[tH+1-\isqrt(4tH),\ tH+1+\isqrt(4tH)\bigr]\cap[2,S].
\]
Every prime divisor of $N$ at most $S$ lies in
\[
 \mathcal D=\left\{d\in\Z:\ d\in\bigcup_{t=1}^{T}I_t,
                         \ d\bmod M\in\mathcal R\right\}.
\]
Under the verified input premises, $N$ is prime if and only if none of the
integers in $\mathcal D$ divides $N$.  In particular,
$\mathcal D=\varnothing$ is a primality certificate requiring no further
trial divisibility checks.
\end{proposition}

\begin{proof}
For $\ell\mid N$ with $\ell\le S$, the point-order condition and Lagrange's
theorem give $\#E(\F_\ell)=tH$ for some positive integer $t$.
Hasse's upper bound gives $tH\le U$, hence $t\le T$.  The two-sided Hasse
bound is equivalent to
\[
 (\ell+1-tH)^2\le4\ell
 \quad\Longleftrightarrow\quad
 (\ell-tH-1)^2\le4tH.
\]
Since all endpoints other than the square root are integers, this is
precisely membership in $I_t$.  The checked arithmetic witnesses give the
residue condition.  If $N$ is composite, it has a prime divisor in
$\mathcal D$; conversely, every divisor in $\mathcal D$ is between $2$ and
$\sqrt N$, hence is proper.  This proves the equivalence.
\end{proof}

\paragraph{Algorithm.}
Verify the point, its prime factors, all unit guards, and the arithmetic
witnesses.  Construct the residue classes rather than trusting an asserted
list.  Generate the integer intervals $I_t$, intersect them with each
residue progression, remove duplicates, and test the resulting divisors.
A discovered divisor returns \texttt{COMPOSITE}; complete exclusion returns
\texttt{PRIME}.  Exhausting a finite computation budget before completing
the set returns \texttt{INCONCLUSIVE}.  The bounded implementation executes
one multiplicative and one elliptic input certificate; the general
statement permits any independently verified source of the residue
restriction.

\begin{corollary}[A constant-size final closure regime]\label{cor:cert-one-window}
If $2H>U$ and $M>2\isqrt(4H)$, there is at most one integer to test per
residue class in $\mathcal R$.
\end{corollary}
\begin{proof}
The first inequality gives $T\le1$.  The untruncated interval $I_1$ has
endpoint difference $2\isqrt(4H)<M$, so each residue class modulo $M$
occurs at most once.
\end{proof}

This corollary concerns the final closure after the certificates are
provided.  It does not bound the cost of finding them.  If $H$ is comparable
to $\sqrt N$, a modulus of size comparable to $N^{1/4}$ can leave only
boundedly many final divisibility checks per residue class; it need not itself exceed
$\sqrt N$.  For smaller orders $H$, the number $T$ can be large and direct
window enumeration can be impractical.

\paragraph{A genuinely joint example.}
Take $N=211$.  The multiplicative witness $a=2$ for $F=5$ satisfies
\[
 2^{210}\equiv1\pmod{211},\quad 2^{42}\equiv107\pmod{211},
 \quad\gcd(106,211)=1.
\]
Thus every prime divisor is $1$ modulo $10$, but the bound $11^2>211$
fails.  On
\[
 E:y^2=x^3-2x+2,\qquad Q=(23,27),
\]
the verified affine point has order $19$.  This order also fails the
uniform elliptic threshold: with $S=14$, one has $U=22\ge19$.
The window computation, however, gives
\[
 T=1,\quad I_1=[20-8,20+8]\cap[2,14]=[12,14].
\]
No integer in $[12,14]$ is $1$ modulo $10$.  The combined certificate
therefore proves $211$ prime with zero final divisibility checks.  The
claim of jointness refers to these two uniform size tests; it is not a
claim that no other elementary proof of this small prime exists.
The verifier additionally compares the interval formula with direct
Hasse-interval divisibility for all orders $2\le H<120$ and hypothetical
divisors $2\le d\le100$, without floating-point arithmetic.

\subsection{A17: exact proof-network planning with shared ancestors}
\label{alg:A17}

\paragraph{Purpose and input.}
Caching a previously proved prime avoids repeated verification.  Choosing
which certificates to use is a separate problem when each integer has
several possible proofs and some proofs share ancestors.  Input a finite
catalogue of proof options.  An option $e$ has a conclusion $v(e)$, a set
of required already proved nodes $D(e)$, a concrete certificate, and a
nonnegative assigned cost $c(e)$.  Fix independently verified anchors
$A$ and a target node.  Costs describe a specified additive model; they
are not automatically predictions of wall time.

\begin{proposition}[Optimal verified-set planning]\label{prop:cert-plan}
Use subsets $V$ of the catalogue's nodes as states, starting with $A$.
For each option with $D(e)\subseteq V$ and $v(e)\notin V$, add the edge
\[
 V\longrightarrow V\cup\{v(e)\},\qquad\text{cost }c(e).
\]
A minimum-cost path to a state containing the target has minimum total
cost among valid derivations in the supplied catalogue, counting each
accepted conclusion once.  No unanchored cycle can be used to justify
itself.
\end{proposition}

\begin{proof}
A path gives an order of derivation in which every prerequisite is already
established.  Conversely, a valid derivation can be placed in a topological
order.  Repeatedly proving an established conclusion is unnecessary in the
nonnegative additive model, so deleting such repetitions leaves a path
of no greater cost.  The minimum path therefore solves the stated
optimization problem.  A cyclic collection with no available first step
does not create a transition from the anchored initial state.
\end{proof}

\paragraph{Algorithm and verification.}
Run Dijkstra's algorithm \cite{cert-dijkstra} on these subset states,
recording a predecessor option for each improved distance.  Reconstruct
the path, then verify each actual certificate in path order before
admitting its conclusion.  A failed certificate invalidates that plan
and must be removed or repaired before replanning.  The optimality claim
assumes a valid catalogue and its fixed cost model; the acceptance check
does not assume that an unverified catalogue is correct.

\paragraph{A shared-ancestor example.}
The anchors consist of $2$.  The catalogue contains valid Pocklington
proofs
\[
 2\longrightarrow3,\qquad3\longrightarrow7,\qquad
 3\longrightarrow13,\qquad\{7,13\}\longrightarrow2003.
\]
The first three use base $2$ and known parts $2,3,3$, respectively; the
last uses base $3$ and $F=7\cdot13=91\mid2002$.  Give these four options
costs $7,1,1,1$.  Also permit direct small-number trial proofs of $7$ and
$13$, each assigned cost $6$.

Optimizing the two branches independently chooses the direct proofs and
spends $6+6+1=13$ units.  The state-based planner instead proves $3$ once
and spends $7+1+1+1=10$ units.  These are deliberately assigned model
costs, not measured timings for these tiny primes.  All selected
certificates are verified, and exhaustive enumeration of all $64$
subsets of the six options confirms the optimum $10$.  A separate
unanchored two-node cycle is rejected.

\paragraph{Cost and limits.}
For $v$ nodes and $h$ options there are at most $2^v$ states and
$h2^v$ transitions.  A standard heap implementation takes
$O(h2^v\log(2^v))$ planning time apart from certificate arithmetic.
This is exact planning for a small portfolio, not a polynomial-time
solution for arbitrary large proof networks.  If real execution costs
depend on cache contents or memory pressure beyond the verified-node set,
the state and cost model must be enlarged before claiming optimality.

\subsection{A18: inheriting a prime certificate through an extracted divisor}
\label{alg:A18}

\paragraph{Purpose and input.}
A candidate that cannot be certified as a whole can still contain a new
certified prime.  The useful input is not merely a failed primality test:
it is a \emph{verified} statement that every prime divisor of an auxiliary
integer $N$ is at least $L$.  Such a statement can come from a valid
partial order certificate, even when $N$ is composite.  An exact divisor
$D$ must also be obtained, for example by a proper gcd exposed during
certificate construction.

\begin{proposition}[Hereditary divisor certification]\label{prop:cert-inherit}
Suppose every prime divisor of $N>1$ is at least $L\ge2$.  If
\[
 1<D\mid N,\qquad D<L^2,
\]
then $D$ is prime.  The lower-bound evidence for $N$ and the exact
divisibility and size checks suffice: no new modular exponentiation
with modulus $D$ is required.  If $L>P$ for a designated input prime $P$,
then the harvested prime is larger than $P$.
\end{proposition}

\begin{proof}
Every prime divisor of $D$ also divides $N$, hence is at least $L$.
If $D$ were composite, two of its prime factors counted with multiplicity
would give $D\ge L^2$.  This contradicts the strict size condition.
The same inheritance gives $D\ge L>P$.
\end{proof}

\paragraph{Algorithm.}
First verify the partial order or other lower-bound certificate on $N$.
Record the resulting $L$ as a proved predicate, not as an unchecked
parameter.  Extract or receive a proposed divisor and check it by exact
division.  Accept its primality only if it lies strictly between $1$
and $L^2$.  For a factorization into several extracted divisors, reuse the
same verified lower bound for every divisor that meets this condition.
Divisors outside the interval remain unresolved and require other
certificates.  The prime status of $N$ itself is never an input assumption.

\paragraph{Exact example.}
Start from the certified prime $P=3$ and the auxiliary lift
\[
 N=154P^2+1=1387.
\]
Base $2$ supplies the valid partial $F=9$ certificate
\[
 2^{1386}\equiv1\pmod{1387},\qquad
 2^{462}\equiv1322\pmod{1387},\qquad
 \gcd(1321,1387)=1.
\]
Thus every prime divisor is $1$ modulo $9$ and, by oddness, is at least
$19$.  The whole-modulus size bound fails since $1387>19^2$.
A further reduced power gives
\[
 2^{693}\equiv512\pmod{1387},\qquad
 \gcd(511,1387)=73,
\]
and exact division yields $1387=73\cdot19$.  Both divisors are below
$19^2=361$, so both $73$ and $19$ are proved prime by inheritance.  The
two child conclusions use zero additional modular powers after the
parent order evidence and factor extraction have been checked.

\paragraph{Cost and limits.}
The extra child verification consists of exact division and size
comparisons; the expensive order checks are shared.  Finding a useful
divisor can nevertheless be difficult.  Most ordinary composites rejected
by a Fermat test do not come with the required partial order certificate,
so this algorithm does not automatically recover primes from all failed
search candidates.  In particular, the previously logged Fermat-failing
large candidates cannot be retroactively accepted by this argument without
new evidence.  The restricted quadratic range of A01 already excludes
composite base-$2$ passes under its hypotheses; the present example uses
a larger coefficient outside that range.  The verifier rejects the
boundary divisor $D=L^2$ and proposed integers that do not divide $N$.

\paragraph{Executed evidence for A13--A18.}
The complete standard-library implementation is included in
\cref{app:certificates}.  Its successful status is
\texttt{ALL\_SIX\_CERTIFICATE\_ALGORITHMS\_VERIFIED}.
The tests cover exact examples, independent small-number factorizations,
592 product-gcd comparisons, overlapping elliptic orders, invalid
denominators, 118 integer Hasse-window comparisons, all $64$ proof-option
subsets in the planning example, and strict inheritance boundaries.
These bounded computations check the stated algorithms and their failure
semantics.  They are not a record-size search or a benchmark establishing
an asymptotic or universal speed advantage.

\section{Integrated experiments and independent checks}
\label{sec:revision-experiments}

\subsection{What was actually executed}

All eighteen modules have executed reference implementations. The
construction experiment starts by proving $P=2^{127}-1$ with the
Lucas--Lehmer recurrence, including a trial-division proof that $127$ is
prime. Its principal outputs are listed in
Table~\ref{tab:revision-constructions}. They are exact certified integers,
not probable-prime passes. No statement is made that these integers were
previously unknown.

\begin{table}[htbp]\centering\small
\caption{Executed construction outputs for the eighteen-module revision.}
\label{tab:revision-constructions}
\begin{tabularx}{\textwidth}{@{}lXr@{}}
\toprule
Module & Exact specification & Decimal digits\\
\midrule
A01 & $180P^2+1$ & 79\\
A02 & $82P^3+1$ & 117\\
A03 & $1104P-1$ and $1104P+1$ & 42 each\\
A04 & Prefix $314159$, suffix $27183$, middle width $90$;
exact progression parameter $j=195$ & 101\\
A05 & $601,1801\mid\Phi_{25}(2)$; $727\mid\Phi_{121}(2)$ & 3, 4, 3\\
A06 & $678P^2-1$ & 80\\
\bottomrule
\end{tabularx}
\end{table}

For A04, the integer and all its leading zeroes inside the middle block
are fixed by the formula and constants in the construction section and
Appendix~\ref{app:constructions}. A04 does not claim that any arbitrarily
chosen prefix and suffix admit a prime in every bounded middle interval.

The arithmetic and certificate modules have complementary demonstrations.
A07 certifies $7919$; A08 certifies $42379$ and $223999$; A11 certifies the
packet $727,1453,2179,2663$. A15 certifies $101$ by fused elliptic order
$105$. A16 certifies $211$ by intersecting two individually insufficient
uniform bounds. A18 certifies $19$ and $73$ as exact divisors of the
composite auxiliary integer $1387$. Other modules operate on these same
examples rather than each claiming a different prime. These overlapping
outputs are not added together as a discovery count.

\subsection{An alternate certificate for the existing large chain}
\label{sec:alternate-chain}

The quadratic members $p_1,\ldots,p_5$ of the continuing chain have
coefficients
\[
 438,\quad13696,\quad5494,\quad13258,\quad22330.
\]
Each coefficient is positive, even, and smaller than $4p_{i-1}-2$.
Their existing full base-$2$ Fermat congruences, together with
$2^{a_i}\not\equiv1\pmod{p_i}$, therefore give A01 certificates.
The revised seven-node verifier checks the short nonunity directly, in
addition to independently checking the original Pocklington gcd.

For the 18,882-digit endpoint
\[
 p_5=22330p_4^2+1,
\]
even that short modular power is unnecessary in a minimal A01 certificate:
the exact size comparison $1<2^{22330}<p_5$ proves nonunity.
The required long calculation is still
$2^{p_5-1}\bmod p_5=1$. Thus A01 provides a shorter
logical certificate for a previously certified large integer; it neither
discovers a new integer nor eliminates the large modular calculation.
If the original Pocklington gcd for the same base is already available,
it too implies short nonunity, since
$2^{a_i}\equiv1$ would imply $2^{a_ip_{i-1}}\equiv1$.

The stopped family $N=2mp_6^2+1$, $1\le m\le10^6$, also satisfies
the A01 coefficient bound. A future full Fermat pass together with short
nonunity would therefore be a proof in this family. Every durable completed
test in the interrupted run was a Fermat failure, so this observation
does not change that run's outcome.

\subsection{Independent oracles and deliberately invalid inputs}

The reference programs do more than reproduce their own formulas.
They compare different algorithms or use exact small-instance oracles.
The principal checks are recorded in Table~\ref{tab:revision-audit}.
An acceptance count below counts certificate attempts, possibly including
several witnesses for the same integer.

{\small
\begin{longtable}{@{}>{\raggedright\arraybackslash}p{0.25\textwidth}p{0.67\textwidth}@{}}
\caption{Finite verification evidence and its precise scope.}
\label{tab:revision-audit}\\
\toprule
Component & Executed check\\
\midrule\endfirsthead
\toprule Component & Executed check\\\midrule\endhead
\bottomrule\endfoot
Original inventory & The unchanged standalone listing certified all
32 original integers again.\\
Continuing chain & Seven exact nodes, through 18,885 digits, passed
Pocklington checks and decimal SHA-256 identity checks. A01 also
verified the five quadratic steps.\\
A01, A02, A06 & Construction verifier compared respectively
4,144, 672, and 1,981 witness attempts with trial division.\\
Lucas and cyclotomic arithmetic & 1,953 matrix-versus-recurrence
checks and 744 modular geometric-sum checks.\\
Independent boundary audit & Separately implemented checks:
8,272 positive and 517 negative attempts; 1,364 and 98 valid
certificates, respectively; zero false prime certificates.\\
A07, A08 & Digit join versus exhaustive words; quadratic root compiler
versus finite differences, including small characteristic.\\
A09, A14 & Exact proper-factor routing; 592 independent batches
compared with the gcd of each leaf.\\
A10 & 193,871 exhaustive canonical products, 768 fixed-seed stress
products through 4,096-bit radix, and 128 modular powers.\\
A11, A13 & Explicit four-number packet and rejected malformed packets;
86 small odd moduli audited by independent factorization.\\
A15, A16 & Invalid curves, wrong orders, nonunit denominators, and
point-cancellation controls; 118 exact Hasse-window comparisons.\\
A12, A17 & All 326 ordered screen subsets and all 64 certificate-edge
subsets in their respective finite examples.\\
A18 & Exact inherited divisors; nondivisors and equality at the
strict square boundary rejected.\\
\end{longtable}}

The mathematical proofs establish the general implications. Finite tests
check that the implementations match those implications on stated inputs;
they are not substitutes for proofs. Particularly useful controls are the
composite integers $91$, $35$, and $125$, which satisfy tempting
weakened versions of A01 or A06. The signed reduction counterexample
$N=5$, $x=24$ demonstrates why its product-domain restriction cannot be
silently enlarged.

The older elapsed times in the baseline sections remain historical
observations. The repeated baseline standard-library verification took
135.055 seconds in this session. The seven-node verification used GMP
through the optional Python interface; its exact observed duration is
reported in the execution ledger. These different workloads and arithmetic
backends are not a comparative benchmark. The small new prototypes complete
quickly on the present environment, but no record-scale speed claim is
extrapolated from them.

\section{How the modules can be combined}
\label{sec:revision-composition}

\subsection{Three concrete workflows}

\paragraph{Known-prime quadratic search.}
Choose a certified $p$ and a finite even-coefficient window.
Compile small-factor exclusions using A08, or apply A09 when many
explicit candidates are already available. A12 may choose a subset of
screen groups under an explicit cost model. For a surviving $ap^2+1$
inside A01's range, a full Fermat pass and short nonunity finish the
proof. A02 or the baseline order criteria are alternatives when the
chosen candidate family has different structure. The constructor need
not invoke every module.

\paragraph{Human-specified digit pattern.}
Use A04 for fixed end blocks or A07 for a more general digit language.
Both enforce an appropriate known factor of $N-1$ while constructing
candidates. Apply the small-factor screens, then an exact order
certificate. A13 can combine independently supplied primary witnesses;
A14 can verify their nonvanishing obligations through one successful
product gcd. The digit pattern shapes the output, while the order
certificate proves its primality.

\paragraph{Partial arithmetic and elliptic evidence.}
Verify a universal arithmetic residue constraint and one or more elliptic
point orders. A15 combines useful primary components into a larger known
order; A16 turns that order into exact Hasse windows and intersects
them with the arithmetic classes. If a proper factor is discovered during
the computation, A18 can certify a sufficiently small divisor using a
previously verified universal lower bound. A17 plans dependencies among
already available finite certificate proposals. It does not manufacture
the missing witnesses or predict their search cost.

\subsection{The meaningful sense of a new composition}

There are three distinct mathematical operations here. Candidate
construction imposes constraints on $N$ itself. Order fusion increases a
proved order in a finite group. Constraint intersection narrows the set
of possible prime divisors of $N$. Their combination is useful because
they act on different objects, but their types and hypotheses must be
preserved.

For example, A16's $N=211$ demonstration begins with
$\ell\equiv1\pmod{10}$ for every prime divisor and with an elliptic point
of exact order $19$. A hypothetical divisor
$\ell\le\lfloor\sqrt{211}\rfloor=14$ must lie in $[12,14]$.
This interval contains no integer congruent to $1$ modulo $10$, so the
intersection is empty. Neither uniform size bound alone suffices.
The new verification object is the combination of the two checked
constraints, not a guess that the candidate is unusually likely prime.

These procedures are implemented additions to this manuscript.
The coefficient inequalities, finite closures, and composition rules
are proved in full. Their historical novelty is not settled by a search
for matching terminology, and no priority claim is made. In particular,
classical Pocklington, Lucas, cyclotomic-order, CRT, dynamic-programming,
and elliptic principles remain credited to their established sources.

\section{Efficiency, digit blocks, and the Lucas--Lehmer comparison}
\label{sec:revision-efficiency}

Let $L=\bits(N)$, and let $\mathcal M(L)$ denote the cost of multiplying
$L$-bit integers. A full modular exponent whose exponent has $O(L)$ bits
requires $O(L)$ modular arithmetic steps by binary powering. With a
suitable reduction implementation, the corresponding cost is conventionally
expressed as $O(L\mathcal M(L))$, with gcd, setup, and certificate-factor
costs accounted for separately. A01 can remove a long reduced-power
obligation in its specific coefficient range; it does not reduce the
remaining full exponent to $O(\log L)$ steps.

For a Mersenne candidate $2^r-1$ with prime exponent $r$, Lucas--Lehmer
uses $r-2$ iterations of $s\leftarrow s^2-2$ modulo that candidate.
The modulus permits particularly efficient special reduction, and
implementations use fast multiplication \cite{revision-gimps-math}.
The present paper contains no matched benchmark establishing that its
general lifting certificates equal or beat that specialized test.
Having the same broad arithmetic order of growth does not establish
equal constants, memory use, or optimized throughput.

Large integers can indeed be processed in blocks. GMP represents them
using limbs, and its large multiplication algorithms split inputs,
transform or evaluate pieces, multiply, and recombine
\cite{revision-gmp}. A10 adds a proved reduction identity for a particular
modulus shape. These are ways to compute exact global arithmetic.
They do not reduce primality to independent primality decisions on decimal
chunks: $2323=23\cdot101$ has two prime two-digit chunks, whereas the prime
$1009$ has chunks $10$ and $09$, neither prime.

There is also a difference between reducing the number of arithmetic
operations and reducing total work. A11 changes four small-modulus
verifications into one larger-modulus packet. Its demonstration reduces
the count of squarings from $24$ to $6$, while the illustrative squared-bit
cost increases from $3054$ to $11094$. A12's $279$ to $156$ improvement
is in assigned model units, not measured seconds. A17 optimizes a finite
certificate proposal graph and is exponential in the number of nodes.
These limitations are part of the algorithms' stated contracts.

The practical next implementation step is therefore a matched benchmark
for a chosen candidate family and certificate route, including failed
candidates, sieve setup, arithmetic width, and independent verification.
This revision supplies exact prototypes and proofs for that work.
It does not claim that applying more filters, doubling the candidate's
digit length after each failure, or renaming classical operations yields
an automatic supply of giant primes.


\section{Reproducing the expanded manuscript}
\label{sec:revision-reproduce}

The complete deliverable consists of one PDF and this one TEX file.
Compilation needs ordinary LaTeX packages and no external figure,
bibliography database, numerical data file, or included source file.
Run \texttt{pdflatex Certified\_Prime\_Lifting.tex} at least twice,
or use \texttt{latexmk -pdf Certified\_Prime\_Lifting.tex}.

\subsection{Extracting the exact programs}

Appendix~\ref{app:code} retains the original verifier byte for byte.
Appendices~\ref{app:constructions}, \ref{app:arithmetic}, and
\ref{app:certificates} contain the three six-module verifiers;
Appendix~\ref{app:later-chain} verifies the continuing large chain;
Appendix~\ref{app:boundary-audit} supplies the independently implemented
boundary audit. Every listing has its own imports, constants, and entry
point. Extract them into separate directories because two programs write
an adjacent \texttt{results.json}. Save and run the following extractor
beside this TEX file:

\begin{lstlisting}
from pathlib import Path
import re

source = Path("Certified_Prime_Lifting.tex").read_text(encoding="utf-8")
items = [
    ('verifier', 'baseline', 'verify_baseline.py'),
    ('constructions', 'constructions', 'verify_constructions.py'),
    ('arithmetic', 'arithmetic', 'verify_arithmetic.py'),
    ('certificates', 'certificates', 'verify_certificates.py'),
    ('later-chain', 'later_chain', 'verify_later_chain.py'),
    ('boundary-audit', 'boundary_audit', 'independent_boundary_audit.py'),
]
for label, folder, filename in items:
    pattern = (r"\\begin\{lstlisting\}\[label=lst:" + re.escape(label)
               + r"\]\n(.*?)\\end\{lstlisting\}")
    match = re.search(pattern, source, re.S)
    if match is None:
        raise ValueError("Missing listing: " + label)
    target = Path("verifiers") / folder / filename
    target.parent.mkdir(parents=True, exist_ok=True)
    target.write_text(match.group(1), encoding="utf-8")
    print(target)
\end{lstlisting}

The three new module verifiers and the boundary audit use only Python's
standard library. Run, for example,
\begin{lstlisting}[language=bash]
python3 verifiers/constructions/verify_constructions.py
python3 verifiers/arithmetic/verify_arithmetic.py
python3 verifiers/certificates/verify_certificates.py
python3 verifiers/boundary_audit/independent_boundary_audit.py
\end{lstlisting}
Their required final status strings are, respectively,
\begin{quote}\small\ttfamily
A01\_A06\_ALL\_EXACT\_CHECKS\_PASSED\\
ALL\_SIX\_ARITHMETIC\_ALGORITHMS\_VERIFIED\\
ALL\_SIX\_CERTIFICATE\_ALGORITHMS\_VERIFIED\\
INDEPENDENT\_BOUNDARIES\_VERIFIED
\end{quote}

For the larger inherited inventories, run
\begin{lstlisting}[language=bash]
python3 verifiers/baseline/verify_baseline.py
python3 verifiers/later_chain/verify_later_chain.py
\end{lstlisting}
The respective success states are
\begin{quote}\small\ttfamily
ALL\_32\_PRIMES\_CERTIFIED\\
ALL\_SEVEN\_PRIMES\_CERTIFIED
\end{quote}
The latter program defaults to standard
Python integers. If \texttt{gmpy2} is separately installed, the optional
\texttt{--gmp} flag accelerates the same exact checks. The mathematical
certificate does not depend on that optional backend.

\subsection{Execution ledger and source identities}

The reference environment used Python 3.12.14 on x86-64 Linux.
The unchanged baseline was rerun successfully in 135.055 seconds.
The seven-node chain, including the additional A01 checks, was rerun
successfully with the optional GMP backend in 38.024
seconds. These are observations from different workloads, not a timing
comparison between algorithms. All three final module listings and the
independent boundary listing were extracted from the integrated TEX and
executed successfully.

The SHA-256 values below identify the exact UTF-8 program text, including
its final newline. They detect accidental edits and support reproducing
this revision; they do not replace the mathematical verification.

\noindent\textbf{baseline}\par
{\ttfamily\scriptsize 2dbc720444cefb09c060df3f6c4fd16c359abfc2b9df5a65a28574d7de05a68a}\par\smallskip
\noindent\textbf{constructions}\par
{\ttfamily\scriptsize 13a39c5f4470f508f4d012c15b3041c61c1b04d9ddad6a4d9a5d820f75f9db7b}\par\smallskip
\noindent\textbf{arithmetic}\par
{\ttfamily\scriptsize 0e92435d4089f559e87aac9b7c16cfdd09b0210f71b7c41497db96650c4f2154}\par\smallskip
\noindent\textbf{certificates}\par
{\ttfamily\scriptsize ac90cdca390824689358c7a79bd264ac83b4d2e8042820f7acef995bd506f144}\par\smallskip
\noindent\textbf{later\_chain}\par
{\ttfamily\scriptsize 8ac7bdfe891112e210eb5b366ada19ec2aa6c42cc3d574e1ae10849508f59509}\par\smallskip
\noindent\textbf{boundary\_audit}\par
{\ttfamily\scriptsize 15f0a3e548dc45759610250742f9c8aa91314bfd623dd6d1d82d0fbee57d7435}\par\smallskip

\clearpage
\begin{thebibliography}{99}
\small\raggedright\setlength{\itemsep}{3pt}\setlength{\parskip}{0pt}
\addcontentsline{toc}{section}{References}

\bibitem{pocklington}
H. C. Pocklington, \emph{The determination of the prime or composite nature
of large numbers by Fermat's theorem}, Proceedings of the Cambridge
Philosophical Society \textbf{18} (1914), 29--30.

\bibitem{bls}
J. Brillhart, D. H. Lehmer, and J. L. Selfridge,
\emph{New primality criteria and factorizations of $2^m\pm1$}, Mathematics
of Computation \textbf{29} (1975), no. 130, 620--647.
\url{https://doi.org/10.1090/S0025-5718-1975-0384673-1}.

\bibitem{morrison}
M. A. Morrison, \emph{A note on primality testing using Lucas sequences},
Mathematics of Computation \textbf{29} (1975), no. 129, 181--182.

\bibitem{maurer}
U. M. Maurer, \emph{Fast generation of prime numbers and secure public-key
cryptographic parameters}, Journal of Cryptology \textbf{8} (1995), no. 3,
123--155. \url{https://crypto.ethz.ch/publications/Maurer95a.html}.

\bibitem{gk}
S. Goldwasser and J. Kilian, \emph{Primality testing using elliptic curves},
Journal of the ACM \textbf{46} (1999), no. 4, 450--472.
\url{https://doi.org/10.1145/320211.320213}.

\bibitem{sutherland}
A. V. Sutherland, \emph{18.783 Elliptic Curves, Lecture 11: Primality
proving}, Massachusetts Institute of Technology, 7 March 2022.
\url{https://math.mit.edu/classes/18.783/2022/LectureSlides11.pdf}.

\bibitem{broker}
R. Br\"oker and P. Stevenhagen, \emph{Constructing elliptic curves of prime
order}, in \emph{Computational Arithmetic Geometry}, Contemporary
Mathematics \textbf{463}, American Mathematical Society, 2008, 17--28.
\url{https://arxiv.org/abs/0712.2022}.

\bibitem{lehmertest}
C. Caldwell, \emph{Primality proving: $N-1$ tests; $N+1$ tests and the
Lucas--Lehmer test}, PrimePages exposition. Accessed 5 November 2026.
\url{https://t5k.org/prove/prove3_1.html};
\url{https://t5k.org/prove/prove3_2.html}.

\bibitem{gimps}
GIMPS, \emph{List of known Mersenne prime numbers}. Source snapshot consulted
5 November 2026. \url{https://www.mersenne.org/primes/}.

\bibitem{primegrid}
Y. Gallot et al., \emph{$b^{2^n}-b^{2^{n-1}}+1$}, PrimeGrid discussion
thread 6062, beginning 27 December 2014.
\url{https://www.primegrid.com/forum_thread.php?id=6062}.

\bibitem{prothrecord}
PrimePages, \emph{$10223\cdot2^{31172165}+1$}, database entry 122473,
proven status and decimal digit count. Accessed 5 November 2026.
\url{https://t5k.org/primes/page.php?id=122473}.

\bibitem{primelist}
PrimePages, \emph{The largest known primes}, list snapshot dated
5 November 2026, 14:37:39 UTC.
\url{https://t5k.org/primes/lists/all.txt}.

\bibitem{prst}
PRST project, \emph{Primality tests for numbers of special forms},
software repository. The bounded record-size experiment used the binary
labelled 14.0.1888. \url{https://github.com/patnashev/prst}.

\bibitem{maynard}
J. Maynard, \emph{Dense clusters of primes in subsets}, Compositio
Mathematica \textbf{152} (2016), no. 7, 1517--1554.
\url{https://arxiv.org/abs/1405.2593}.

\bibitem{corpus}
Keio, \emph{Recurrences, Absolute Differences, and Collatz Orbits:
Algebra, Prime Patterns, and Zero Distributions}, manuscript
\texttt{[072] unified\_collatz\_recurrences.tex}, in
\emph{Verifiable Mathematical Reasoning Corpus},
commit \texttt{5936298ed4bacb7fb1423b889f02278fa521ea16}.
Relevant sections: ``Pascal--Franel weights and a certified ten-prime
trajectory'' and ``Primality-certificate barriers in Mersenne inverse
orbits.''
\href{https://github.com/QEEDDD/Verifiable-Mathematical-Reasoning-Corpus/tree/5936298ed4bacb7fb1423b889f02278fa521ea16}{Pinned repository snapshot}.
This manuscript is cited for the explicit constructions audited here,
not as blanket validation of the repository.

\bibitem{saito-irr}
K. Saito, \emph{Mills' constant is irrational}, Mathematika \textbf{71}
(2025), no. 3, e70027.
\url{https://doi.org/10.1112/mtk.70027};
\url{https://arxiv.org/abs/2404.19461v2}.

\bibitem{saito-var}
K. Saito, \emph{Transcendency of variants of Mills' constant},
arXiv:2508.16068v3, 7 December 2025. The version is fixed because the
theorem statements changed across revisions.
\url{https://arxiv.org/abs/2508.16068v3}.


\bibitem{philemon2021}
Ariko Stephen Philemon, \emph{Primality of numbers of the form
$ap^k+1$}, arXiv:2005.02327v5, 10 April 2021.
\url{https://arxiv.org/abs/2005.02327v5}.
Theorem 4.4 supplies the earlier $a<p$ quadratic special case discussed
here; the coefficient barriers in this revision are proved independently.

\bibitem{arith-hs}
E.~Horowitz and S.~Sahni,
\emph{Computing partitions with applications to the knapsack problem},
Journal of the ACM \textbf{21} (1974), no.~2, 277--292.
\url{https://doi.org/10.1145/321812.321823}.
Author-hosted text:
\url{https://www.cise.ufl.edu/~sahni/papers/computingPartitions.pdf}.

\bibitem{arith-gm}
T.~Granlund and P.~L.~Montgomery,
\emph{Division by invariant integers using multiplication},
Proceedings of the ACM SIGPLAN 1994 Conference on Programming Language
Design and Implementation (1994).
\url{https://doi.org/10.1145/178243.178249}.
\url{https://gmplib.org/~tege/divcnst-pldi94.pdf}.

\bibitem{arith-flint-crt}
FLINT developers, \emph{Multimodular reduction and Chinese remaindering},
official documentation, accessed 6 November 2026.
\url{https://flintlib.org/doc/mpn_extras/multi_crt.html}.

\bibitem{arith-flint-sqrt}
FLINT developers, \emph{fq\_sqrt: square roots in finite fields},
official documentation, accessed 6 November 2026.
\url{https://flintlib.org/doc/fq/fq_sqrt.html}.


\bibitem{cert-batchgcd}
FactHacks, \emph{Batch gcd}, primary project documentation accompanying
the 29C3 talk \emph{FactHacks: RSA factorization in the real world}.
\url{https://facthacks.cr.yp.to/batchgcd.html}.
See also \emph{Batch trial division},
\url{https://facthacks.cr.yp.to/batchtrial.html}.

\bibitem{cert-dijkstra}
E. W. Dijkstra, \emph{A note on two problems in connexion with graphs},
Numerische Mathematik \textbf{1} (1959), 269--271.
\url{https://doi.org/10.1007/BF01386390}.


\bibitem{revision-gmp}
GNU MP developers, \emph{GNU MP Manual: Nomenclature and Types;
FFT Multiplication}, official documentation, accessed 6 November 2026.
\url{https://gmplib.org/manual/Nomenclature-and-Types};
\url{https://gmplib.org/manual/FFT-Multiplication}.

\bibitem{revision-gimps-math}
GIMPS, \emph{Mathematics and Research Strategy}, official project
documentation, last updated 16 October 2024; accessed 6 November 2026.
\url{https://www.mersenne.org/various/math.php}.

\end{thebibliography}

\clearpage
\appendix
\section{Complete standalone certificate verifier}
\label{app:code}

The following single program constructs and certifies all 32 integers.
Large literals are confined to the norm representations, elliptic coordinates,
and metadata hashes. Their inclusion makes the TEX source self-contained;
there are no downloads or external certificate files. For readability in
print, long program lines may wrap visually; extraction from the TEX source
preserves the original Python line breaks.

\begin{lstlisting}[label=lst:verifier]
"""Complete constants and symbolic constructions for the 32-prime corpus.
No N decimal expansions, external data files, or primality tests are used.
Q coordinates and norm representations are the only large literals.
Factor tuples are (prime-id-or-small-integer, exponent, witness base).
This file constructs integers; a separate verifier checks primality.
"""
from math import prod

MERSENNE_EXPONENTS = (521, 1279, 4423)
PROTH_SEED = (401713, 3300, 3)
CHAIN_COEFFICIENTS = (438, 13696, 5494, 13258)
COLLATZ_R = 98723
PROTH_OUTPUT = ("mersenne_to_proth", "M1279", 2158, 3)
# (name, factors, lower bound: F+1, or 2F+1 for odd F)
N_MINUS_ONE = (('p0', ((2, 256, 11), (7, 256, 2)), 'plus_one'),
 ('p1', (('p0', 2, 2),), 'plus_one'),
 ('p2', (('p1', 2, 2),), 'plus_one'),
 ('p3', (('p2', 2, 2),), 'plus_one'),
 ('p4', (('p3', 2, 2),), 'plus_one'),
 ('power4', (('M1279', 4, 2),), 'plus_one'),
 ('mixed', (('M4423', 1, 2), ('Proth1000', 1, 2)), 'plus_one'),
 ('near_square', (('M4423', 1, 2),), 'odd'),
 ('linear_2214_p3', (('p3', 1, 2),), 'plus_one'),
 ('Phi3_M521_r260', ((2, 2, 2), (5, 1, 3), (13, 1, 2), ('M521', 1, 2)), 'plus_one'),
 ('Phi3_M521_r78299', ((13, 1, 2), (19, 1, 2), (317, 1, 2), ('M521', 1, 2)), 'plus_one'),
 ('Phi7_M521_r78299',
  ((13, 1, 2), (19, 1, 2), (317, 1, 2), ('M521', 1, 2), ('Phi3_M521_r78299', 1, 2)),
  'plus_one'),
 ('Phi3_M1279_r4683', ((3, 1, 2), (7, 1, 5), (223, 1, 2), ('M1279', 1, 2)), 'plus_one'),
 ('Phi3_M1279_r48909',
  ((3, 1, 2), (7, 1, 2), (17, 1, 2), (137, 1, 2), ('M1279', 1, 2)),
  'plus_one'),
 ('Phi7_M1279_r48909',
  ((3, 1, 5),
   (7, 1, 2),
   (17, 1, 2),
   (137, 1, 2),
   ('M1279', 1, 2),
   ('Phi3_M1279_r48909', 1, 2)),
  'plus_one'))

N_MINUS_ONE += (("Collatz_L0", (("M1279",1,2),), "plus_one"),
                ("Collatz_L1", (("M1279",1,2),), "plus_one"))

# (name, factors, delta, d): N=F^3-delta*F^2-d*F+1
CUBIC = (('cube_m1279', (('M1279', 1, 2),), 0, 1971),
 ('cube_m4423', (('M4423', 1, 2),), 0, 15751),
 ('cube_proth', (('Proth1000', 1, 2),), 2, 1315),
 ('cube_product', (('M1279', 1, 2), ('Proth1000', 1, 2)), 0, 6529))

# (name, seed, multiplicity, Lucas A, Lucas Q, criterion)
LUCAS = (("lucas_m1279", "M1279", 2, 3, 1, "rank"),
         ("negative_cube_Proth1000", "Proth1000", 1, 9, 1, "negative_cube"))
# (name, F-id, G-id, j offset, Fermat base, Lucas A, Lucas Q)
CRT = ("TwoSidedCRT", "M1279", "Proth1000", 963, 2, 5, 1)
# (Mersenne exponent, multiplier, whether Phi_7 is also certified)
CYCLOTOMIC = ((521,260,False),(521,78299,True),
              (1279,4683,False),(1279,48909,True))

CM_C_A = int(
    "310631762588305878343787163773812865717521638792091730228877763789034495"
    "954192907319663768591665987731375923749696423700820569088044808879886822"
    "258964235435536062255948339117868667989870666141079852134124567709538406"
    "111829961059967601471937794099279664347282513320846348431308128221727386"
    "304456250215307011466999782784002159857277076988608387962314279762358437"
    "670945365766441423864654270209319245536986659580841360980358909041299475"
    "77993450358413679645451258349128766354907608278402272633703964049120"
)
CM_C_B = int(
    "646394448418340813587086397455518241093626160973699044592409376946072812"
    "927542095303134986263719079210603866810486549745106402316447944647298652"
    "273501513907108885381875336005647835807391259856661179973876073735053139"
    "548478658554567266954915667679810463290320134325911326318563117849384230"
    "803896882501919744227857224518960422919908511746619077843495351690238831"
    "032042640196794609549495307918566284724809952407098795998911706752193668"
    "3319414552705087525196779080910885219957232976518677041086882554367"
)
CM_M_A = int(
    "125556301470467014213504119369214331566866510773858216486523063427163211"
    "730579729456184100200351588736361996066221109587475994042391413331307591"
    "4889861556182731228830864030057819504902195660902"
)
CM_M_B = int(
    "171575966758849216072637266161917841026461861951504617224221626925195888"
    "939795367706262827187804862499320222886436892319532435202804868666865774"
    "8637802679564998927566867053397023169176231059619"
)
CM1006_QX = int(
    "673328134373210046801457714014713279690189557668434309245922982534583515"
    "518159955843849693092589324214765684678235330557550525951029113487797472"
    "196288128368778431073692784112749595006186913490078007588504681706912446"
    "418675148430876068556139380879014814699613260676910273795841331548120786"
    "714646884099019442644671804074540609948429838015740516018093319758648102"
    "487808272682646570058786465803905579113640157932982668119697766822014200"
    "787584782907388197983483305308865883094283217040531813846955673280461999"
    "193689761980489930703563809886838460170213392325252669699327746831297726"
    "525858239266521888258778815306775417997899734778913812868273864954841532"
    "715556439956437201214505846689194679547303918549642110545341905703372353"
    "494214164696999257898989140686838100928179988046540138545993346187657388"
    "150002305622123517073683259161470309562935686280479929943675553884761917"
    "550888736676924031296152313975089887877975306914337333514367707611561032"
    "6812799201259876751021104657679903827666429159825211599439314951999606"
)
CM1006_QY = int(
    "295791267571015429399469841243025383457072761688802342862952596919286469"
    "638058812827729877323766676884829956394387139145544258435255240879250721"
    "561464817943649508278360054484248576537898356358549201494701916930496028"
    "494849275294493178497330039739734043145259689049222217345889656286031738"
    "589274684125611613988099130701869179880900506558570203627692684361803010"
    "644485615962123336772323563635516818534383619904642309448548278185411614"
    "741003751218094683567875455994619298970591757245688914985223908555365385"
    "191002341200848191744848529850397073202431827243246170146490498303281303"
    "302191106984361571820378666192831359056899116184768546128375004664325306"
    "968536635670478982981186658495835983465316179159592959814113806340668351"
    "938351911982581859637616142991972505184218472236304654236874288510854233"
    "465376310771296059472653935104672004428066494745951036411407358257415199"
    "362625082227105501881474747111492858205509802319100847367937472511625204"
    "2199945654442977439992970221081806596866208549288285688589352935870132"
)

CM1993_QX = int(
    "672120230134223320956150163200363857064374771257409266849368813848297679"
    "034741202886841462132176236840867912598995394451725726996937213675418436"
    "570148051371769361857073007986498823694298158084469098721463990653817048"
    "733634626074342596924250892494078144297296753095978210007913212230665841"
    "357919237578285944025243864996731956910473047000710999509736544045105848"
    "756677133390378390825082359935691094230286336613531916938185170467344044"
    "801745940518857225083938406178966907437461486694552583611601083862174362"
    "139160726089288571738099772783043032800334941466983644687852340975025700"
    "905852511815436687924445748219185339376674037493793778650645930596484894"
    "404510948356944666924606941590703553276859837799050782132772527636402356"
    "157683727854792153173940527810658572724321401111588543812199484504463181"
    "712088496729161295293930814169289853167090564260899935963893644172260214"
    "835036423639354797056904973861178562423403311837156174767980692022334415"
    "414303491262038101302309105359828001720751845490023899052527053732865151"
    "814694111026556229095157553286257537911161415417805867389728627807409628"
    "097485480360550300371826182388923232906346093903868912366970549590943401"
    "470747871492233849222103551801725991999595497158831383587164120349076704"
    "582316132266766671349944628268470822748973447084665845358708391843680815"
    "819332915321421665830885714987404101678098675035265125241922361493551995"
    "122501608626219270802543472826524075986326815509946677650049604188214391"
    "832136816717635126007925819402677254091585743964638036736438918062959742"
    "602008351929284253096774317119142701066557140015939256313082562451702634"
    "879602343682427895384256335154570882175090281908193107436044069398912646"
    "692654511366746951099209336761146954445225333104805110810181958461698388"
    "927245660700793996686782939926988870390728888296291653287460016411070057"
    "019372660810662943672627986183016837345045502125554655735721390677414009"
    "395504707977919722583539877382571807143185553703977075306328899024587738"
    "776616053576990241402045484717928966151620115604"
)
CM1993_QY = int(
    "210977068548549840502732220376689921800604169476658382330752966310619249"
    "902936487811763153197629792246655747353157308557907333368240594629861103"
    "449606438910892853930218434389328617934026170080446264339925472953371113"
    "464566475773676054732476648941204111295186860446023299701974839579684307"
    "614421984035401627194991690326491482019105583175928714834352247882370423"
    "634675221681688282392889314648872247213207541183346435774398839628735381"
    "859477362087896069646138380972503245085835456844472879621095068168692325"
    "201886690652150927813363865926553401444861111469799977336093013078354038"
    "646665383414117098330446517519634298904704064958659710077815187309364140"
    "283471406784217485032534531866749539621242106273492339593305912257144688"
    "622347527280750445377932771078782441436949653303414089884213331790558927"
    "360669277003510076013123230064979551094670858006637204556894292602984143"
    "141083587053999172880490445949722508953495162060756296408311963952806759"
    "914372927728396699557894195366870335929873878132682639916469266473563782"
    "978987926692680773107776011031115482294604425103498095130011554053336456"
    "493159182425591870231805615305394285341004905111701525526769262643708370"
    "078535257339959840665468498249817191654327533404324002279714917573741764"
    "695013205232125482951474646551864312855335682162464862403841433154445651"
    "490663000564102355354335132200854200277086685386375355000441401280993931"
    "434441384009004196728699935593190611529814542399636229955760374676294062"
    "353964090750340068409796987563264746658638292612430114008140695756107100"
    "273989083944232629052894723102845197277656773864179801062270880288942946"
    "487297570561412650441577516638029411298356975428240208276454613468056599"
    "888199094269709900174866528547896313576052834308534834438918035328007086"
    "566429193895474588242449248209564626898127810163416150957086219006491231"
    "977990978960584573546327635774523996982833685086555247739475491347090708"
    "986839995424905219449964749465292794892534035457140619072480036938397240"
    "995236499644750265400078634240992957074608326605"
)

CM770_QX = int(
    "104583438413007648105217157475695888305136794693377697778364328360156120"
    "004472282379727644779850450538519882159502164263590860200533527366051482"
    "396743927716836905198158267777263915376531965643426686391626093110781596"
    "755382462728242127271443576860784856868096229101972140064476408657445972"
    "548966323723879940729069233201295747514812066609515941488607043515729866"
    "004176940178386365801595256785628315143733157217218997389046590152266130"
    "485050260825562798124104300553146134386273648204301317246571909526588093"
    "525042476014413823036403933090484749076792500196962339460983132433191291"
    "689605199273217039797115850486493056925249981387058831728896130123429961"
    "105918717599525692527442473655990718636126842918288696592581928823295316"
    "64536405134119972073305183809149054641013462744195"
)
CM770_QY = int(
    "127390853015957650159978698277001669324476288351016561489290989753008104"
    "396318603573060361976158117283447985641319567797889625170203951506101567"
    "717947346500186197104037585232736204478626662966472793435959597313178299"
    "921307459603932400804512056444181142693422099417170580107581457006595032"
    "518827991242133385491320067654898012644616576071290414305940270342972960"
    "668395814884072604592097776802525673110485899608946159473201987235922391"
    "049744363851218275823781144186613974961870667905515073349012882455950818"
    "912900425405766399269492867608281592643526966616855545327830655329351910"
    "320532178060892699178696326903636738973069516507992576629415369717959948"
    "253589533946113815658886880691231033220844551634206602015034725431478536"
    "05782308494379427759187594771855977973191210724248"
)

# (name, seed-id, norm D, a, b, t, curve A, curve B, affine point)
CM = (
 ("elliptic_reverse_cm", "Proth1000", 1, CM_C_A, CM_C_B,
  3018, -4, 0, (CM1006_QX, CM1006_QY)),
 ("elliptic_reverse_cm_1650", "Proth1000", 1, CM_C_A, CM_C_B,
  2**1650+94, -4, 0, (CM1993_QX, CM1993_QY)),
 ("elliptic_reverse_cm_minus3_638", "M1279", 3, CM_M_A, CM_M_B,
  2**638+794, 0, 1, (CM770_QX, CM770_QY)),
)

def factor_value(factors, numbers):
    return prod((numbers[q] if isinstance(q,str) else q)**e
                for q,e,_base in factors)

def build_numbers():
    n = {"M%d"%e: 2**e-1 for e in MERSENNE_EXPONENTS}
    n["Proth1000"] = PROTH_SEED[0]*2**PROTH_SEED[1]+1
    n["p0"] = 14**512-14**256+1
    for i,k in enumerate(CHAIN_COEFFICIENTS,1):
        n["p%d"%i] = k*n["p%d"%(i-1)]**2+1
    for name,factors,delta,d in CUBIC:
        F = factor_value(factors,n)
        n[name] = F**3-delta*F**2-d*F+1
    n["power4"] = 2*1241*n["M1279"]**4+1
    n["mixed"] = 2*361*n["M4423"]*n["Proth1000"]+1
    n["near_square"] = 2*n["M4423"]*(n["M4423"]-868)+1
    n["linear_2214_p3"] = 2214*n["p3"]+1
    name,seed,e,_a = PROTH_OUTPUT
    n[name] = n[seed]*2**e+1
    n["lucas_m1279"] = 2*388*n["M1279"]**2-1
    F = n["Proth1000"]
    n["negative_cube_Proth1000"] = F**3-2*F**2-9053*F-1
    name,Fid,Gid,offset,_a,_A,_Q = CRT
    F,G = n[Fid],n[Gid]
    modulus = 2*F*G
    rho = 1+2*F*((-pow(F,-1,G))%G)
    n[name] = rho+(G//F+1+offset)*modulus
    for e,r,paired in CYCLOTOMIC:
        b = r*n["M%d"%e]
        n["Phi3_M%d_r%d"%(e,r)] = b*b+b+1
        if paired:
            n["Phi7_M%d_r%d"%(e,r)] = sum(b**j for j in range(7))
    for name,seed,D,a,b,t,A,B,Q in CM:
        assert a*a+D*b*b == n[seed]
        n[name] = n[seed]*t*t+2*a*t+1
    n["Collatz_L0"] = 6*COLLATZ_R*n["M1279"]+1
    n["Collatz_L1"] = 8*COLLATZ_R*n["M1279"]+1
    return n

# Digits and SHA-256 of ASCII decimal digits WITHOUT a newline.
EXPECTED = {'M1279': (386, 'aaa42d44bf59eb3b901c42a4f1b4404b4d483d07bcf4dfb920cb1518da372239'),
 'M4423': (1332, '5c6c10ad2e971f207b1a7b31e1bfc3b1f608fb99f01f8f3d7bcda27474174bd9'),
 'Proth1000': (1000, '830f39a37d5887cd387e6d7faca5c0b495434a11f35bc424f5e1d458adc2aa87'),
 'TwoSidedCRT': (1999, 'b5b20a0e1e1e9879b276ac949818757395becc9c93b4ce502f1aeae528e53dde'),
 'cube_m1279': (1156, '00cf7903f0a9859d9d30d59ff6736e6c32877864ec3395a868689f32e5822345'),
 'cube_m4423': (3995, 'a4d0e01af571dc1bff0af971e4333e2b69ca7b5291ba50e467795d8f8771d5fe'),
 'cube_product': (4153, '75633ee2ef9d19c089b5c56999ac741d16ee0bc6e53ff45f056613f72fbeb294'),
 'cube_proth': (2998, '368200f424894a48b345556771bb59680ac5d558483646739332e955b4976080'),
 'lucas_m1279': (773, 'a8c8e34a06e557738eda7d13eec4a0d921874032d0711f205fed8efed5647c67'),
 'mersenne_to_proth': (1035, 'a012ccd657c3e9fa00c45848ba312b9ce2c5eeed97ea6ccc3efc7de2e2a26164'),
 'mixed': (2334, 'bc694e73d89eb6473bfbfd80e5f2239cc38c6494a735b942929eab97c2be2856'),
 'near_square': (2664, '49f9a2da9fdd46a784988ce3274abd0706901f7d43f1b3598d7000d8755a9dc2'),
 'power4': (1544, '9ae48b1c2900b09afb20cd70383851f5a620dc44846f24cc727714920c918a78'),
 'M521': (157, '06514590f94c1d4595274279dfb1568077269f22e4258a5c9f928fe07e832e0e'),
 'p0': (587, 'a338e3a2cccacc978def8f28b08c248d39d5cd1542a47f8f0cf627d03febfb0c'),
 'p1': (1177, 'c852e5af4a1166124ffff6084fb1d986da411fe5b210aace6739009180e12c3f'),
 'p2': (2357, 'f1c08bab7b1e51ec1eef6dc053d1d98a2a13eb74e323c7273ac65c2dd2c58785'),
 'p3': (4718, 'c3e84997a4b2a949282840d14981c52bd98d66982c022a38528038779cd32566'),
 'p4': (9439, '82ddc0fab24f40e0c8c777c2f2e2138fec8d9bc505e21eb414858111ca910089'),
 'elliptic_reverse_cm': (1006, 'ad37b747ef780d2277a07dcc27f3c46306ed93468ba94405a4739a40a75ae16a'),
 'elliptic_reverse_cm_1650': (1993, '2c65661c212fed01711609fdb236ec578ce7393fc575d2798ca010c8b46c4c18'),
 'elliptic_reverse_cm_minus3_638': (770, '1dec1d9b2886ded668571100864e4cd4a01d226c9641f1dddfc1ee324bac6290'),
 'negative_cube_Proth1000': (2998, '003683f716c33e8388e84922b33ec5c417d2b0866eebc25da205ac8028cede92'),
 'linear_2214_p3': (4721, '0d509027ec97653dc3dda75e7cfc2de75f32365004d625dd5510beff7788e9bc'),
 'Phi3_M521_r78299': (324, '3d6ea87798e2ac072eab8dcf6eab4e2adcd88eccebf7b8db128fe2503bb4c52d'),
 'Phi7_M521_r78299': (971, '293ca4a619d7f565236f87e8f156411a5f8f757fcbef1eb3190b43804ab49850'),
 'Phi3_M521_r260': (319, '43f721ffb5e37bc3393bc3bad8841e51cfc90a4b444be7febb9fa2f6fff155d7'),
 'Phi3_M1279_r48909': (780, '9b30f8ed0708b9a32807f0113b0988832f5ac750bec62d19668abd793a61c356'),
 'Phi7_M1279_r48909': (2339, 'c168e7e14a3bbc74c1a83cf5a03c1d728e410d109b930d395754a07d43f3cac1'),
 'Phi3_M1279_r4683': (778, 'db4eb5d98b75f8fc9ed4d912be961ed5cf6c48396075269e3f1f89265eb147b0')}

EXPECTED.update({
 "Collatz_L0": (391, "720bb661d9d67e6f31922febbffe8a88410a12492a9f12d20c9f932435c565bb"),
 "Collatz_L1": (391, "3f2baf960fdd26c8277fbaf35e045906de1a38b86e80cafa2c80145383e0b3b4"),
})


# Complete standard-library verification. Run with Python 3.11 or later.
import hashlib
import json
import math
import sys
import time

sys.set_int_max_str_digits(0)


def require(ok, message):
    if not ok:
        raise ValueError(message)


def small_prime(n):
    return n >= 2 and all(n % d for d in range(2, math.isqrt(n)+1))


def mm(x, y, n):
    a,b,c,d = x
    e,f,g,h = y
    return ((a*e+b*g)%n, (a*f+b*h)%n,
            (c*e+d*g)%n, (c*f+d*h)%n)


def lucas_u(k, a, b, n):
    r, m = (1,0,0,1), (a%n,(-b)%n,1,0)
    while k:
        if k & 1:
            r = mm(r,m,n)
        k >>= 1
        if k:
            m = mm(m,m,n)
    return r[2]


def ec_mul(k, point, n, ca):
    def add(p, q):
        if p is None:
            return q
        if q is None:
            return p
        x,y = p
        u,v = q
        if x == u:
            if (y+v)%n == 0:
                return None
            require(y == v, 'Exceptional affine case')
            top, bot = 3*x*x+ca, 2*y
        else:
            top, bot = v-y, u-x
        require(math.gcd(bot,n) == 1, 'Nonunit denominator')
        slope = top*pow(bot,-1,n)%n
        z = (slope*slope-x-u)%n
        return z, (slope*(x-z)-y)%n
    r = None
    for bit in bin(k)[2:]:
        r = add(r,r)
        if bit == '1':
            r = add(r,point)
    return r


def verify_all():
    start = time.perf_counter()
    numbers = build_numbers()
    require(set(numbers) == set(EXPECTED), 'Inventory mismatch')
    require(len(numbers) == len(set(numbers.values())) == 32,
            'Expected exactly 32 distinct integers')
    for name,n in numbers.items():
        s = str(n)
        require((len(s),hashlib.sha256(s.encode('ascii')).hexdigest())
                == EXPECTED[name], 'Digit or hash mismatch: '+name)
    proved, report = set(), []

    def accept(name, method, t0):
        require(name not in proved, 'Duplicate acceptance')
        proved.add(name)
        row = dict(id=name, digits=EXPECTED[name][0], method=method,
                   seconds=round(time.perf_counter()-t0,6))
        report.append(row)
        print(json.dumps(row), flush=True)

    def prime_value(ref):
        if isinstance(ref,int):
            require(ref < 1000000 and small_prime(ref), 'Small factor')
            return ref
        require(ref in proved, 'Unproved dependency: '+ref)
        return numbers[ref]

    def minus_orders(n, factors):
        require(n > 1 and n%2 == 1, 'Invalid N')
        f, qs, full = 1, set(), {}
        for ref,e,a in factors:
            q = prime_value(ref)
            require(q not in qs and e >= 1, 'Invalid factor list')
            qs.add(q)
            require((n-1)%q**e == 0, 'Factor not in N-1')
            f *= q**e
            z = pow(a,(n-1)//q,n)
            if a not in full:
                full[a] = pow(z,q,n)
            require(full[a] == 1 and math.gcd(z-1,n) == 1,
                    'N-1 order witness failed')
        return f

    def plus_orders(n, ref, e, a, b):
        p = prime_value(ref)
        f, d = p**e, a*a-4*b
        require(n > 1 and n%2 == 1 and (n+1)%f == 0, 'N+1 shape')
        require(math.gcd(n,2*b*d) == 1, 'Lucas degeneracy')
        require(lucas_u(n+1,a,b,n) == 0, 'Full Lucas residue')
        require(math.gcd(lucas_u((n+1)//p,a,b,n),n) == 1,
                'Reduced Lucas gcd')
        return f

    for e in MERSENNE_EXPONENTS:
        t0 = time.perf_counter()
        require(small_prime(e) and e%2, 'Mersenne exponent')
        name, s = 'M%d'%e, 4
        for _ in range(e-2):
            s = (s*s-2)%numbers[name]
        require(s == 0, 'Lucas-Lehmer residue')
        accept(name,'Lucas-Lehmer',t0)

    t0 = time.perf_counter()
    k,e,a = PROTH_SEED
    n = numbers['Proth1000']
    require(k%2 == 1 and 0 < k < 2**e and n == k*2**e+1,
            'Proth seed shape')
    require(pow(a,(n-1)//2,n) == n-1, 'Proth seed witness')
    accept('Proth1000','Proth',t0)

    for name,factors,mode in N_MINUS_ONE:
        t0 = time.perf_counter()
        n = numbers[name]
        f = minus_orders(n,factors)
        require(mode in ('plus_one','odd'), 'Bound type')
        if mode == 'odd':
            require(f%2 == 1, 'Odd factor required')
        bound = 2*f+1 if mode == 'odd' else f+1
        require(bound*bound > n, 'Insufficient factored part')
        accept(name,'N-1 order and divisor bound',t0)

    for name,factors,delta,d in CUBIC:
        t0 = time.perf_counter()
        n = numbers[name]
        f = minus_orders(n,factors)
        require(f >= 3 and f%2 == 1 and 0 < d < f and d%2,
                'Cubic parameter range')
        require(delta == (0 if f%4 == 3 else 2), 'Cubic delta')
        require(n == f**3-delta*f*f-d*f+1, 'Positive cubic shape')
        accept(name,'Positive cubic theorem',t0)

    name,ref,e,a = PROTH_OUTPUT
    t0 = time.perf_counter()
    k, n = prime_value(ref), numbers[name]
    require(k%2 == 1 and k < 2**e and n == k*2**e+1,
            'Proth lift shape')
    require(pow(a,(n-1)//2,n) == n-1, 'Proth lift witness')
    accept(name,'Proth',t0)

    for name,ref,e,a,b,kind in LUCAS:
        t0 = time.perf_counter()
        n = numbers[name]
        f = plus_orders(n,ref,e,a,b)
        if kind == 'negative_cube':
            d, delta = 9053, 0 if f%4 == 3 else 2
            require(f%2 == 1 and 0 < d < f and d%2,
                    'Negative cubic parameters')
            require(n == f**3-delta*f*f-d*f-1 and n > 1,
                    'Negative cubic shape')
        else:
            require(kind == 'rank' and (f-1)**2 > n,
                    'Lucas size bound')
        accept(name,'Lucas '+kind,t0)

    t0 = time.perf_counter()
    name,fi,gi,offset,a,la,lb = CRT
    n, f, g = numbers[name], prime_value(fi), prime_value(gi)
    require(f%2 and g%2 and math.gcd(f,g) == 1, 'CRT inputs')
    require(minus_orders(n,((fi,1,a),)) == f, 'CRT minus')
    require(plus_orders(n,gi,1,la,lb) == g, 'CRT plus')
    modulus = 2*f*g
    rho = 1+2*f*((-pow(f,-1,g))%g)
    require(1 < rho < modulus and rho%(2*f) == 1
            and rho%g == g-1, 'CRT residue')
    require(n < modulus**2 and (rho*rho > n or n%rho != 0),
            'CRT divisor exclusion')
    accept(name,'Two-sided CRT',t0)

    for name,ref,d,u,v,t,ca,cb,q in CM:
        t0 = time.perf_counter()
        p, n = prime_value(ref), numbers[name]
        require(n > 3 and n%2 == 1 and t > 0, 'Elliptic input range')
        require(u > 0 and v > 0 and u%2 == 0 and v%2 == 1,
                'Norm parities')
        require(u*u+d*v*v == p and n == p*t*t+2*u*t+1,
                'Norm construction')
        require(math.gcd(n,6*(4*ca**3+27*cb**2)) == 1,
                'Bad curve reduction')
        require((math.isqrt(p)-1)**4 > n, 'Hasse bound')
        x,y = q
        require(0 <= x < n and 0 <= y < n
                and (y*y-x**3-ca*x-cb)%n == 0, 'Affine point')
        require(ec_mul(p,q,n,ca) is None, 'Elliptic order')
        accept(name,'Elliptic known prime order',t0)

    q0, q1 = numbers['Collatz_L0'], numbers['Collatz_L1']
    require(q1 > q0 > numbers['M1279'] and q0%2 == q1%2 == 1,
            'Collatz pair ordering')
    require(3*q1+1 == 4*q0, 'Exact Collatz three-step relation')
    require(proved == set(numbers), 'Uncertified integer remains')
    result = dict(status='ALL_32_PRIMES_CERTIFIED', count=len(proved),
                  total_seconds=round(time.perf_counter()-start,6),
                  python=sys.version, checks=report)
    print(json.dumps(result,indent=2), flush=True)
    return result


if __name__ == '__main__':
    verify_all()
\end{lstlisting}


\clearpage
\section{Construction algorithms A01--A06}
\label{app:constructions}

This complete program proves its common Mersenne seed, runs the bounded searches, and checks all construction certificates and controls.

\begin{lstlisting}[label=lst:constructions]
"""A01--A06: exact certificates and bounded checks, Python standard library.

All accepted large integers depend on an independently repeated Lucas--Lehmer
certificate for 2**127-1.  Fermat passes are never accepted without the stated
structural theorem.  No internet, probable-prime acceptance, or random input.
"""
from math import gcd, isqrt
from pathlib import Path
import hashlib
import json
import time


def need(ok, message):
    if not ok:
        raise ValueError(message)


def trial_prime(n):
    if n < 2:
        return False
    if n % 2 == 0:
        return n == 2
    return all(n % d for d in range(3, isqrt(n) + 1, 2))


def mersenne_seed(e):
    need(e > 2 and trial_prime(e), "prime Mersenne exponent")
    p = (1 << e) - 1
    s = 4
    for _ in range(e - 2):
        s = (s*s - 2) % p
    need(s == 0, "Lucas--Lehmer residue")
    return p


def mm(x, y, n):
    a,b,c,d = x
    e,f,g,h = y
    return ((a*e+b*g)%n, (a*f+b*h)%n,
            (c*e+d*g)%n, (c*f+d*h)%n)


def lucas_u(k, a, b, n):
    """Lower-left entry of [[a,-b],[1,0]]**k, modulo n."""
    need(k >= 0 and n > 1, "Lucas arguments")
    z = (1,0,0,1)
    x = (a % n, (-b) % n, 1, 0)
    while k:
        if k & 1:
            z = mm(z, x, n)
        x = mm(x, x, n)
        k >>= 1
    return z[2]


def jacobi(a, n):
    need(n > 0 and n % 2, "Jacobi modulus")
    a %= n
    s = 1
    while a:
        while a % 2 == 0:
            a //= 2
            if n % 8 in (3,5):
                s = -s
        a,n = n,a
        if a % 4 == n % 4 == 3:
            s = -s
        a %= n
    return s if n == 1 else 0


def metadata(n):
    s = str(n)
    return dict(n=s, digits=len(s), bits=n.bit_length(),
                sha256=hashlib.sha256(s.encode()).hexdigest())


def fermat_pair(n, a, b):
    need(1 < b < n, "witness range")
    f = pow(b, n-1, n)
    short = pow(b, a, n)
    return f, short


def a01(p, a, b=2):
    """p must already be certified prime; parity-strengthened plus closure."""
    need(p > 2 and p % 2 and a > 0 and a % 2 == 0,
         "A01 input parity")
    need(a < 4*p-2, "A01 strict coefficient bound")
    n = a*p*p+1
    f, short = fermat_pair(n, a, b)
    if f != 1:
        return dict(status="COMPOSITE", reason="Fermat failure")
    if short == 1:
        return dict(status="INCONCLUSIVE", reason="short power is one")
    return dict(status="CERTIFIED", a=a, base=b, full_residue=f,
                short_residue=str(short), **metadata(n))


def cubic_factors(p, a):
    n = a*p**3+1
    bound = (a*p+4)//(2*p+1)
    for t in range(2, bound+1, 2):
        d = (t*p)**2-4*(a*p-t)
        if d < 0:
            continue
        z = isqrt(d)
        if z*z != d or (t*p-z) % 2:
            continue
        u,v = (t*p-z)//2, (t*p+z)//2
        if u >= 2 and v >= 2 and u % 2 == v % 2 == 0:
            r,s = u*p+1, v*p+1
            need(r*s == n, "reconstructed cubic product")
            return dict(t=t, u=u, v=v, factor=str(r), cofactor=str(s))
    return None


def a02(p, a, b=2):
    """One Fermat pair plus a finite, coefficient-sized trace exclusion."""
    need(p > 2 and p % 2 and a > 0 and a % 2 == 0,
         "A02 input parity")
    n = a*p**3+1
    f, short = fermat_pair(n, a, b)
    if f != 1:
        return dict(status="COMPOSITE", reason="Fermat failure")
    if short == 1:
        return dict(status="INCONCLUSIVE", reason="short power is one")
    factors = cubic_factors(p,a)
    if factors is not None:
        return dict(status="COMPOSITE", reason="trace factorization", **factors)
    bound = (a*p+4)//(2*p+1)
    return dict(status="CERTIFIED", a=a, base=b, full_residue=f,
                short_residue=str(short), trace_upper_bound=bound,
                traces_checked=bound//2, **metadata(n))


def pocklington(n, p, exponent=1, base=2):
    f = p**exponent
    need(n > 1 and n % 2 and (n-1) % f == 0,
         "Pocklington factor identity")
    need((2*f+1)**2 > n, "Pocklington parity bound")
    residue = pow(base, n-1, n)
    g = gcd(pow(base, (n-1)//p, n)-1, n)
    return residue == 1 and g == 1


def lucas_known_factor(n, p, a, b=1, exponent=1):
    f=p**exponent
    need((n+1) % f == 0 and (2*f-1)**2 > n,
         "Lucas known factor and size")
    d = a*a-4*b
    if gcd(n,2*b*d) != 1:
        return False
    return (lucas_u(n+1,a,b,n) == 0 and
            gcd(lucas_u((n+1)//p,a,b,n),n) == 1)


def a06(p, a, u, v=1):
    """Short-index Lucas certificate for a*p*p-1; p is already certified."""
    need(p > 2 and p % 2 and a > 0 and a % 2 == 0,
         "A06 input parity")
    need(a < 4*p+2, "A06 strict coefficient bound")
    need(isqrt(a)**2 != a, "A06 excludes square coefficients")
    n = a*p*p-1
    d = u*u-4*v
    if gcd(n,2*v*d) != 1:
        return dict(status="INCONCLUSIVE", reason="degenerate parameters")
    short = lucas_u(a,u,v,n)
    short_gcd = gcd(short,n)
    if 1 < short_gcd < n:
        return dict(status="COMPOSITE", factor=str(short_gcd))
    if short_gcd != 1:
        return dict(status="INCONCLUSIVE", reason="short rank not excluded")
    full = lucas_u(n+1,u,v,n)
    if full:
        return dict(status="INCONCLUSIVE", reason="full Lucas condition fails")
    return dict(status="CERTIFIED", a=a, lucas_A=u, lucas_B=v,
                discriminant=d, full_residue=full,
                short_residue=str(short), short_gcd=short_gcd, **metadata(n))


SMALL = [q for q in range(3,101,2) if trial_prime(q)]


def sieve(n):
    return all(n == q or n % q for q in SMALL)


def search_a01(p, stop=10000):
    checked = 0
    for a in range(2,min(stop,4*p-2),2):
        n = a*p*p+1
        if not sieve(n):
            continue
        checked += 1
        r = a01(p,a)
        if r['status'] == 'CERTIFIED':
            r['survivors_checked'] = checked
            return r
    raise RuntimeError("bounded A01 search had no certificate")


def search_a02(p, stop=10000):
    checked = 0
    for a in range(2,stop,2):
        if not sieve(a*p**3+1):
            continue
        checked += 1
        r = a02(p,a)
        if r['status'] == 'CERTIFIED':
            r['survivors_checked'] = checked
            return r
    raise RuntimeError("bounded A02 search had no certificate")


def search_a03(p, stop=50000):
    checked = 0
    for j in range(1,stop+1):
        k = 6*j
        minus,plus = k*p-1,k*p+1
        need(plus < (2*p-1)**2, "A03 common size bound")
        if not sieve(minus) or not sieve(plus):
            continue
        checked += 1
        if pow(2,plus-1,plus) != 1 or pow(2,minus-1,minus) != 1:
            continue
        if not pocklington(plus,p):
            continue
        for u in range(3,64):
            if jacobi(u*u-4,minus) == -1 and lucas_known_factor(minus,p,u):
                return dict(status="BOTH_CERTIFIED", k=k,
                            plus=metadata(plus), minus=metadata(minus),
                            plus_base=2, lucas_A=u, lucas_B=1,
                            parameters_examined=j, surviving_pairs=checked,
                            minus_full=lucas_u(minus+1,u,1,minus),
                            minus_reduced_gcd=gcd(lucas_u(k,u,1,minus),minus),
                            plus_full=pow(2,plus-1,plus),
                            plus_reduced_gcd=gcd(pow(2,k,plus)-1,plus))
    raise RuntimeError("bounded A03 search had no certified pair")


def search_a04(p, prefix=314159, suffix=27183, suffix_digits=5,
               middle_digits=90, stop=10000):
    b = 10
    f = p*p
    shift = b**suffix_digits
    offset = prefix*b**(suffix_digits+middle_digits)+suffix
    need(gcd(shift,f) == 1, "A04 invertible suffix shift")
    need(0 <= suffix < shift and prefix > 0, "A04 block widths")
    r = ((1-offset)*pow(shift,-1,f)) % f
    max_x = min(b**middle_digits-1, (f*f-1-offset)//shift)
    need(r <= max_x, "A04 nonempty certified interval")
    checked = 0
    for j in range(stop):
        x = r+j*f
        if x > max_x:
            break
        n = offset+x*shift
        if n % 2 == 0 or not sieve(n):
            continue
        checked += 1
        if pocklington(n,p,2):
            need(str(n).startswith(str(prefix)), "prefix reconstruction")
            need(str(n).endswith(str(suffix).zfill(suffix_digits)),
                 "suffix reconstruction")
            return dict(status="CERTIFIED", prefix=prefix, suffix=suffix,
                        suffix_digits=suffix_digits, middle_digits=middle_digits,
                        middle=str(x), least_middle_residue=str(r), j=j,
                        known_factor=str(f), base=2, surviving_candidates=checked,
                        full_residue=pow(2,n-1,n),
                        reduced_gcd=gcd(pow(2,(n-1)//p,n)-1,n), **metadata(n))
    raise RuntimeError("bounded A04 search had no certificate")


def cyclotomic_prime_power(b,p,t):
    need(b >= 2 and p > 2 and p % 2 and t >= 1, "A05 input range")
    x = b**(p**(t-1))
    return sum(x**j for j in range(p))


def geometric_sum_mod(x,k,n):
    """Return (x**k, 1+...+x**(k-1)) modulo n without division."""
    power,total = 1,0
    block_power,block_total = x % n,1
    while k:
        if k & 1:
            total=(total+power*block_total)%n
            power=power*block_power % n
        block_total=block_total*(1+block_power) % n
        block_power=block_power*block_power % n
        k >>= 1
    return power,total


def a05(b,p,t,r):
    """p is a certified odd prime; r is a supplied exact divisor."""
    need(b >= 2 and p > 2 and p % 2 and t >= 1, "A05 input range")
    bound = 2*p**t+1
    need(1 < r, "A05 nontrivial candidate")
    need(gcd(r,p) == 1, "A05 exceptional-prime exclusion")
    need(r < bound*bound, "A05 strict square bound")
    x=pow(b,p**(t-1),r)
    residue=geometric_sum_mod(x,p,r)[1]
    need(residue == 0, "A05 exact cyclotomic divisibility")
    out=dict(status="CERTIFIED", b=b, p=p, t=t,
             lower_factor_bound=bound, cyclotomic_residue=residue,
             calculation="modular geometric sum without division", **metadata(r))
    if p**t * b.bit_length() <= 20000:
        c=cyclotomic_prime_power(b,p,t)
        need(c % r == 0, "A05 small expanded-value cross-check")
        out.update(cyclotomic_value=str(c),cofactor=str(c//r))
    return out


def mine_a05(b,p,t):
    need(trial_prime(p), "A05 seed prime")
    m = p**t
    bound = 2*m+1
    out=[]
    for r in range(2*m+1,bound*bound,2*m):
        x=pow(b,p**(t-1),r)
        if geometric_sum_mod(x,p,r)[1] == 0 and gcd(r,p) == 1:
            out.append(a05(b,p,t,r))
    return out


def search_a06(p, stop=10000):
    checked=0
    for a in range(2,min(stop,4*p+2),2):
        if isqrt(a)**2 == a:
            continue
        n=a*p*p-1
        if not sieve(n):
            continue
        checked+=1
        if pow(2,n-1,n) != 1:
            continue
        for u in range(3,64):
            if jacobi(u*u-4,n) != -1:
                continue
            r=a06(p,a,u)
            if r['status'] == 'CERTIFIED':
                r['survivors_checked']=checked
                return r
    raise RuntimeError("bounded A06 search had no certificate")


def rejected(call):
    try:
        call()
    except ValueError:
        return True
    return False


def controls():
    # Independent Lucas recurrence checks, including negative parameters.
    matrix_checks=0
    for n in (35,77,101):
        for aa in range(-2,5):
            for bb in (-1,1,2):
                x,y=0,1
                for k in range(31):
                    need(lucas_u(k,aa,bb,n) == x % n,
                         "matrix Lucas versus direct recurrence")
                    x,y=y,aa*y-bb*x
                    matrix_checks+=1
    geometric_checks=0
    for n in (35,77,101):
        for x in range(2,10):
            for k in range(31):
                pair=geometric_sum_mod(x,k,n)
                need(pair==(pow(x,k,n),sum(x**j for j in range(k))%n),
                     "geometric sum versus direct finite sum")
                geometric_checks+=1
    # Boundary and missing-hypothesis counterexamples, not just malformed data.
    need(91 == 7*13 and pow(9,90,91)==1 and pow(9,10,91)==9,
         "A01 boundary counterexample")
    need(rejected(lambda:a01(3,10,9)), "A01 must reject equality boundary")
    need(a01(3,6,54)['status']=='INCONCLUSIVE', "A01 short nonunity required")
    need(pow(191,216,217)==1 and pow(191,8,217)!=1,
         "A02 positive Fermat pair on a composite")
    bad=a02(3,8,191)
    need(bad['status']=='COMPOSITE' and bad['reason']=='trace factorization',
         "A02 trace exclusion must detect 217")
    need(int(bad['factor'])*int(bad['cofactor'])==217, "A02 exact factors")
    need(lucas_u(36,1,1,35)==0 and gcd(lucas_u(4,1,1,35),35)==1,
         "A06 square-coefficient counterexample")
    need(rejected(lambda:a06(3,4,1)), "A06 square guard")
    need(lucas_u(126,1,1,125)==0 and gcd(lucas_u(14,1,1,125),125)==1,
         "A06 boundary counterexample")
    need(rejected(lambda:a06(3,14,1)), "A06 strict bound guard")
    need(cyclotomic_prime_power(4,3,2)==4161 and 4161 % 57 == 0,
         "A05 exceptional-prime composite example")
    need(rejected(lambda:a05(4,3,2,57)), "A05 exceptional-prime guard")
    need(rejected(lambda:a05(2,5,2,1082401)), "A05 strict size guard")

    plus_cases=minus_cases=cubic_cases=0
    plus_accepts=minus_accepts=cubic_accepts=0
    for p in (3,5,7,11,13,17,19,23,29,31):
        need(trial_prime(p), "small seed")
        for a in range(2,4*p-2,2):
            n=a*p*p+1
            prime=trial_prime(n)
            for b in range(2,min(n,16)):
                r=a01(p,a,b)
                plus_cases+=1
                if r['status']=='CERTIFIED':
                    plus_accepts+=1
                    need(prime, "A01 versus trial division")
        for a in range(2,4*p+2,2):
            if isqrt(a)**2==a:
                continue
            n=a*p*p-1
            prime=trial_prime(n)
            for u in (1,3,4,5,6,7,8):
                r=a06(p,a,u)
                minus_cases+=1
                if r['status']=='CERTIFIED':
                    minus_accepts+=1
                    need(prime, "A06 versus trial division")
        if p <= 19:
            for a in range(2,26,2):
                n=a*p**3+1
                prime=trial_prime(n)
                for b in range(2,10):
                    r=a02(p,a,b)
                    cubic_cases+=1
                    if r['status']=='CERTIFIED':
                        cubic_accepts+=1
                        need(prime, "A02 versus trial division")
    return dict(status="ALL_CONTROLS_PASSED", matrix_recurrence_checks=matrix_checks,
                geometric_sum_checks=geometric_checks,
                A01_cases=plus_cases, A01_acceptances=plus_accepts,
                A02_cases=cubic_cases, A02_acceptances=cubic_accepts,
                A06_cases=minus_cases, A06_acceptances=minus_accepts,
                essential_guard_counterexamples=7)


def main():
    started=time.perf_counter()
    p=mersenne_seed(127)
    out=dict(seed=dict(formula="2^127-1",lucas_lehmer_residue=0,**metadata(p)))
    for name,fun in [('A01',search_a01),('A02',search_a02),
                     ('A03',search_a03),('A04',search_a04),
                     ('A06',search_a06)]:
        t=time.perf_counter()
        out[name]=fun(p)
        out[name]['observed_seconds']=time.perf_counter()-t
        print(name,out[name]['status'],flush=True)
    t=time.perf_counter()
    out['A05']=dict(status="CERTIFIED_DIVISORS",
                   cases=mine_a05(2,5,2)+mine_a05(2,11,2))
    out['A05']['observed_seconds']=time.perf_counter()-t
    need([int(x['n']) for x in out['A05']['cases'][:2]]==[601,1801],
         "A05 specified divisors")
    for x in out['A05']['cases']:
        need(trial_prime(int(x['n'])), "A05 independent small trial division")
    # Different classical certificate implications for the refined rules.
    need(pocklington(int(out['A01']['n']),p,2),
         "A01 also certified by classical Pocklington")
    need(pocklington(int(out['A02']['n']),p,3),
         "A02 also certified by classical Pocklington")
    need(lucas_known_factor(int(out['A06']['n']),p,
                           out['A06']['lucas_A'],out['A06']['lucas_B'],2),
         "A06 also certified by classical Lucas prime-power criterion")
    out['classical_cross_checks']={
        'A01':'Pocklington with known factor p^2',
        'A02':'Pocklington with known factor p^3',
        'A06':'Lucas with known factor p^2 and reduced index a*p'}
    out['controls']=controls()
    out['status']='A01_A06_ALL_EXACT_CHECKS_PASSED'
    out['observed_total_seconds']=time.perf_counter()-started
    path=Path(__file__).with_name('results.json')
    path.write_text(json.dumps(out,indent=2)+'\n')
    print(json.dumps(out,indent=2))


if __name__=='__main__':
    main()
\end{lstlisting}

\clearpage
\section{Arithmetic and search algorithms A07--A12}
\label{app:arithmetic}

This complete program includes the digit join, parameter sieve, two product trees, signed reduction, packet certificate, and finite planner.

\begin{lstlisting}[label=lst:arithmetic]
#!/usr/bin/env python3
"""Exact, bounded reference experiments for algorithms A07--A12.

Python standard library only. No probable-prime acceptance, network access,
unbounded prime search, or external data. The structural kernel is tested on
integer products, not asserted to make a composite number prime.
"""
from __future__ import annotations
from functools import lru_cache
from itertools import permutations, product
from math import gcd, isqrt, prod
from pathlib import Path
import hashlib
import json
import random
import time


def require(condition, message):
    if not condition:
        raise ValueError(message)


def trial_prime(n):
    if n < 2:
        return False
    return all(n % d for d in range(2, isqrt(n) + 1))


def small_primes(bound):
    return [n for n in range(2, bound + 1) if trial_prime(n)]


def certificate(n, prime_factors):
    """Find and then exactly verify an elementary Pocklington certificate."""
    factor = 1
    for q, e in prime_factors.items():
        require(trial_prime(q) and e >= 1, "uncertified seed factor")
        factor *= q ** e
    require(n > 2 and n % 2 and (n-1) % factor == 0,
            "invalid factored-part input")
    require(factor * factor > n, "known factor too small")
    witnesses = {}
    for q in prime_factors:
        for a in range(2, min(n, 100)):
            if (pow(a, n-1, n) == 1 and
                    gcd(pow(a, (n-1)//q, n)-1, n) == 1):
                witnesses[str(q)] = a
                break
        if str(q) not in witnesses:
            return None
    for qs, a in witnesses.items():
        q = int(qs)
        require(pow(a, n-1, n) == 1, "full residue failed")
        require(gcd(pow(a, (n-1)//q, n)-1, n) == 1,
                "reduced gcd failed")
    return {"n": n, "F": factor, "factorization": prime_factors,
            "witnesses": witnesses}


def digits_value(word, base):
    r = 0
    for d in word:
        r = base * r + d
    return r


def word_ok(word):
    return all(a != b for a, b in zip(word, word[1:]))


def digit_mitm(digit_sets, base, factor, wheel):
    """A07: sparse join; no table indexed by every residue modulo factor."""
    split = len(digit_sets)//2
    tail = len(digit_sets) - split
    left = [w for w in product(*digit_sets[:split]) if word_ok(w)]
    right = [w for w in product(*digit_sets[split:]) if word_ok(w)]
    buckets = {}
    for w in right:
        value = digits_value(w, base)
        buckets.setdefault(value % factor, []).append((w, value))
    candidates = []
    bucket_hits = 0
    for a in left:
        prefix = digits_value(a, base)
        target = (1 - prefix * pow(base, tail, factor)) % factor
        for b, suffix in buckets.get(target, []):
            bucket_hits += 1
            if a[-1] == b[0]:
                continue
            n = prefix * base**tail + suffix
            if max(wheel) < n < factor*factor:
                candidates.append(n)
    candidates.sort()
    survivors = [n for n in candidates if all(n % q for q in wheel)]
    return {"left_words": len(left), "right_words": len(right),
            "right_buckets": len(buckets), "bucket_hits": bucket_hits,
            "congruence_candidates": candidates, "survivors": survivors}


def experiment_a07():
    alphabet = (1, 3, 7, 9)
    sets = [alphabet] * 4
    f = 107
    wheel = (3, 5, 7, 11)
    out = digit_mitm(sets, 10, f, wheel)
    brute = [digits_value(w, 10) for w in product(alphabet, repeat=4)
             if word_ok(w)]
    congr = sorted(n for n in brute if n % f == 1 and
                   max(wheel) < n < f*f)
    require(out["congruence_candidates"] == congr,
            "MITM omitted or added a word")
    require(out["survivors"] == [n for n in congr
                                if all(n % q for q in wheel)],
            "wheel mismatch")
    certs = [certificate(n, {f: 1}) for n in out["survivors"]]
    require(all(c is not None for c in certs), "example not certified")
    require(all(trial_prime(c["n"]) for c in certs), "trial cross-check")
    out.update({"alphabet": alphabet, "length": 4,
                "language": "no equal adjacent digits", "F": f,
                "unrestricted_words": 4**4, "language_words": len(brute),
                "wheel_primes": wheel, "certificates": certs})
    return out


def polynomial_value(coefficients, t, modulus=None):
    value = 0
    for a in reversed(coefficients):
        value = value*t+a
        if modulus is not None:
            value %= modulus
    return value


def difference_heads(coefficients, modulus):
    row = [polynomial_value(coefficients, i, modulus)
           for i in range(len(coefficients))]
    heads = []
    while row:
        heads.append(row[0])
        row = [(b-a) % modulus for a, b in zip(row, row[1:])]
    return heads


def roots_by_differences(coefficients, q):
    """A08: updates must be performed from low to high index."""
    heads = difference_heads(coefficients, q)
    roots = []
    values = []
    for t in range(q):
        values.append(heads[0])
        if heads[0] == 0:
            roots.append(t)
        for j in range(len(heads)-1):
            heads[j] = (heads[j] + heads[j+1]) % q
    require(values == [polynomial_value(coefficients, t, q)
                       for t in range(q)], "finite-difference mismatch")
    return roots


def polynomial_sieve(coefficients, first, last, primes):
    factors = {}
    roots = {}
    for q in primes:
        residues = roots_by_differences(coefficients, q)
        roots[str(q)] = residues
        for r in residues:
            t = first + (r-first) % q
            for t in range(t, last+1, q):
                n = polynomial_value(coefficients, t)
                if n > q:
                    factors.setdefault(t, q)
    survivors = [t for t in range(first, last+1) if t not in factors]
    return roots, factors, survivors


def modular_sqrt(a, q):
    """Tonelli--Shanks on a prime field, with a checked returned root."""
    require(trial_prime(q), "root modulus not a proved small prime")
    a %= q
    if q == 2 or a == 0:
        return a
    if pow(a, (q-1)//2, q) != 1:
        return None
    if q % 4 == 3:
        r = pow(a, (q+1)//4, q)
    else:
        odd, v = q-1, 0
        while odd % 2 == 0:
            odd //= 2
            v += 1
        z = 2
        while pow(z, (q-1)//2, q) != q-1:
            z += 1
        c = pow(z, odd, q)
        r = pow(a, (odd+1)//2, q)
        t = pow(a, odd, q)
        m = v
        while t != 1:
            j, u = 0, t
            while u != 1 and j < m:
                u = u*u % q
                j += 1
            require(j < m, "Tonelli--Shanks invariant failed")
            b = pow(c, 1 << (m-j-1), q)
            r = r*b % q
            c = b*b % q
            t = t*c % q
            m = j
    require(r*r % q == a, "invalid square root")
    return r


def quadratic_lift_roots(factor, q):
    """Roots of factor*(2*t*t+2*t+2)+1 modulo q."""
    require(q > 1 and trial_prime(q), "invalid root-field prime")
    if q == 2 or factor % q == 0:
        return []
    delta = (-3-2*pow(factor,-1,q)) % q
    root = modular_sqrt(delta,q)
    if root is None:
        return []
    inv2 = pow(2,-1,q)
    return sorted({((-1+root)*inv2)%q,((-1-root)*inv2)%q})


def experiment_a08():
    f = 1009
    coefficients = (2*f+1, 2*f, 2*f)
    qs = small_primes(31)
    roots, factors, survivors = polynomial_sieve(coefficients, 0, 20, qs)
    compiled_roots = {str(q):quadratic_lift_roots(f,q) for q in qs}
    require(compiled_roots == roots, "quadratic root compiler mismatch")
    require(survivors == [t for t in range(21)
                          if all(polynomial_value(coefficients, t) % q
                                 for q in qs)], "sieve mismatch")
    for t, q in factors.items():
        n = polynomial_value(coefficients, t)
        require(trial_prime(q) and 1 < q < n and n % q == 0,
                "invalid rejection witness")
    certs = []
    for t in survivors:
        n = polynomial_value(coefficients, t)
        c = certificate(n, {f: 1})
        require((c is not None) == trial_prime(n), "bounded certificate check")
        if c is not None:
            c["t"] = t
            certs.append(c)
    # Degree exceeds some residue characteristics: no factorial inversion.
    extra_checks = 0
    for coefficients2 in [(3,-4,0,5,2), (0,1,0,-1), (7,)*7]:
        for q in (2,3,5,7,11):
            roots_by_differences(coefficients2, q)
            extra_checks += 1
    return {"F": f, "coefficients_low_to_high": coefficients,
            "t_interval": [0,20], "wheel_primes": qs, "roots": roots,
            "quadratic_root_compiler_matches": True,
            "rejected_parameters": {str(k):v for k,v in sorted(factors.items())},
            "surviving_parameters": survivors, "certificates": certs,
            "additional_polynomial_characteristic_checks": extra_checks,
            "finite_difference_additions": 2*sum(qs)}


def product_tree(values):
    require(values and all(n > 0 for n in values), "invalid product leaves")
    levels = [list(values)]
    while len(levels[-1]) > 1:
        a = levels[-1]
        levels.append([prod(a[j:j+2]) for j in range(0,len(a),2)])
    return levels


def remainder_tree(value, tree):
    values = [value % tree[-1][0]]
    for level in reversed(tree[:-1]):
        values = [values[j//2] % n for j,n in enumerate(level)]
    return values


def locate_factor(n, primes, qtree, level=None, index=0):
    """Descend only nodes whose product has nontrivial gcd with n."""
    if level is None:
        level = len(qtree)-1
    d = gcd(n, qtree[level][index])
    if d == 1:
        return None
    if level == 0:
        q = primes[index]
        require(n % q == 0, "bad factor route")
        return q
    for child in (2*index, 2*index+1):
        if child < len(qtree[level-1]):
            q = locate_factor(n, primes, qtree, level-1, child)
            if q is not None:
                return q
    raise ValueError("lost factor in product tree")


def batch_sieve(candidates, primes):
    require(all(trial_prime(q) for q in primes), "nonprime wheel entry")
    ntree = product_tree(candidates)
    qtree = product_tree(primes)
    residues = remainder_tree(qtree[-1][0], ntree)
    out = []
    for n, r in zip(candidates, residues):
        g = gcd(n, r)
        require(g == gcd(n, qtree[-1][0]), "batch gcd mismatch")
        if g == 1:
            out.append({"n":n,"gcd":g,"status":"UNRESOLVED"})
            continue
        q = locate_factor(n, primes, qtree)
        if q < n:
            require(n % q == 0, "improper rejection")
            out.append({"n":n,"gcd":g,"status":"COMPOSITE", "factor":q})
        else:
            require(q == n and trial_prime(q), "small prime equality")
            out.append({"n":n,"gcd":g,"status":"SMALL_PRIME", "factor":q})
    return out, len(ntree), len(qtree)


def experiment_a09():
    ns = [1713,7919,341,49,17,10403,42379,223999]
    qs = small_primes(43)
    out, nl, ql = batch_sieve(ns, qs)
    require([x["n"] for x in out if x["status"] == "COMPOSITE"] ==
            [1713,341,49], "unexpected factor routing")
    require([x["n"] for x in out if x["status"] == "SMALL_PRIME"] == [17],
            "small wheel prime misclassified")
    # A composite can survive every wheel prime.
    require(10403 == 101*103 and
            next(x for x in out if x["n"]==10403)["status"]=="UNRESOLVED",
            "survival incorrectly promoted")
    return {"candidates":ns,"wheel_primes":qs,"wheel_product":prod(qs),
            "candidate_tree_levels":nl,"wheel_tree_levels":ql,
            "leaf_gcds":out}


def cofactor_reduce(x, k, nbits, sign):
    """A10: exact one-correction reduction in the stated product range."""
    require(k >= 1 and nbits >= 1 and sign in (-1,1), "bad modulus shape")
    radix = 1 << nbits
    modulus = k*radix+sign
    require(modulus >= 3 and 0 <= x <= (modulus-1)**2, "outside proved range")
    high, low = x >> nbits, x & (radix-1)
    quotient, digit = divmod(high, k)
    value = low + (digit << nbits) - sign*quotient
    if sign == 1 and value < 0:
        value += modulus
    elif sign == -1 and value >= modulus:
        value -= modulus
    require(0 <= value < modulus, "one-correction range failed")
    return value


def cofactor_pow(a, exponent, k, nbits, sign):
    require(exponent >= 0, "negative exponent outside this algorithm")
    modulus = k*(1<<nbits)+sign
    a %= modulus
    r = 1
    for bit in bin(exponent)[2:]:
        r = cofactor_reduce(r*r,k,nbits,sign)
        if bit == '1':
            r = cofactor_reduce(r*a,k,nbits,sign)
    return r


def experiment_a10():
    exhaustive = 0
    for k in range(1,10):
        for nbits in range(1,5):
            for sign in (-1,1):
                modulus=k*(1<<nbits)+sign
                if modulus<3:
                    continue
                for u in range(modulus):
                    for v in range(modulus):
                        require(cofactor_reduce(u*v,k,nbits,sign)==u*v%modulus,
                                "exhaustive product failure")
                        exhaustive += 1
    rng=random.Random(20261106)
    stress=0
    for nbits in (64,256,1024,4096):
        for k in (1,3,17,2**31-1):
            for sign in (-1,1):
                modulus=k*(1<<nbits)+sign
                for _ in range(24):
                    u=rng.randrange(modulus);v=rng.randrange(modulus)
                    require(cofactor_reduce(u*v,k,nbits,sign)==u*v%modulus,
                            "large product failure")
                    stress+=1
                for exponent in (0,1,17,257):
                    require(cofactor_pow(3,exponent,k,nbits,sign)==
                            pow(3,exponent,modulus),"power failure")
    # Outside the range x <= (N-1)^2, a single correction can fail.
    k,nbits,sign,x=1,2,1,24
    modulus=5
    high,low=x>>nbits,x&3
    t,s=divmod(high,k)
    uncorrected=low+(s<<nbits)-t
    corrected_once=uncorrected+modulus if uncorrected<0 else uncorrected
    require(corrected_once==-1 and x%modulus==4,"boundary counterexample")
    rejected=False
    try:
        cofactor_reduce(x,k,nbits,sign)
    except ValueError:
        rejected=True
    require(rejected,"range guard missing")
    return {"exhaustive_canonical_product_checks":exhaustive,
            "stress_product_checks":stress,"max_radix_bits":4096,
            "power_checks":4*4*2*4,
            "outside_range_counterexample":{"N":5,"x":24,
                 "one_correction_result":corrected_once,"true_remainder":4},
            "runtime_speed_claim":"NONE"}


def crt(residues, moduli):
    require(len(residues)==len(moduli) and moduli,"empty CRT packet")
    z,modulus=0,1
    for residue,n in zip(residues,moduli):
        require(gcd(modulus,n)==1,"CRT inputs not pairwise coprime")
        step=((residue-z)*pow(modulus,-1,n))%n
        z+=modulus*step
        modulus*=n
    require(all(z%n==r%n for r,n in zip(residues,moduli)),"CRT mismatch")
    return z,modulus


def verify_packet(p,moduli,z):
    """A11 proof depends only on order witnesses, not local trial division."""
    if not trial_prime(p) or p<3 or not moduli:
        return False
    if any(n<=1 or n>=p**4 for n in moduli):
        return False
    modulus=prod(moduli)
    if not 0<=z<modulus:
        return False
    y=pow(z,p,modulus)
    return pow(y,p,modulus)==1 and gcd(y-1,modulus)==1


def binary_cost(exponent):
    require(exponent>=1,"positive exponent needed")
    return exponent.bit_length()-1, exponent.bit_count()-1


def experiment_a11():
    p=11
    multipliers=[6,12,18,22]
    moduli=[k*p*p+1 for k in multipliers]
    local=[pow(2,k,n) for k,n in zip(multipliers,moduli)]
    z,modulus=crt(local,moduli)
    require(verify_packet(p,moduli,z),"correct packet failed")
    require(all(trial_prime(n) for n in moduli),"independent small primality")
    y=pow(z,p,modulus)
    require(all(y%n==pow(r,p,n) for r,n in zip(local,moduli)),
            "packet residue extraction")
    bad=[243]+moduli[1:]
    # Keep the same candidate construction, including k=2 for 243.
    badlocal=[pow(2,k,n) for k,n in zip([2]+multipliers[1:],bad)]
    badz,badm=crt(badlocal,bad)
    require(not verify_packet(p,bad,badz),"composite packet accepted")
    require(not verify_packet(9,moduli,z),"composite seed accepted")
    require(not verify_packet(p,moduli,0),"zero witness accepted")
    require(not verify_packet(p,moduli,1),"identity witness accepted")
    require(not verify_packet(p,[p**4],0),"size boundary accepted")
    sq,mul=binary_cost(p)
    widths=[n.bit_length() for n in moduli]
    return {"p":p,"multipliers":multipliers,"moduli":moduli,
            "local_roots":local,"packet_modulus":modulus,"packet_witness":z,
            "intermediate_y":y,"final_power":pow(y,p,modulus),
            "final_gcd":gcd(y-1,modulus),"packet_status":"ALL_CERTIFIED",
            "component_bits":widths,"packet_bits":modulus.bit_length(),
            "individual_binary_operations":{"squares":2*sq*len(moduli),
                                            "multiplies":2*mul*len(moduli)},
            "packet_binary_operations":{"squares":2*sq,"multiplies":2*mul},
            "schoolbook_square_work_individual":2*sq*sum(b*b for b in widths),
            "schoolbook_square_work_packet":2*sq*modulus.bit_length()**2,
            "negative_controls":5,"runtime_speed_claim":"NONE"}


def optimal_screen_plan(rejections,costs,terminal_cost,count):
    """A12: finite population, deterministic outcomes, prescribed cost model."""
    all_ids=frozenset(range(count))
    s=len(rejections)
    @lru_cache(None)
    def remaining(mask):
        alive=all_ids
        for j,R in enumerate(rejections):
            if mask>>j&1:
                alive=alive-R
        return alive
    @lru_cache(None)
    def solve(mask):
        alive=remaining(mask)
        best=(terminal_cost*len(alive),())
        for j in range(s):
            if not mask>>j&1:
                tail,order=solve(mask|1<<j)
                offer=(costs[j]*len(alive)+tail,(j,)+order)
                if offer[0]<best[0]:
                    best=offer
        return best
    cost,order=solve(0)
    return cost,order,solve.cache_info().currsize


def plan_cost(order,rejections,costs,terminal_cost,count):
    alive=set(range(count))
    total=0
    for j in order:
        total+=costs[j]*len(alive)
        alive-=rejections[j]
    total+=terminal_cost*len(alive)
    return total,sorted(alive)


def experiment_a12():
    f=1009
    ns=[f*(2*t*t+2*t+2)+1 for t in range(21)]
    qs=[17,3,13,5,7]
    costs=[5,1,4,1,2]
    terminal=40
    rejections=[frozenset(t for t,n in enumerate(ns) if n%q==0) for q in qs]
    opt,order,states=optimal_screen_plan(rejections,costs,terminal,len(ns))
    best_exhaustive=None
    checked=0
    for length in range(len(qs)+1):
        for perm in permutations(range(len(qs)),length):
            cost,_=plan_cost(perm,rejections,costs,terminal,len(ns))
            best_exhaustive=cost if best_exhaustive is None else min(best_exhaustive,cost)
            checked+=1
    require(opt==best_exhaustive,"DP not globally optimal in model")
    baseline,rest=plan_cost(range(len(qs)),rejections,costs,terminal,len(ns))
    actual,rest_opt=plan_cost(order,rejections,costs,terminal,len(ns))
    require(actual==opt,"plan reconstruction cost")
    return {"population":"1009*(2*t*t+2*t+2)+1, t=0..20", "count":len(ns),
            "screen_primes":qs,"abstract_costs":costs,"terminal_cost":terminal,
            "rejection_sets":{str(q):sorted(R) for q,R in zip(qs,rejections)},
            "optimal_order":[qs[j] for j in order],"optimal_cost":opt,
            "baseline_order":qs,"baseline_cost":baseline,
            "terminal_parameters":rest_opt,"DP_states":states,
            "ordered_subsets_exhaustively_checked":checked,
            "planning_cost_in_objective":False,
            "wall_time_or_future_population_claim":"NONE"}


def main():
    start=time.perf_counter()
    results={"A07":experiment_a07(),"A08":experiment_a08(),
             "A09":experiment_a09(),"A10":experiment_a10(),
             "A11":experiment_a11(),"A12":experiment_a12()}
    results["status"]="ALL_SIX_ARITHMETIC_ALGORITHMS_VERIFIED"
    results["reference_runtime_seconds"]=time.perf_counter()-start
    results["source_sha256"]=hashlib.sha256(Path(__file__).read_bytes()).hexdigest()
    output=Path(__file__).with_name("results.json")
    output.write_text(json.dumps(results,indent=2,sort_keys=True)+"\n")
    print(json.dumps(results,indent=2,sort_keys=True))


if __name__=="__main__":
    main()
\end{lstlisting}

\clearpage
\section{Certificate algorithms A13--A18}
\label{app:certificates}

This complete program verifies all input evidence before applying the certificate transformations and hybrid closure.

\begin{lstlisting}[label=lst:certificates]
"""Exact, standard-library examples for Algorithms A13--A18.

These are bounded demonstrations of certificate transformations, not a
large-prime search.  Every acceptance uses explicit integer conditions.
No probable-prime result is accepted as a proof.
"""
from math import gcd, isqrt, prod
from heapq import heappop, heappush
import json
import time


def need(ok, message):
    if not ok:
        raise ValueError(message)


def trial_prime(n):
    return n >= 2 and all(n % d for d in range(2, isqrt(n) + 1))


def factor_small(n):
    need(n >= 1, "positive factorization input")
    out, q = {}, 2
    while q * q <= n:
        while n % q == 0:
            out[q] = out.get(q, 0) + 1
            n //= q
        q += 1
    if n > 1:
        out[n] = out.get(n, 0) + 1
    return out


def factored_value(factors):
    need(bool(factors), "nonempty factorization")
    for q, e in factors.items():
        need(trial_prime(q) and isinstance(e, int) and e >= 1,
             "certified prime factors and positive exponents")
    return prod(q ** e for q, e in factors.items())


def pocklington_partial(n, factors, witnesses):
    """Return a prime-divisor lower bound; n need not be prime."""
    need(n > 1 and n % 2 == 1, "odd modulus greater than one")
    f = factored_value(factors)
    need((n - 1) % f == 0, "known part divides n-1")
    residues = []
    for q in sorted(factors):
        a = witnesses[q]
        full = pow(a, n - 1, n)
        reduced = pow(a, (n - 1) // q, n)
        d = gcd(reduced - 1, n)
        need(full == 1 and d == 1, "Pocklington order obligation")
        residues.append([q, a, full, reduced, d])
    lower = 2 * f + 1 if f % 2 else f + 1
    return {"N": n, "F": f, "prime_divisor_lower_bound": lower,
            "residues": residues, "proves_whole_N": lower * lower > n}


def verify_factored_order(n, h, z):
    factors = factor_small(h)  # only the small demonstration inputs
    need(h > 1 and (n - 1) % h == 0, "factored exponent")
    need(pow(z, h, n) == 1, "annihilation by H")
    gs = [gcd(pow(z, h // q, n) - 1, n) for q in factors]
    need(all(d == 1 for d in gs), "exact H-order over every prime divisor")
    lower = 2 * h + 1 if h % 2 else h + 1
    return {"H": h, "z": z, "gcds": gs,
            "prime_divisor_lower_bound": lower,
            "proves_whole_N": lower * lower > n}


def prime_power_fusion(n, witnesses):
    """A13: separate full Pocklington witnesses -> one H-order witness."""
    e = n - 1
    need(n > 1 and n % 2 and witnesses, "fusion input")
    h, z, components = 1, 1, []
    for q, a in sorted(witnesses.items()):
        need(trial_prime(q) and e % q == 0, "selected certified factor")
        need(pow(a, e, n) == 1, "input full power")
        need(gcd(pow(a, e // q, n) - 1, n) == 1, "input q-order")
        t = 1
        while e % (t * q) == 0:
            t *= q
        u = pow(a, e // t, n)
        h *= t
        z = z * u % n
        components.append({"q": q, "prime_power": t, "a": a,
                           "projection": u})
    record = verify_factored_order(n, h, z)
    record.update({"N": n, "components": components})
    return record


def batch_units(n, values):
    """A14: one product gcd on success; exact leaf diagnosis on failure."""
    need(n > 1 and values, "nonempty unit batch")
    values = [x % n for x in values]

    def build(lo, hi):
        if hi - lo == 1:
            return (lo, hi, values[lo], None, None)
        mid = (lo + hi) // 2
        left, right = build(lo, mid), build(mid, hi)
        return (lo, hi, left[2] * right[2] % n, left, right)

    tree = build(0, len(values))
    bad, calls, root_gcd = [], 0, None

    def walk(node):
        nonlocal calls, root_gcd
        calls += 1
        d = gcd(node[2], n)
        if calls == 1:
            root_gcd = d
        if d == 1:
            return
        if node[1] - node[0] == 1:
            bad.append([node[0], d])
            return
        walk(node[3])
        walk(node[4])

    walk(tree)
    proper = sorted({d for _, d in bad if 1 < d < n})
    status = ("ALL_UNITS" if not bad else
              "PROPER_FACTOR_FOUND" if proper else
              "NONUNIT_WITHOUT_PROPER_FACTOR")
    return {"N": n, "values": values, "root_gcd": root_gcd,
            "failed_leaves": bad, "proper_factors": proper,
            "tree_gcd_calls": calls, "status": status}


class NonunitDenominator(ValueError):
    def __init__(self, factor):
        super().__init__("elliptic denominator is not a unit")
        self.factor = factor


def on_curve(p, n, a, b):
    return p is None or (p[1] * p[1] - p[0] ** 3 - a * p[0] - b) % n == 0


def ec_add(p, q, n, a, b):
    """Affine operations with explicit conditions valid at every prime divisor."""
    need(on_curve(p, n, a, b) and on_curve(q, n, a, b), "curve membership")
    if p is None:
        return q
    if q is None:
        return p
    x, y = p
    u, v = q
    if x == u:
        if (y + v) % n == 0:
            return None
        need(y == v, "unresolved equal-x exceptional case")
        numerator, denominator = 3 * x * x + a, 2 * y
    else:
        numerator, denominator = v - y, u - x
    d = gcd(denominator, n)
    if d != 1:
        raise NonunitDenominator(d)
    slope = numerator * pow(denominator, -1, n) % n
    rx = (slope * slope - x - u) % n
    r = (rx, (slope * (x - rx) - y) % n)
    need(on_curve(r, n, a, b), "addition result on curve")
    return r


def ec_mul(k, p, n, a, b):
    need(k >= 0, "nonnegative scalar")
    r = None
    for bit in bin(k)[2:]:
        r = ec_add(r, r, n, a, b)
        if bit == "1":
            r = ec_add(r, p, n, a, b)
    return r


def ec_order_record(n, a, b, p, factors):
    need(n > 1, "modulus greater than one")
    f = factored_value(factors)
    need(gcd(n, 6) == 1, "small characteristic excluded")
    need(gcd(4 * a ** 3 + 27 * b * b, n) == 1, "smooth reductions")
    need(p is not None and on_curve(p, n, a, b), "affine input point")
    need(ec_mul(f, p, n, a, b) is None, "elliptic annihilation")
    reduced = []
    for q in sorted(factors):
        t = ec_mul(f // q, p, n, a, b)
        need(t is not None, "strongly nonzero reduced multiple")
        reduced.append([q, list(t)])
    return {"order": f, "point": list(p), "reduced_multiples": reduced}


def elliptic_fusion(n, a, b, records):
    """A15: project overlapping certified orders to primary components."""
    checked, best = [], {}
    for i, (p, factors) in enumerate(records):
        checked.append(ec_order_record(n, a, b, p, factors))
        for q, e in factors.items():
            if q not in best or e > best[q][0]:
                best[q] = (e, i)
    r, h, projections = None, 1, []
    for q, (e, i) in sorted(best.items()):
        t = q ** e
        f = checked[i]["order"]
        z = ec_mul(f // t, records[i][0], n, a, b)
        need(z is not None, "nonzero primary component")
        r = ec_add(r, z, n, a, b)
        h *= t
        projections.append([q, e, i, list(z)])
    compact = ec_order_record(n, a, b, r, {q: e for q, (e, _) in best.items()})
    # Stronger integral sufficient bound already used in the main manuscript.
    proves = (isqrt(h) - 1) ** 4 > n
    return {"N": n, "curve": [a, b], "inputs": checked,
            "projections": projections, "fused": compact,
            "proves_whole_N": proves}


def crt_merge(m, residues, modulus, allowed):
    """Generalized CRT, preserving all possible signs/classes."""
    g = gcd(m, modulus)
    new_m = m // g * modulus
    new = set()
    for r in residues:
        for s in allowed:
            if (s - r) % g:
                continue
            v = modulus // g
            t = 0 if v == 1 else ((s - r) // g * pow(m // g, -1, v)) % v
            new.add((r + m * t) % new_m)
    return new_m, sorted(new)


def hasse_windows(n, h):
    """Every prime divisor <= sqrt(n) lies in this exact union of intervals."""
    need(n > 1 and h > 1, "positive modulus and nontrivial order")
    s = isqrt(n)
    tmax = (s + 1 + isqrt(4 * s)) // h
    intervals = []
    for t in range(1, tmax + 1):
        centre, radius = t * h + 1, isqrt(4 * t * h)
        lo, hi = max(2, centre - radius), min(s, centre + radius)
        if lo <= hi:
            intervals.append([lo, hi])
    return {"sqrt_floor": s, "T": tmax, "intervals": intervals}


def constrained_divisors(n, modulus, residues, intervals, budget=100000):
    """A16 closure AFTER the residue and elliptic evidence has been checked."""
    candidates = set()
    for lo, hi in intervals:
        for r in residues:
            first = lo + (r - lo) % modulus
            for d in range(first, hi + 1, modulus):
                if len(candidates) >= budget:
                    return {"status": "INCONCLUSIVE_BUDGET"}
                candidates.add(d)
    candidates = sorted(candidates)
    factors = [d for d in candidates if n % d == 0]
    return {"candidates": candidates, "divisor_checks": len(candidates),
            "status": "COMPOSITE" if factors else "PRIME",
            "proper_divisors": factors}


def verify_hybrid_certificate(n, arithmetic_factors, arithmetic_witnesses,
                              a, b, point, elliptic_factors):
    """A16 entry point: construct every predicate from checked evidence."""
    arithmetic = pocklington_partial(n, arithmetic_factors, arithmetic_witnesses)
    elliptic = ec_order_record(n, a, b, point, elliptic_factors)
    f = arithmetic["F"]
    modulus = 2 * f if f % 2 else f
    m, residues = crt_merge(1, [0], modulus, [1])
    windows = hasse_windows(n, elliptic["order"])
    closure = constrained_divisors(n, m, residues, windows["intervals"])
    return {"N": n, "arithmetic": arithmetic, "elliptic": elliptic,
            "windows": windows, "modulus": m, "residues": residues,
            "closure": closure}


def proof_plan(nodes, anchors, target, edges):
    """A17: minimum nonnegative additive cost in the verified-set graph."""
    index = {x: i for i, x in enumerate(nodes)}
    mask = lambda values: sum(1 << index[x] for x in set(values))
    start, goal = mask(anchors), 1 << index[target]
    prepared = [(mask(e["deps"]), 1 << index[e["out"]], e) for e in edges]
    dist, parent, queue = {start: 0}, {}, [(0, start)]
    while queue:
        cost, state = heappop(queue)
        if cost != dist[state]:
            continue
        if state & goal:
            path, end = [], state
            while state != start:
                old, name = parent[state]
                path.append(name)
                state = old
            return {"cost": cost, "path": path[::-1],
                    "visited_states": len(dist), "final_mask": end}
        for required, out, edge in prepared:
            need(edge["cost"] >= 0, "nonnegative costs")
            if state & out or state & required != required:
                continue
            new, value = state | out, cost + edge["cost"]
            if value < dist.get(new, float("inf")):
                dist[new] = value
                parent[new] = state, edge["name"]
                heappush(queue, (value, new))
    return None


def verify_plan(plan, anchors, edges):
    known = set(anchors)
    need(all(trial_prime(p) for p in known), "verified anchors")
    by_name = {e["name"]: e for e in edges}
    for name in plan["path"]:
        edge = by_name[name]
        need(set(edge["deps"]) <= known, "no circular proof dependency")
        n = edge["out"]
        if edge["kind"] == "trial":
            need(trial_prime(n), "trial certificate")
        else:
            factors, a = edge["factors"], edge["a"]
            need(set(factors) <= known, "every factor already proved")
            record = pocklington_partial(n, factors, {q: a for q in factors})
            need(record["proves_whole_N"], "Pocklington size bound")
        known.add(n)
    return sorted(known)


def inherit_divisor(n, lower, d):
    """A18: lower is a VERIFIED bound on every prime divisor of n."""
    need(1 < d < n and n % d == 0, "exact proper divisor")
    need(lower >= 2 and d < lower * lower, "strict inherited size bound")
    return {"divisor": d, "cofactor": n // d,
            "inherited_lower_bound": lower, "status": "PRIME"}


def must_reject(fn, message):
    try:
        fn()
    except ValueError:
        return
    raise ValueError("negative control accepted: " + message)


def run_all():
    began = time.perf_counter()
    out = {}
    fusion31 = prime_power_fusion(31, {2: 3, 3: 21})
    fusion1009 = prime_power_fusion(1009, {2: 11, 3: 2})
    need(3 * 21 % 31 == 1, "raw witness multiplication loses all orders")
    need(fusion31["proves_whole_N"] and fusion1009["proves_whole_N"],
         "A13 advertised certificates")
    fusion_checks = 0
    for n in range(3, 400, 2):
        witnesses = {}
        for q in factor_small(n - 1):
            for a in range(2, min(n, 18)):
                if pow(a, n - 1, n) == 1 and gcd(pow(a, (n - 1) // q, n) - 1, n) == 1:
                    witnesses[q] = a
                    break
        if not witnesses:
            continue
        rec = prime_power_fusion(n, witnesses)
        for ell in factor_small(n):
            need(pow(rec["z"], rec["H"], ell) == 1, "independent annihilation")
            need(all(pow(rec["z"], rec["H"] // q, ell) != 1
                     for q in witnesses), "independent exact primary orders")
        if rec["proves_whole_N"]:
            need(trial_prime(n), "independent A13 primality cross-check")
        fusion_checks += 1
    must_reject(lambda: prime_power_fusion(31, {2: 1}), "bad fusion witness")
    out["A13"] = {"examples": [fusion31, fusion1009],
                  "independent_small_moduli_checked": fusion_checks}

    h, z = fusion1009["H"], fusion1009["z"]
    positive = batch_units(1009, [pow(z, h // q, 1009) - 1 for q in [2, 3]])
    split = batch_units(91, [2, 7, 13, 5])
    unsuitable = batch_units(101, [0, 2, 3, 4])
    need(positive["status"] == "ALL_UNITS" and positive["tree_gcd_calls"] == 1,
         "one gcd positive path")
    need(split["root_gcd"] == 91 and split["proper_factors"] == [7, 13],
         "saturated product gcd localized")
    need(unsuitable["status"] == "NONUNIT_WITHOUT_PROPER_FACTOR", "not a compositeness claim")
    unit_checks = 0
    for n in range(2, 150):
        for values in [[2, 3, 5, 7], [0, n - 1, n + 1], [n, 2 * n, n * n],
                       [n // 2, n // 3, n // 5, n // 7]]:
            rec = batch_units(n, values)
            expected = [[i, gcd(v, n)] for i, v in enumerate(values) if gcd(v, n) > 1]
            need(rec["failed_leaves"] == expected, "gcd tree agrees with every leaf")
            unit_checks += 1
    out["A14"] = {"positive": positive, "factor_localization": split,
                  "unsuitable_witness": unsuitable, "independent_batches": unit_checks}

    n, a, b = 101, -4, -2
    p, q = (22, 16), (40, 39)
    ec = elliptic_fusion(n, a, b, [(p, {3: 1, 5: 1}), (q, {7: 1})])
    need(ec["fused"]["order"] == 105 and ec["proves_whole_N"], "fused elliptic certificate")
    raw = ec_add(p, q, n, a, b)
    need(raw == (69, 53), "advertised direct coprime sum")
    ec_order_record(n, a, b, raw, {3: 1, 5: 1, 7: 1})
    # An overlapping-order case where a naive sum loses the factor 5.
    p15 = ec_mul(7, raw, n, a, b)
    p35 = ec_mul(3, raw, n, a, b)
    naive = ec_add(p15, p35, n, a, b)
    need(ec_mul(21, naive, n, a, b) is None, "overlap cancellation control")
    overlap = elliptic_fusion(n, a, b, [(p15, {3: 1, 5: 1}),
                                      (p35, {5: 1, 7: 1})])
    need(overlap["fused"]["order"] == 105, "LCM fusion avoids overlap cancellation")
    negative_p = (p[0], (-p[1]) % n)
    need(ec_add(p, negative_p, n, a, b) is None, "raw point-sum cancellation")
    same_order = elliptic_fusion(n, a, b, [(p, {3: 1, 5: 1}),
                                         (negative_p, {3: 1, 5: 1})])
    need(same_order["fused"]["order"] == 15, "primary selection prevents total cancellation")
    must_reject(lambda: ec_order_record(101, 0, 0, (1, 1), {2: 1}), "singular curve")
    must_reject(lambda: ec_order_record(n, a, b, p, {7: 1}), "incorrect point order")
    must_reject(lambda: elliptic_fusion(-1, 0, 0, [((0, 0), {2: 1})]),
                "nonpositive elliptic modulus")
    try:
        ec_add((0, 1), (5, 14), 35, 0, 1)
    except NonunitDenominator as error:
        need(error.factor == 5, "proper denominator factor")
    else:
        raise ValueError("nonunit elliptic addition was accepted")
    need(trial_prime(n), "independent A15 cross-check")
    out["A15"] = {"coprime_fusion": ec, "direct_sum": list(raw),
                  "overlapping_fusion": overlap,
                  "naive_overlapping_sum_order_divides": 21,
                  "opposite_equal_order_points_cancel_raw": True,
                  "nonunit_control_factor": 5}

    n = 211
    hybrid = verify_hybrid_certificate(n, {5: 1}, {5: 2},
                                       -2, 2, (23, 27), {19: 1})
    order, windows, closure = (hybrid["arithmetic"], hybrid["windows"],
                               hybrid["closure"])
    need(not order["proves_whole_N"], "arithmetic uniform bound alone insufficient")
    s = isqrt(n)
    need(19 <= s + 1 + isqrt(4 * s), "elliptic uniform bound alone insufficient")
    need(windows["intervals"] == [[12, 14]], "exact Hasse window")
    need(closure["status"] == "PRIME" and closure["divisor_checks"] == 0,
         "empty hybrid intersection")
    need(trial_prime(n), "independent A16 cross-check")
    window_checks = 0
    for h in range(2, 120):
        ntest = 10000
        ww = hasse_windows(ntest, h)
        represented = {r for lo, hi in ww["intervals"] for r in range(lo, hi + 1)}
        direct = set()
        for r in range(2, isqrt(ntest) + 1):
            radius = isqrt(4 * r)
            low, high = max(1, r + 1 - radius), r + 1 + radius
            if ((low + h - 1) // h) * h <= high:
                direct.add(r)
        need(represented == direct, "integer window equivalence cross-check")
        window_checks += 1
    must_reject(lambda: verify_hybrid_certificate(n, {5: 1}, {5: 1},
                                                  -2, 2, (23, 27), {19: 1}),
                "unverified arithmetic predicate")
    must_reject(lambda: verify_hybrid_certificate(n, {5: 1}, {5: 2},
                                                  -2, 2, (23, 28), {19: 1}),
                "unverified elliptic predicate")
    hybrid["window_equivalence_checks"] = window_checks
    out["A16"] = hybrid

    nodes = [2, 3, 7, 13, 2003]
    edges = [
        {"name": "prove_3", "deps": [2], "out": 3, "cost": 7,
         "kind": "pock", "factors": {2: 1}, "a": 2},
        {"name": "7_via_3", "deps": [3], "out": 7, "cost": 1,
         "kind": "pock", "factors": {3: 1}, "a": 2},
        {"name": "13_via_3", "deps": [3], "out": 13, "cost": 1,
         "kind": "pock", "factors": {3: 1}, "a": 2},
        {"name": "7_direct", "deps": [], "out": 7, "cost": 6, "kind": "trial"},
        {"name": "13_direct", "deps": [], "out": 13, "cost": 6, "kind": "trial"},
        {"name": "goal", "deps": [7, 13], "out": 2003, "cost": 1,
         "kind": "pock", "factors": {7: 1, 13: 1}, "a": 3},
    ]
    plan = proof_plan(nodes, [2], 2003, edges)
    need(plan is not None and plan["cost"] == 10, "optimal shared proof plan")
    known = verify_plan(plan, [2], edges)
    # Independent exhaustive enumeration of edge subsets in this finite example.
    optimum = float("inf")
    for subset in range(1 << len(edges)):
        selected = [e for i, e in enumerate(edges) if subset >> i & 1]
        closure_nodes = {2}
        while True:
            newer = closure_nodes | {e["out"] for e in selected if set(e["deps"]) <= closure_nodes}
            if newer == closure_nodes:
                break
            closure_nodes = newer
        if 2003 in closure_nodes:
            optimum = min(optimum, sum(e["cost"] for e in selected))
    need(optimum == plan["cost"], "exhaustive optimality cross-check")
    cycle = [{"name": "c1", "deps": [7], "out": 3, "cost": 0},
             {"name": "c2", "deps": [3], "out": 7, "cost": 0}]
    need(proof_plan([2, 3, 7], [2], 7, cycle) is None, "unanchored cycle rejected")
    out["A17"] = {"nodes": nodes, "anchors": [2], "catalogue": edges,
                  "plan": plan, "verified_primes": known,
                  "independent_branch_model_cost": 13,
                  "costs_are_assigned_model_units_not_wall_times": True,
                  "edge_subsets_checked": 1 << len(edges),
                  "unanchored_cycle_rejected": True}

    n = 1387
    partial = pocklington_partial(n, {3: 2}, {3: 2})
    need(not partial["proves_whole_N"], "auxiliary modulus not certified prime")
    residue = pow(2, (n - 1) // 2, n)
    d = gcd(residue - 1, n)
    need(d == 73 and n // d == 19, "exact factor extraction")
    inherited = [inherit_divisor(n, partial["prime_divisor_lower_bound"], x)
                 for x in [d, n // d]]
    need(all(trial_prime(x["divisor"]) and x["divisor"] > 3 for x in inherited),
         "independent harvested-prime cross-check")
    must_reject(lambda: inherit_divisor(6859, 19, 361), "strict inherited size boundary")
    must_reject(lambda: inherit_divisor(n, 19, 17), "nondivisor")
    out["A18"] = {"auxiliary_N": n, "seed": 3, "coefficient": 154,
                  "partial_order_certificate": partial,
                  "factor_extraction_residue": residue, "factor": d,
                  "harvested": inherited,
                  "new_modular_powers_for_inherited_divisors": 0}
    out["status"] = "ALL_SIX_CERTIFICATE_ALGORITHMS_VERIFIED"
    out["audit_runtime_seconds"] = time.perf_counter() - began
    return out


if __name__ == "__main__":
    print(json.dumps(run_all(), indent=2, sort_keys=True))
\end{lstlisting}

\clearpage
\section{Exact verification of the continuing seven-node chain}
\label{app:later-chain}

This program reconstructs every chain node, proves its known factors in dependency order, and checks both the original Pocklington route and the alternate A01 obligations where applicable.

\begin{lstlisting}[label=lst:later-chain]
"""Exact later seven-node chain; optional --gmp, no external data files."""
import argparse
import hashlib
import json
import math
import sys
import time

if hasattr(sys, "set_int_max_str_digits"):
    sys.set_int_max_str_digits(0)

def require(ok, message):
    if not ok:
        raise ValueError(message)

def verify(use_gmp=False):
    if use_gmp:
        import gmpy2
        integer, power, gcd = gmpy2.mpz, gmpy2.powmod, gmpy2.gcd
    else:
        integer, power, gcd = int, pow, math.gcd
    start = time.perf_counter()
    p = integer(14)**512 - integer(14)**256 + 1
    expected = [
        (587, "a338e3a2cccacc978def8f28b08c248d39d5cd1542a47f8f0cf627d03febfb0c"),
        (1177, "c852e5af4a1166124ffff6084fb1d986da411fe5b210aace6739009180e12c3f"),
        (2357, "f1c08bab7b1e51ec1eef6dc053d1d98a2a13eb74e323c7273ac65c2dd2c58785"),
        (4718, "c3e84997a4b2a949282840d14981c52bd98d66982c022a38528038779cd32566"),
        (9439, "82ddc0fab24f40e0c8c777c2f2e2138fec8d9bc505e21eb414858111ca910089"),
        (18882, "7407dab961f68eb15854c74700a8007bf30d89410a6beeb1d2fc7d365f4b7ce8"),
        (18885, "1fd5cb8c8470b74c24c163501a2775693b19707db62c65e6b074a1e56ff7cbd3"),
    ]
    nodes = [(p, [(integer(2),256,11), (integer(7),256,2)])]
    for k,e in [(438,2),(13696,2),(5494,2),(13258,2),(22330,2),(5936,1)]:
        n = k*p**e + 1
        nodes.append((n, [(p,e,2)]))
        p = n
    proved = {integer(2), integer(7)}
    rows = []
    for i,(n, factors) in enumerate(nodes):
        began = time.perf_counter()
        f = integer(1)
        for q,e,a in factors:
            require(q in proved and e > 0, "unproved factor")
            f *= q**e
        require((n-1)%f == 0 and f*f > n, "factor-size condition")
        for q,e,a in factors:
            require(power(integer(a),n-1,n) == 1, "Fermat congruence")
            z = power(integer(a),(n-1)//q,n)
            require(gcd(z-1,n) == 1, "Pocklington gcd")
        decimal = str(n)
        digest = hashlib.sha256(decimal.encode("ascii")).hexdigest()
        require((len(decimal),digest) == expected[i], "identity mismatch")
        alternate = None
        if 1 <= i <= 5:
            q = factors[0][0]
            k = (n-1)//(q*q)
            require(q > 2 and q % 2 and k > 0 and k % 2 == 0,
                    "A01 parity")
            require(k < 4*q-2, "A01 strict coefficient bound")
            short = power(integer(2), k, n)
            require(short != 1, "A01 short nonunity")
            # The full base-2 Fermat condition was verified above.
            alternate = {
                "algorithm":"A01", "coefficient":int(k),
                "short_power_is_one":False,
                "short_nonunity_from_size":bool(integer(1)<<int(k) < n),
                "candidate_bits":int(n.bit_length())}
        proved.add(n)
        rows.append({"index":i,"digits":len(decimal),"sha256":digest,
                     "alternate_quadratic_certificate":alternate,
                     "seconds":time.perf_counter()-began})
    return {"status":"ALL_SEVEN_PRIMES_CERTIFIED",
            "arithmetic":"gmpy2" if use_gmp else "Python integers",
            "seconds":time.perf_counter()-start,"nodes":rows}

if __name__ == "__main__":
    parser = argparse.ArgumentParser()
    parser.add_argument("--gmp",action="store_true")
    args = parser.parse_args()
    print(json.dumps(verify(args.gmp),indent=2))

\end{lstlisting}

\clearpage
\section{Independent checks of the coefficient boundaries}
\label{app:boundary-audit}

This audit uses direct Lucas recurrence and exact small trial division to check the positive and negative coefficient closures independently of the construction program's matrix implementation.

\begin{lstlisting}[label=lst:boundary-audit]
"""Independent bounded audit of the two quadratic closure lemmas.
Small primality truth uses exhaustive trial division. Lucas values here use
the direct recurrence, independently of fast matrix/doubling implementations.
"""
from math import gcd, isqrt
import json, time

def require(ok, text):
    if not ok:
        raise ValueError(text)

def prime(n):
    return n >= 2 and all(n%d for d in range(2,isqrt(n)+1))

def lucas_pair_direct(n,a,c,b):
    old,cur=0,1
    short=0 if a==0 else None
    for j in range(1,n+2):
        if j==a:
            short=cur
        if j==n+1:
            return short,cur
        old,cur=cur,(c*cur-b*old)%n
    raise ValueError("unreachable")

def run():
    start=time.perf_counter()
    pos_cases=pos_passes=neg_cases=neg_passes=0
    for p in range(3,100,2):
        if not prime(p):
            continue
        for a in range(2,4*p-2,2):
            n=a*p*p+1
            truth=prime(n)
            for b in (2,3,5,9):
                pos_cases+=1
                passed=pow(b,n-1,n)==1 and pow(b,a,n)!=1
                require(not passed or truth,"positive false certificate")
                pos_passes+=passed
    for p in (3,5,7,11,13,17,19):
        for a in range(2,4*p+2,2):
            if isqrt(a)**2==a:
                continue
            n=a*p*p-1
            truth=prime(n)
            for c,b in ((1,1),(3,1),(4,1),(5,1),(5,2)):
                disc=c*c-4*b
                if gcd(n,2*b*disc)!=1:
                    continue
                ua,ufull=lucas_pair_direct(n,a,c,b)
                neg_cases+=1
                passed=ufull==0 and gcd(ua,n)==1
                require(not passed or truth,"negative false certificate")
                neg_passes+=passed
    require(91==7*13 and pow(9,90,91)==1 and pow(9,10,91)!=1,
            "positive boundary control")
    for a,n in ((4,35),(14,125)):
        ua,ufull=lucas_pair_direct(n,a,1,1)
        require(not prime(n) and gcd(n,-6)==1 and ufull==0
                and gcd(ua,n)==1,"negative missing-guard control")
    # A single correction is not guaranteed for arbitrary x < N^2.
    n,radix,k,x=3,2,1,8
    big,low=divmod(x,radix)
    t,s=divmod(big,k)
    y=low+s*radix-t
    corrected=y+n if y<0 else y
    require(x<n*n and corrected!=x%n,"fold domain control")
    return {"status":"INDEPENDENT_BOUNDARIES_VERIFIED",
            "positive_attempts":pos_cases,"positive_certificates":pos_passes,
            "negative_attempts":neg_cases,"negative_certificates":neg_passes,
            "false_prime_certificates":0,
            "guard_counterexamples":[91,35,125],
            "fold_out_of_domain":{"N":3,"x":8,"one_correction":corrected,
                                  "correct_remainder":x%n},
            "seconds":time.perf_counter()-start}

if __name__=="__main__":
    print(json.dumps(run(),indent=2))

\end{lstlisting}

\end{document}
